% ****** Start of file apssamp.tex ******
%
%   This file is part of the APS files in the REVTeX 4.2 distribution.
%   Version 4.2a of REVTeX, December 2014
%
%   Copyright (c) 2014 The American Physical Society.
%
%   See the REVTeX 4 README file for restrictions and more information.
%
% TeX'ing this file requires that you have AMS-LaTeX 2.0 installed
% as well as the rest of the prerequisites for REVTeX 4.2
%
% See the REVTeX 4 README file
% It also requires running BibTeX. The commands are as follows:
%
%  1)  latex apssamp.tex
%  2)  bibtex apssamp
%  3)  latex apssamp.tex
%  4)  latex apssamp.tex
%
\documentclass[%
reprint,
superscriptaddress,
%groupedaddress,
%unsortedaddress,
%runinaddress,
%frontmatterverbose, 
%preprint,
%preprintnumbers,
nofootinbib,
%nobibnotes,
%bibnotes,
 amsmath,amssymb,
% aps,
prx,
%prb,
%rmp,
%prstab,
%prstper,
floatfix,
]{revtex4-2}

%\usepackage{amssymb}
%\usepackage{multirow}
\usepackage{mathtools}


\usepackage{lipsum,bbm,dsfont} 

\usepackage{graphicx}% Include figure files
\usepackage{dcolumn}% Align table columns on decimal point
\usepackage{bm}% bold math
\usepackage[hidelinks,colorlinks=true,linkcolor=blue,citecolor=blue]{hyperref}
\usepackage{xcolor}
\usepackage{dsfont}
\usepackage{graphicx}
\usepackage{adjustbox}

\usepackage{caption} 
\captionsetup{justification=justified, singlelinecheck=false}

\usepackage{array}
\newcolumntype{P}[1]{>{\centering\arraybackslash}p{#1}}

%\bibliographystyle{apsrev4-1}
\usepackage{braket}
%\usepackage{hyperref}% add hypertext capabilities
%\usepackage[mathlines]{lineno}% Enable numbering of text and display math
%\linenumbers\relax % Commence numbering lines

%\usepackage[showframe,%Uncomment any one of the following lines to test 
%%scale=0.7, marginratio={1:1, 2:3}, ignoreall,% default settings
%%text={7in,10in},centering,
%%margin=1.5in,
%%total={6.5in,8.75in}, top=1.2in, left=0.9in, includefoot,
%%height=10in,a5paper,hmargin={3cm,0.8in},
%]{geometry}
\newcommand{\dket}[1]{\lvert #1 \rangle\!\rangle}
\newcommand{\dbra}[1]{\langle\!\langle #1 \rvert}

\begin{document}

\preprint{APS/123-QED}

%\title{Gaussian Formalism for $n$-Modes Diagonal Interacting with One or Two Qubits}

%\title{Gaussian Cat States of $n$-Modes with Ancilla Qubits: General Dynamics, Noise and Measurements}

%\title{Gaussian Cat States of $n$-Modes with Ancilla Qubits: Open Dynamics and Measurements}

%\title{$n$-Modes vs $m$-Qubits: Diagonal Open Gaussian Dynamics and Noisy Measuraments}

%\title{Phase methods for superposition of Gaussian Satate with ancillary qubit. }

%\title{Superposition of Phase Spaces: Solution of $n$-Modes Open Gaussian Dynamics with Qubit Diagonal Interactions}

%\title{Quantum Superposition of Phase Spaces: Gaussian Processes of Non-Gaussian States}

%\title{Quantum Superpositions of Phase Spaces: Gaussian-Branched Cat States}

\title{Superpositions of Quantum Gaussian Processes}

\author{Lorenzo Braccini}
\affiliation{Department of Physics and Astronomy, University College London, Gower Street, WC1E 6BT London, United Kingdom}
%Lines break automatically or can be forced with \\

\author{Sougato Bose}
\affiliation{Department of Physics and Astronomy, University College London, Gower Street, WC1E 6BT London, United Kingdom}

\author{Alessio Serafini}
\affiliation{Department of Physics and Astronomy, University College London, Gower Street, WC1E 6BT London, United Kingdom}


\date{\today}% It is always \today, today,
             %  but any date may be explicitly specified
%\keywords{Suggested keywords}%Use showkeys class option if keyword
                              %display desired
\begin{abstract}
%We extend the Gaussian formalism of Continuous Variable (CV) systems to any superposition of Gaussian states of $n$ modes entangled with $N$ qubits, i.e. ``Gaussian-Branched Cat States" (GBCS). This is achieved by associating a phase space of $n$ modes with each element of the joint destiny matrix of the qubits, that is, a ``superposition of phase spaces''. Any GBCS is fully characterised by its phase-spaces quantities, which are sets of $2^{2N}$ $2n$-dimensional vectors and matrices (which contain the first and second moments of the branches) and the reduced density matrix of the qubits. By restricting to Gaussian Hamiltonians with diagonal qubit interactions, any time evolution -- unitary, open, monitored, and with feedback -- can be mapped to a set of Ordinary Differential Equations (ODEs) governing the time evolution of these quantities, which solutions give an exact description of the general dynamics of GBCS. Any measurements, both on the qubits and the modes, ideal or noisy, can be easily expressed in terms of phase-space quantities. A stern-gerlach interferometer of a levitated mass undergoing diffusive dynamics (sensing an unknown force) and two qubits interacting via a noisy and monitored resonator (generating measurement-based entanglement) are given as paradigmatic examples.

We generalise the Gaussian formalism of Continuous Variable (CV) systems to describe their entanglement with Discrete Variable (DV) systems, leading to superposition of CV Gaussian states. A new class of CV-DV entangled states, named Gaussian-Branched Cat States (GBCSs), yield an analytical formalism to describe quantum hybrid systems. GBCSs~are fully characterised by their superposed phase-space parameters: sets of generalised complex first moments and covariance matrices, along with the DV reduced density matrix (phases and contrasts). These states arises in all the instances where Gaussian dynamics, operations, and measurements, are performed conditionally on a DV state. The time evolution of the GBCS phase space parameters allows --  via a new set of equations in closed form -- to analytically treat a large set of unitary and open dynamics, generated by Gaussian Hamiltonians labelled by DV eigenvalues. Conditional operations, such as displacements and rotations, and joint Gaussian measurements (homodyne/heterodyne jointly with DV projectors), can be both described as maps of GBCS's parameters. A phase-space perturbation theory is given to extend the analysis to non-orthogonal DV super-operators, e.g. DV decay. We showcase our general formalism with two paradigmatic examples of experimental modelling: i) the transient behaviour of a driven parametric amplified resonator used for dispersive coupling read-out of a qubit; ii) the generation of the Wigner negativity and force sensing of a levitated nanoparticle undergoing Stern-Gerlach interferometry in a diffusive environment.

\hspace{0.5cm}


%We extend the Gaussian formalism of Continuous Variable (CV) systems to Gaussian-Branched Cat States (GBCSs), defined as $n$-mode $N$-qubit entangled states where the terms of the superposition are tensor products of Gaussian states and a set of orthogonal qubit states. This is achieved by associating a phase space of $n$ modes with each element of the $N$--qubits' joint destiny matrix, that is, a ``superposition of phase spaces''. Any GBCS is fully characterised in $2^{2N}$ phase spaces by sets of $2n$-dimensional vectors and matrices (containing the first and second moments of the branches) and the reduced density matrix of the qubits. By restricting to Gaussian Hamiltonians with diagonal qubit interactions, any time evolution -- unitary, open, monitored, and with feedback -- can be mapped to the time evolution of such phase-space quantities. Any qubit or Gaussian measurement, ideal or noisy, can be easily expressed in terms of these parameters. We showcase our general formalism with two paradigmatic examples: i) the certification of the Wigner negativity of a levitated nanoparticle undergoing noisy Stern-Gerlach interferometry; ii) the evaluation of two-qubit entanglement after interaction through a monitored leaky resonator.
\end{abstract}

\maketitle




\section{Introduction \label{sec:intro}}


The interplay between Discrete-Variable (DV) quantum systems, such as qubits, and Continuous-Variable~(CV) ones, such as quantum modes, is pivotal to most of the current research in quantum technology and experiments, ranging from the detection of quantum correlations to state-of-the-art quantum computers~\cite{hensen_loophole_2015,Reglade_quantum_2024}. Hybrid CV-DV systems can be used to engineer a vast number of quantum protocols, from non-classical state preparation in the CV to computation in the DV.

In CV systems, quantum Gaussian processes \textit{per se} have found application in a wide number of protocols~\cite{serafini_quantum_2017,quantum_braunstein_2003,weedbrook_gaussian_2012}, such as entanglement generation~\cite{peres_horodecki_2000,giedke_entanglement_2003,garttner_entanglement_2023}, quantum key distribution~\cite{rudolph_teleportation_2001,grosshans_cryptography_2002} and quantum control routines~\cite{wiseman_optimal_2005,wiseman_quantum_2010}. Describing quantized bosonic field operators, CVs have found experimental implementations in quantum optics~\cite{walls_quantum_1994,bowen_experimental_2003,bowen_teleport_2003}, circuit electrodynamics~\cite{blais_circuit_2021,flurin_generating_2012,fedorov_unconditional_2021}, atomic systems~\cite{leibfried_quantum_2003}, collective spins for atomic ensembles~\cite{hammerer_quantum_2010}, and levitated nanoparticles~\cite{ballestero_levitodynamics_2021,delic_cooling_2020,tebbenjohanns_motional_2020,magri_real_2021}. The Gaussian structure of states, interactions, noises, operations, and measurements, allows for an exact description at the level of the phase space (avoiding truncations), dispensing one with the cumbersome task of solving quantum Gaussian processes in the infinite-dimensional Hilbert space~\cite{serafini_quantum_2017}. In fact, Gaussian processes can be formulated in terms of evolutions and maps of a vector of expectation values (first moments) and a matrix of uncertainties (covariance matrix). Uniquely determining a Gaussian states, these state parameters have a clear operational description as they are directly related to quantities that are often measured in the experimental practice. 

\begin{figure}[!b]
    \centering
    \includegraphics[width=1\linewidth]{Cats_Diagram.pdf}
    \caption{Representation of a superposition of Gaussian processes generating a Gaussian-Branched Cat States.}
    \label{fig:enter-label}
\end{figure}


\newpage However, in CV systems, non-Gaussian states are crucial to unlock most quantum advantages, since Gaussian can in principle be mimicked by classical probability distributions: for instance, non Gaussian elements are needed for universal quantum computation~\cite{lloyd_quantum_1999,Jeong_Effective_2002,Ralph_Quantum_2003,albarelli_resource_2018,Eriksson2024} and error correction~\cite{gottesman_encoding_2001,campagne-ibarcq_quantum_2020,dahan_creation_2023}. Non-Gaussian state represents excellent quantum sensors, which could test the macroscopic limit of quantum mechanics and non-classical gravitational effects~\cite{romero-isart_large_2011,roda_macroscopic_2024,bose_spin_2017,marletto_gravitationally-induced_2017}. A classic way to generate non-Gaussianity and delocalisation is to create a distinguishable superposition of CV Gaussian states, often named Schr\"odinger cat states~\cite{Shore_jaynes_1993}. They naturally arise in CV-DV hybrid systems from fundamental light-matter interactions and, typically, these hybrid states are of the form
\begin{equation}
    \ket{\psi}= \frac{1}{\sqrt{2}}\left( \ket{g_\uparrow} \ket{\uparrow} + \ket{g_\downarrow}\ket{\downarrow} \right) ,
\end{equation}
where $\ket{\uparrow}$ and $\ket{\downarrow}$ are some DV states and $\ket{g_i}$ labels some CV Gaussian state.


DV-CV cat states have arisen in classic foundational experiments~\cite{Monroe_Schrodinger_1996,brune_observing_1996} and are still central to recent protocols~\cite{cavailles_demosntration_2018,schrodinger_bild_2023,Realization_Xiaozhou_2025}, as they posseses quantum advantages for both technological applications and foundational tests~\cite{Sanders_2012,Munro_weekforce_2002,Stob_violation_2007,Leghtas_Hardware_2013,Mirrahimi_Dynamically_2014,Bergmann_quantum_2016,wbhn-v1sw}. It is remarkable how often quantum systems have shown the capability of generating such states, for instance, in the fields of non-linear quantum optics~\cite{gerry_quantum_1997,serafini_minimum_2004} and optomechanics~\cite{bose_preparation_1997,bose_scheme_1999,qvarfort_time-evolution_2020,marchese2026cascadedoptomechanicalsensingsmall}, atomic~\cite{machluf_coherent_2013,Amit_T_2019,margalit_realization_2020,keil_stern_2021} and matter-wave~\cite{bose_qubit_2006,Romero_Isart_quantum_2011,scala_matter_2013} interferometry, and quantum field theory~\cite{Funai_gaussian_2025,sun2026representationtheoryinhomogeneousgaussian}. Their applications represents the building blocks of modern quantum computers in numerous architectures~\cite{xiang_hybrid_2013}, including superconducting~\cite{strong_wallraff_2004,coherent_chiorescu_2004,blais_quantum_2007,coupling_majer_2007}, solid-state~\cite{strong_samkharadze_2018,strong_samkharadze_2018,coherent_mi_2018,coherent_landig_2018,strong_kubo_2010,strong_chen_2018,realization_liu_2023}, and ion traps systems~\cite{cirac_quantum_1995,sorensen_quantum_1999,leibfried_quantum_2003,haljan_spin_2005,Fluhmann2019,gan_hybrid_2020}.  The frequent appearance of these states is due to the fact that, often, the leading-order DV-CV nonlinearties -- hence, the ones common in practical set-ups~\cite{florian_ion_2001,webb_resilient_2018,kaku2025suddendecoherenceresonantparticle,tobar2025quantummetasurfacesprobesvacuum,ayyash2025multimodequbitconditionaloperationsgeneralized,Ayyash_driven} -- are the ones that are quadratic in the CV observables $\hat{X}$ and $\hat{P}$ (or $\hat{a}$ and $\hat{a}^\dag$) and function of a single DV operator. We shall term this set of  Hamiltonians as operator-valued Gaussian interactions. 

In some instances, both CV and DV subsystems play a central role in the hybrid system, while, in other instances, the presence of either subsystem is hidden behind effective Hamiltonians as direct qubit-qubit~\cite{Kentaro_crosscross_2021} or non-Gaussian CV interactions~\cite{semiao_kerr_2009,Nigg_black_2012,Leghtas_deterministic_2013}. Yet, fundamentally and operationally, such interactions are almost always mediated by CV subsystems, such as positions of atoms or states of resonators. Solutions for experimental predictions of these hybrid systems are often achieved numerically by truncating the CV Hilbert space (in the $n$ basis or other appropriate spectral decompositions) or limited to specific analytical examples.

Despite the obvious importance of these states, to the best of our knowledge, no generalization of the Gaussian CV formalism exists to analytically treat entangled CV-DV states under open-system dynamical conditions, measurements and operations, which generate superpositions of Gaussian states.
We therefore introduce ``Gaussian-Branched Cat States" (GBCS): $n$-mode and $N$-dimensional DV entangled states, where the terms of the superposition are tensor products of Gaussian states and a set of orthogonal qubit states. GBCS admit an operational description where a phase space of $n$ CV degrees of freedom is associated with each element of an $N\times N$ DV density matrix forming, so to speak, a ``superposition of phase spaces''. In this representation, any GBCS is uniquely characterised by sets of $2n$-dimensional complex vectors and matrices (containing the first moments and covariance matrix of each branch), and by the reduced density matrix of the DV. In practice, GBCS arise in all protocols where Gaussian operations, dynamics, and measurements, are performed conditionally to the state of a DV subsystem, alongside any quantum Gaussian processes and DV operations, which we shall term \textit{``Superpositions of Quantum Gaussian Processes"}. 

%One may ask if nature can generate GBCS, and if experimental protocols of CV-DV hybrid system can be analytically described by the evolutions of the phase space parameters characterizing these states. Remarkably, as it has been the case for Gaussian states in several physical systems~\cite{walls_quantum_1994,bowen_experimental_2003,bowen_teleport_2003,blais_circuit_2021,flurin_generating_2012,fedorov_unconditional_2021,ballestero_levitodynamics_2021,delic_cooling_2020,tebbenjohanns_motional_2020,magri_real_2021}, this is indeed formally possible and also found to be frequent in experimental practice: 

Being based on Gaussian building blocks, the chosen representation of GBCS allows one to directly apply a large set of 
Gaussian phase-space results, firstly developed in quantum optics~\cite{serafini_quantum_2017} and summarized in Sec.~\ref{sec:intro_phase}. In Sec.~\ref{sec:sup_phase_space}, the most generals GBCSs is defined via its phase spaces parametrization. At the dynamical level, GBCS solve without truncation the most general operator-valued Gaussian interaction between $n$ modes and $N$-level system via the time evolution of its superposed phase-space parameters, formulated in Sec~\ref{sec:dynamics} for both unitary and open dynamics. The two central case studies of this work are the dispersive coupling ($\hat{a}^{\dag} \hat{a} \otimes \hat{\sigma}_z$) and the Stern-Gerlach coupling ($\hat{X} \otimes \hat{\sigma}_z$), describing qubits in off-resonance time-dependent electromagnetic fields and gradients, respectively~\cite{xiang_hybrid_2013}.  Sec.~\ref{sec:operations} deals with conditional operations, such as rotations and displacements, which -- with the inclusion of local qubit and/or Gaussian operations -- have been shown to be universal for computation of either the CV and for the DV subsystems~\cite{Hybrid_loock_2008,krastanov_universal_2015,Heeres_Implementing_2017,Eickbusch_fast_2022,grosso_controlled_2025}. This state parametrization naturally gives a phase-space picturization of such non-Gaussian processes and provides a geometrical understanding of Wigner fringes,  illustrated in Sec.~\ref{sec:wigner}.  
Gaussian and qubit (ideal and noisy) joint measurements can be described as maps of the phase-space parameters (Sec.~\ref{sec:measure}). In Sec~\ref{sec:perturbation}, we place GBCS in the context of a perturbation expansion, thus allowing to describe non-diagonal DV super-operators such as DV decaying noise. By analytically describing the leading order nonlinearity and noises, along with the measurement and operations accessible in experiments, the proposed formalism is applicable to solve structured dynamics and make predictions for experiments in realistic conditions. In Sec.~\ref{sec:physical_examples}, our methodology is applied to the modelling of physical examples: (i) the dispersive read-out of a qubit via a resonator with a parametric amplifier; (ii) a Stern-Gerlach interferometer for a levitated nanoparticle in a diffusive environment. Conclusions are drawn in Sec.~\ref{sec:conclusion}, with comparison to other known methodologies. \textcolor{red}{This work is purely analytical. }
%\textcolor{red}{A Python} software library and a Mathematica Notebook implementing the dynamics, operations, and measurements studied in this work are provided on \href{https://github.com/LorenzoBracciniUCL/GBCS}{GitHub}, including visualisation of time-dependent Wigner functions.




%Such interactions frequently arise in nature in many physical systems and play essential roles in entanglement between CV subsystem and qubits. For example, the interaction describes the electronic or magnetic coupling between a qubit and the electromagnetic field, which -- in the off-resonance dissipative regime -- is operator-valued and Gaussian . The , which couples an atom with the gradient of the electromagnetic field, also leads to an interaction of this form . The treatment presented here allows one to solve the arising dynamics between any number of modes and qubits, for open, monitored and feedback dynamics, \textcolor{red}{cite jappan second order potential}

%The probability distributions that describe a quantum state over such a space are known as characteristic functions, and for a Gaussian quantum state they are complex Gaussians, fully determined by . Thus, the time evolution of the system can be map to two Ordinary Differential Equations (ODEs) describing the evolution of the first and second moments. This procedure can be extended to Markovian open quantum dynamics (master equations) and continuous monitoring and feedback (stochastic master equations). The phase space method allows one to give a simple closed form for outcomes of Gaussian measurements and correlations -- for instance, homodyne measurements and values of entanglement witnesses. 







\section{Introduction to Gaussian States \\ \& Phase Space Methods \label{sec:intro_phase}}


In this section, with the purpose of introducing notation and laying out the Gaussian techniques that will be used, a review of Gaussian phase space methods for CV systems is presented (for more details and derivation, see Ref. \cite{serafini_quantum_2017}). As we shall see, this representation of Gaussian quantum states allows for the analytical treatment of a large set of dynamics, measurements, and operations, which will be key to the remainder of this work. 

The expectation value of a quantum operator $\hat{O}$ is denoted with $\braket{\hat{O}} = \text{Tr} [ \hat{O} \hat{\varrho} ]$, where $\hat{\varrho} $ is the  density matrix of the system under exam (operators will wear a hat, $\hat{\cdot}$, while phase space quantities will not). Consider a set of $n$ continuous degrees of freedom, with canonical operators $\hat{X}_j$ and $\hat{P}_j$. The Hilbert space associated with the $n$-dimensional CV system (modes) is $\mathcal{H} = L^2 (\mathds{R}^n)$. The dimensionless canonical operators $\hat{x}_j=\hat{X}_{j}\sqrt{\frac{m_j\omega_j}{\hbar}}$ and $\hat{p}_j=\hat{P}_{j}\sqrt{\frac{1}{m_j\omega_j\hbar}}$, where $m_j$ and $\omega_j$ are the mass and frequency associated with the degree of freedom, have commutation relations $[\hat{x}_j, \hat{p}_k] = i\delta_{jk}$.  Thus, it is possible to define the vector of operators
\begin{equation}
\label{eq:def_first_second}
    \hat{r} := \begin{pmatrix}
        \hat{x}_1 &
        \hat{p}_1 &
        \hat{x}_2 & 
        ... &
        \hat{x}_n &
        \hat{p}_n 
    \end{pmatrix}^{\rm T} \; , 
\end{equation}
such that the commutator relations reads $[ \hat{r}, \hat{r}^{\rm T} ] = i \Omega_n $, where  $\Omega_n = \bigoplus_{i = 1}^n \Omega_1$ is the $n$ dimensional symplectic form, with
\begin{equation}
    \Omega_1 = \begin{pmatrix}
        0 & 1 \\
        -1 & 0
    \end{pmatrix} \;.
\end{equation}
It may be useful to recall the definition of the symplectic group as the set of real matrices $S \in Sp_{2, \mathds{R}}$ such that $S \Omega S^{\rm T} = \Omega$. Furthermore, the symplectic inner product is antisymmetric, i.e., $r_1^{\rm T} \Omega r_2 = - r_2^{\rm T} \Omega r_1 $ for any pair of vectors, and $\Omega^{-1} = \Omega^{\rm T} = - \Omega$. The ladder operators are defined as $\hat{a}_i = (\hat{x}_i + i \hat{p}_i)/ \sqrt{2}$, and one can move from one representation to the other by the use of the transformation $U_n = \bigoplus_{i = 1}^n U_1$ such that the vectors  $\hat{a} = U_n \hat{r} = (\hat{a}_1, \hat{a}_1^\dag, ... , \hat{a}_n, \hat{a}_n^\dag)^{T}$, with
\begin{equation}
\label{eq:change_basis}
    U_1 = \frac{1}{\sqrt{2}}
   \begin{pmatrix}
        1 & i \\
        1 & -i
    \end{pmatrix}   \implies  \begin{pmatrix}
        \hat{a} \\
        \hat{a}^\dag
    \end{pmatrix} = U_1 \begin{pmatrix}
        \hat{x} \\
        \hat{p}
    \end{pmatrix}\;.
\end{equation}
In this compact notation, the first and second moments of a quantum state $\hat{\varrho}$ are defined as the expectation values $\bar{r} := \text{Tr} [ \hat{r} \hat{\varrho}]$ and $\sigma := \text{Tr} [ \{ (\hat{r}-\bar{r}), (\hat{r}-\bar{r})^{\rm T} \} \hat{\varrho} ] $, respectively, where $\{ \cdot, \cdot \}$ represents anti-commutation. The former is a $2n$-dimensional real vector and the latter is a $2n \times 2n$ matrix, also known as the covariance matrix, which is real, symmetric, and positive semi-definite. For instance, in the case of $n=1$ mode, such parameters are
\begin{equation}
\label{eq:def_first_second}
    \bar{r}  = \begin{pmatrix}
        \braket{x}\\
        \braket{p}
    \end{pmatrix} , \hspace{1cm}
    \sigma = \begin{pmatrix}
        2\Delta_{x^2} & \Delta_{\{x,p\}} \\
        \Delta_{\{x,p\}} & 2\Delta_{p^2}
    \end{pmatrix} \;,
\end{equation}
where $\Delta_A = \braket{\hat{A}^2} - \braket{\hat{A}}^2$. In this formalism, the Heisenberg uncertainty principle reads $\sigma + i \Omega \geq 0$.
 
The first and second moments uniquely define a quantum Gaussian state of the $n$ modes through its phase-space description. The phase space methods relays on the set of displacement operators
\begin{equation}
    \hat{D}_{\xi} = {\rm e}^{i \xi^{\rm T} \Omega \hat{r}} \; , 
\end{equation}
where $\xi \in \mathds{R}^{2n}$, which are the operators that displace the vacuum to a coherent state. The displacement operators $\hat{D}_\xi$ are orthonormal with respect to the Hilbert--Schmidt inner product, such that any bounded operator $\hat{O}$ on the space can be written according to the Fourier–Weyl relation
\begin{equation}
    \hat{O} = \frac{1}{(2\pi)^n} \int_{\mathds{R}^{2n}} \hspace*{-0.2cm}\text{d} \xi\,  \text{Tr} \left[ \hat{D}_{-\xi} \hat{O} \right] \hat{D}_{\xi} \;. 
\end{equation}
If the operator $\hat{O}$ is a density matrix $\hat{\varrho}$, the relationship above defines the characteristic function $\chi(\xi) = \text{Tr} \left[ \hat{D}_{-\xi} \hat{\varrho} \right]$, which represents the quantum state via
\begin{equation}
    \hat{\varrho} =  \frac{1}{(2 \pi)^n} \int_{\mathds{R}^{2n}} \text{d} \xi  \; \chi (\xi) \hat{D}_{\xi} \;,
\end{equation}
where $r$ is the vector of the $2n$ phase-space variables. The characteristic function is generally complex and, by definition, $\text{Tr} [\hat{\varrho}] = \chi (0) = 1$. The characteristic function of a general Gaussian state ($\chi_g$) is fully determined by its first moments and covariance matrix, as per
\begin{equation}
\label{eq:chara_general_gaussian}
    \chi_g(\xi) = \exp \left( -\frac{1}{4} \xi^{\rm T} \Omega^{\rm T} \sigma   \Omega \xi + i \xi^{\rm T} \Omega^{\rm T} \bar{r} \right) \;.
\end{equation}
The Wigner function can be computed as the Fourier transform of the characteristic one,\footnote{i.e., $\mathcal{W}(\xi)= \frac{1}{2^n \pi^{2 n}} \int_{\mathbb{R}^{2 n}} d \xi' {\rm e}^{i \xi'^{\rm T} \Omega \xi} \chi (\xi')$} leading to 
\begin{equation}
\label{eq:wigner_general_gaussian}
    \mathcal{W}_g (\xi) = \frac{2^n}{\pi^n \sqrt{\text{det} \sigma }}{\rm e}^{-\left(\xi - \bar{r} \right)^{\rm T}  \sigma ^{-1} \left(\xi - \bar{r} \right)} \;,
\end{equation}
which is always a real and positive Gaussian if the state is Gaussian (a possible definition of a Gaussian state). 

Crucially, phase space methods allow us to map the problem of evaluating the Markovian time evolution of the quantum state $\hat{\varrho}$ of $n$ infinite-dimensional Hilbert spaces, i.e. $L^2({\mathbbm R^n})$, to a Partial Differential Equation (PDE) for $\chi (\xi)$, equivalent to Schrödinger's for a closed system or to a diffusive master equation under white noise. Seeking conciseness of the expressions, let $\tilde{\xi} = \Omega \xi$. Specifically, from the definition of the characteristic representation of operators, one finds the mappings
\begin{align}
    \partial_{\tilde{\xi}^p} \chi (\tilde{\xi}) &\longleftrightarrow  \frac{i}{2} (\hat{r}^p  \hat{\varrho} + \hat{\varrho} \hat{r}^p ) \;,
    \label{eq:phase_space_corresp1} \\
    \tilde{\xi}^p \chi (\tilde{\xi}) &\longleftrightarrow  \Omega^{qp} (\hat{r}^q \hat{\varrho} -  \hat{\varrho} \hat{r}^q) \;, \label{eq:phase_space_corresp2}
\end{align}
where the superscripts ($p$ and $q$) labels the element of the vectors and matrices, and the Einstein notation is understood (and will be throughout). These (well known) identities and their repetitive used (see Appendix~\ref{app:map_phase_space}) can be used to map unitary and Lindbladian quantum evolutions to PDEs governing the evolution of $\chi$ (and of the Wigner function, which is its Fourier transform). When restricting to the Gaussian dynamics -- that is, to linear or quadratic Hamiltonians and linear couplings to the environment -- of Gaussian states, the PDE has solutions of the form of Eq.~(\ref{eq:chara_general_gaussian}), with a set of ODEs that govern the time evolution of the phase space parameters ($\sigma$ and $\bar{r}$). Such ODEs are often analytically integrable, providing an exact solution to Gaussian quantum dynamics. 

For instance, the most general Gaussian Hamiltonian of $n$ modes is
\begin{align}
\label{eq:Hamiltonian_max}
    \frac{\hat{H}_g}{\hbar \omega} &=   \frac{1}{2} \hat{r}^{\rm T} H_n \hat{r} - r_n^{\rm T} \hat{r}  \;,
\end{align}
where $H_n$ is a $2n \times 2n$ symmetric unitless matrix, capturing the kinetic and local quadratic potential, as well as entangling interactions between the modes, and $r_n$ is a $2n$-dimensional vector representing `forces' (both momentum and position kicks). %, and $\tilde{r} = H_m^{-1} r_m$.  
Consider the time in frequency units, that is, $t \to \tau = \omega t$. Then, the unitary evolution of a Gaussian quantum state $\hat{\varrho}_g$ is given by the von Neumann equation 
\begin{equation}
    \frac{\partial \hat{\varrho}_g}{\partial \tau} = \frac{i}{\hbar \omega} \left[ \hat{\varrho}_g, \hat{H}_g\right] \;. 
\end{equation}
By moving to the phase space representation (Eqs.~\ref{eq:phase_space_corresp1}, \ref{eq:phase_space_corresp2}), the PDE for the time evolution of the characterstic function has solution of the form in Eq.~(\ref{eq:chara_general_gaussian}), with
\begin{equation}
\label{eq:ODE_gaussian}
    \dot{\sigma} = \Omega H_n \sigma - \sigma H_n \Omega \; ,\hspace{0.5cm} \dot{\bar{r}} =  \Omega H_n \bar{r} - \Omega r_n  
    \;,
\end{equation}
where $\dot{A}$ represents the time derivative of quantity $A$. One can solve both equation by defining the symplectic transformation at time $\tau$ generated by $H_n$ as 
\begin{equation}
\label{eq:sympletic_OG}
 S (\tau) = {\rm e}^{\tau \Omega H_n}\;. 
\end{equation}
Then, the solutions to Eq.~(\ref{eq:ODE_gaussian}) for an initial Gaussian state with phase-space parameters $\sigma_0$ and $\bar{r_0}$ are
\begin{align}
\label{eq:cov_gaussian}
\sigma (\tau) &= S(\tau) \sigma_0 S^{\rm T} (\tau) ,\\
     \bar{r}(\tau)  \nonumber
     &= \bar{r}_0(\tau) - ( \tilde{r}_n (\tau) - \tilde{r}_n) \;,
\end{align}
where we defined $\tilde{r}_n =  H_n^{-1} r_n$ and $r_i(\tau) = S (\tau) r_i$ (with $i=0,n$), which solves every quadratic and linear Hamiltonian through a single symplectic transformation. A similar procedure can be applied to Gaussian quantum dynamics that are open (Markovian), by mapping the problem of solving Lindblad-type differential equations to a set of ODEs governing $\bar{r}$ and $\sigma$, (the latter is governed by Lyapunov equations describing decay and diffusive processes). In this instance, the solutions can be expressed always in integral form. Continuous Gaussian monitoring and linear \textit{feedback} of Gaussian states results instead in non-linear Riccati equations for the covariance matrix. Instances of these will arise in unitary dynamics in this work.

%The phase space representation allows one to compute higher order correlations and expectation value of Gaussian in terms of $r$ and $\sigma$ via the partial derivatives of the characteristic function, being the generating function of Gaussian states (see Sec.~\ref{sec:measure} for more details).
Any deterministic Gaussian CP maps -- such as classical mixing, losses, attenuators and amplifiers -- can be expressed via the phase space representation as the map 
\begin{equation}
    \bar{r} \to X \bar{r} + d\hspace{0.5cm} \sigma \to X \sigma X^{\rm T} + Y 
\end{equation}
where $X$ and $Y$ are the two matrices respecting $Y + i \Omega \geq i X \Omega X^{\rm T}$ (to preserve the Heisenberg uncertainty principle) and defining specific maps and $d$ is a displacement.

The phase space representation also allows for the description of measurements often performed in experiments (more details are given in Sec.~\ref{sec:measure}). A general class of Gaussian measurements is represented by the set of general-dyne detections, with Positive Operator Valued Measure (POVM):
\begin{equation}
\label{eq:generadyne}
{\mathds{1}} = \frac{1}{(2\pi)^n} \int_{\mathbbm{R}^{2n}}dr_{o} \hat{D}_{-r_o} \varrho_o \hat{D}_{r_o} ,
\end{equation}
where $n$ is the number of measured modes, $r_{o}$ is the $2n$-dimensional vector of (continuous) measurement outcomes. $\varrho_o$ is a Gaussian state with null first moments and covariance matrix $\sigma_o$: by varying $\sigma_o$ among the physical covariance matrices (i.e., such that $\sigma_o+i\Omega\ge 0$), this class allows one to capture several measurement procedures, such as homodyne and heterodyne measurements, in both ideal and noisy conditions. From the phase-space representation of a Gaussian state, the probability distribution of the outcome $r_o$ of a general-dyne detections is 
\begin{equation}
    \label{eq:prob_generadyne_gaussian}
    P(r_o) =  \frac{{\rm e}^{ - (\bar{r} - r_o)^{\rm T} ( \sigma + \sigma_o )^{-1}  (\bar{r} - r_o) }}{\pi^n  \sqrt{ \text{det} [ \sigma + \sigma_o ] }} \;.
\end{equation}



%The monitoring of linearly-coupled, environmental modes under Markovian, white-noise conditions modifies the standard Lyapunov equation ruling the {\em conditional} state of the open quantum system in the Gaussian regime, turning it in a second-order Riccati equation. It is a remarkable property related to the Gaussian character of the dynamics that the conditional covariance matrix of such monitored states evolves deterministically (i.e., it does not depend on the measurement outcome), whilst the evolution of first moments is stochastic.

\section{ Gaussian-Branched Cat States \\ \& Superpositions of Phase Spaces \label{sec:sup_phase_space}}


The Gaussian phase-space method can be extended to CV-DV hybrid systems by introducing Gaussian-Branched Cat States (GBCSs). In this Section, the main ideas and methodology of this work are presented:  a quantum superposition of phase spaces is used to give the general form of GBCSs. This is laid out for the general case of a $N$-dimensional DV subsystem entangled with $n$ dimensional CV ones, followed by the explicit case of the joint system of a single qubit ($N = 2$) with $n$ modes, to provide a compact example. 

\begin{figure*}[t]
\begin{equation}
\label{eq:density_one_qubit}
    \hat{\varrho} = \frac{1}{(2 \pi)^n} \int_{\mathds{R}^{2n}} \text{d} \tilde{\xi} \;  \hat{D}_{\tilde{\xi}} \otimes \; \begin{pmatrix}
         p_+ \exp \left[-\frac{1}{4} \tilde{\xi}^{\rm T} \sigma^{\text{on}}_+  \tilde{\xi}  +  i  \tilde{\xi}^{\rm T} \bar{r}^{\text{on}}_+   \right] & q  \exp \left[ -\frac{1}{4} \tilde{\xi}^{\rm T} \sigma^{\text{off}}  \tilde{\xi}  +  i  \tilde{\xi}^{\rm T} \bar{r}^{\text{off}}  -   \mathcal{C}  + i  \phi  \right]   \\[8pt]
       q^* \exp \left[ -\frac{1}{4} \tilde{\xi}^{\rm T} \sigma^{\text{off}\;*}  \tilde{\xi}  +  i  \tilde{\xi}^{\rm T} \bar{r}^{\text{off}\;*}  -   \mathcal{C}  - i  \phi  \right]    & p_- \exp \left(-\frac{1}{4} \tilde{\xi}^{\rm T} \sigma^{\text{on}}_-  \tilde{\xi}  +  i  \tilde{\xi}^{\rm T} \bar{r}^{\text{on}}_-   \right) 
       \end{pmatrix}
\end{equation}\caption*{The most general single-qubit by $n$ modes GBCS.\label{fig:density_one_qubit}}
\end{figure*}

Mathematically, the state of a joint system of $n$ CV degrees of freedom and a $N$ dimensional DV is described by the density matrix operator $\hat{\varrho}$ acting on the joint Hilbert space $\mathcal{H}_m \otimes\mathcal{H}_q \sim L^2(\mathds{R}^n) \otimes \mathds{C}^{N} $, where -- here and in the following -- the superscript $m$ and $q$ refers to quantities related to the CV (modes) and DV (qudit) subsystems, respectively. Without loss of generality, given one DV operator $\hat{J}_z$, its eigenstates $\hat{J}_z \ket{j} = j \ket{j}$ -- which spans the DV Hilbert space, with $j \in [ 0, N- 1]$ -- can be used to expand the density matrix of the joint system
\begin{equation}\label{rhojk}
    \hat{\varrho} = \sum_{j, k  = 0}^{N-1}  \hat{\varrho}_{jk} \otimes \ket{j} \bra{k} \;,
\end{equation}
where $ \hat{\varrho}_{jk} = \braket{j|\hat{\varrho}|k}  $ are operators acting on the CV Hilbert space, labelled by the eigenvalue of~$\ket{j}$. The operators $\hat{\varrho}_{jj}$ in diagonal entries are unnormalized density matrices of $n$ modes; however, off-diagonal operators are not Hermitian and, hence, not physical, though the overall hermiticity implies $(\hat{\varrho}_{jk})^\dag = \hat{\varrho}_{kj} $. We associate a phase space with each $ \hat{\varrho}_{jk} $ via the composite Fourier–Weyl relation 
\begin{equation}
\label{eq:density_phase_space}
    \hat{\varrho} = \frac{1}{(2 \pi)^n} \sum_{j, k  = 0}^{N-1} \int_{\mathds{R}^{2n}} \text{d} \xi \; \chi_{jk} (\xi)  \hat{D}_{\xi} \otimes \ket{j} \bra{k} \;,
\end{equation}
where we define the \textit{Branched Characteristic Functions}: \begin{equation}
\label{eq:chara_jk_def}
    \chi_{jk} (\xi) = \text{Tr}_m [ \hat{D}_{- \xi}  \hat{\varrho}_{jk} ] \;.
\end{equation}
Formally, this is possible as the set of operators $\{\hat{D}_{r} \otimes \ket{j} \bra{k} \} $ with $j,k  \in [ 0, N-1] $ and $\xi \in \mathds{R}^{2n}$ is a complete orthonormal basis of the space of operators acting on the CV-DV Hilbert space, i.e. the joint space of trace-class operators $\mathfrak{S} (\mathcal{H}_m \otimes\mathcal{H}_q)$, with respect to the Hilbert--Schmidt product.

 A \textit{Gaussian-Branched Cat State} is a CV-DV entangled state with Gaussian branched characteristic functions:
\begin{equation}
\label{eq:chara_jk}
    \chi_{jk} (\tilde{\xi}) =  \exp \left( -\frac{1}{4} \tilde{\xi}^{\rm T} \sigma_{jk} \tilde{\xi}  +  i  \tilde{\xi}^{\rm T} \bar{r}_{jk} + r^{(0)}_{jk}   \right) \;,
\end{equation}
where $\sigma_{jk}$ are $2n\times 2n$ symmetric matrices, $\bar{r}_{jk}$ are $2n$-dimensional vectors and  $r^{(0)}_{jk} $ are scalars, all labeled by all combinations of $j$ and $k$; note that these parameters will in general be complex. To fix ideas, the most general one qubit by $n$ modes GBCS is shown in Eq.~\ref{eq:density_one_qubit}. We will collectively name them phase-spaces parameters, which fully characterise a GBCS and, rather loosely, we will refer to all $\bar{r}_{jk}$ and $\sigma_{jk}$ as first and second moments, though it should be borne in mind that the off-diagonal parameters do not correspond to physical first and second moments. As we shall see, they are typically complex quantities and, given that $\hat{\varrho}$ is Hermitian, it follows that $\sigma_{jk}^* = \sigma_{kj}$, $\bar{r}_{jk}^* = \bar{r}_{kj}$, and $r_{jk}^{(0)\; *} = r_{kj}^{(0)}$.  

Let us then describe the general properties of GBCS. The matrix $\hat{\chi} (\tilde{\xi} ) = \sum_{jk}\chi_{jk} (\tilde{\xi} ) \ket{j} \bra{k}$ is a Hermitian operator acting on the DV Hilbert space as a function of the phase space variable $\tilde{\xi}$. Of particular interest is the case of $\tilde{\xi} = 0$, which, according to the definition of $\chi_{jk}$ (Eq.~\ref{eq:chara_jk_def}), is the Qudit Reduced Density Matrix (QRDM) $\hat{\varrho}^{q} := \text{Tr}_m [\hat{\varrho}] = \hat{\chi} (0) $, with components
\begin{equation}
     \varrho^{q}_{jk} = \exp(r^{(0)}_{jk}) = \exp \left(- \mathcal{C}_{jk} + i  \phi_{jk} \right),
\end{equation}
where we defined $\mathcal{C}_{jk}$ and $ \phi_{jk} $ as the purely real and imaginary parts of $r^{(0)}_{jk}$, respectively, and we will name them as \textit{contrasts}\footnote{More precisely, $\mathcal{C}_{jk}$ are the exponents of the decay in contrast, which represent sources of decoherence in the QRDM. However, being a dynamical quantity, they include apparent dephasing induced by a superposition in the CV, which can show coherence revival when the CV wavefunction is recombined, see Sec~\ref{sec:physical_examples}, Fig~\ref{fig:example_1}.} and \textit{phases} of the QRDM. Notice that, as will be clear from the final examples, the quantities $\mathcal{C}_{jk}$ provide an expedient and systematic way to identify processes where the qubit's coherence is preserved, spoilt or revived during a superposed Gaussian process. Note that  $\mathcal{C}_{jk} = \mathcal{C}_{kj}$ and $\phi_{jk} = -\phi_{kj}$.

Furthermore, one can find the characteristic function of the modes reduced density matrix, $\hat{\varrho}^m := \text{Tr}_q [\hat{\varrho}] $, such that $\chi^m (\xi) = \sum_{j} \chi_{jj} (\xi)$. On the diagonal terms, the phase space quantises $\sigma_{jj}$ and $r_{jj}$ are real and they respectively represent physical covariance matrices and vectors of first moments, while $\varrho^q_{jj} = {\rm e}^{r_{jj}^{(0)}} = p_{jj}$ are real probabilities. Then, it is possible to interpret $\chi^m$ as the statistical mixture of $n$-mode Gaussian states with first and second moments $r_{jj}$ and $\sigma_{jj}$, weighted according to the probability distribution $\{ p_{jj} \}$ encoded in the diagonal elements of the QRDM. 

It is insightful to explicitly describe the case of $n$-modes with a qubit.  In this case, the DV Hilbert space $\mathds{C}^{2}$ can be represented by a basis $\ket{\pm 1}$, for example, of the Z-Pauli operator $\hat{\sigma}_z \ket{\pm 1} = \pm \ket{\pm 1}$. The most general $n$-mode single-qubit GBCS as given in Eq.~\ref{eq:density_one_qubit}, where the QRDM quantities are, generically, $p_\pm \in \mathds{R}^+$ with $p_+ + p_- = 1$, $q \in \mathds{C}$, and $|q {\rm e}^{-\mathcal{C}}| \leq \sqrt{p_{+} p_{-}} $. Furthermore, we define the quantity $\sigma^{\text{on}}_\pm = \sigma_{\pm1 , \pm 1}$ and $\sigma^{\text{off}} = \sigma_{+ 1 , - 1} =  \sigma_{- 1 , + 1}^*$ as the on- and off-diagonal covariance matrices, and proceed similarly for the first moments $r^{\text{on/off}}$, thus, applying the symmetry imposed by the hermiticity of the joint state.

Note that, in the general case, our formalism is apt to describe several physical systems. As canonically DV variable, this is single large DV physical subsystem (where $N = 2 J +1 $, where $J$ could be the total angular momentum), as well as, an ensemble of $N_q$ qubits, where $N = 2^{N_q}$. For a large spin $J$, for instance, the $\hat{J}_z$ eigenvalues $m \in [-J, J]$ correspond to $j = m + J \in [0, 2J]$ in our notation. Similarly, for $N_q$ qubits, the strings $\vec{s} = (s_1,s_2, \dots, s_{N_q}) \in \{\pm 1\}^{N_q}$ correspond to $j = \sum_{i=1}^{N_q} (1+s_i)\, 2^{\,i-2}$. \textcolor{blue}{It also could describe non-linear purely CV joint systems, where the basis is the fock basis of one of subsystem, for instance, a mechanician resonator coupled to a cavity~\cite{bose_preparation_1997}.}


\section{Dynamics \label{sec:dynamics}}

Any `operator-valued' Gaussian dynamics, which is generated by interactions that are Gaussian Hamiltonians ``labelled" by the DV state, can be analytically solved by an evolving GBCS. In fact, this treatment and state parametrization allows one to map the dynamics of the CV-DV system -- originally described as a Lindblad-type equation for $\hat{\varrho}$ (or $\hat{r}$, in the Heisenberg picture) -- to the time evolution of the phase-spaces parameters of a GBCS. Although this class of Hamiltonian may appear restrictive, this type of interactions arises frequently in experimental practices (as discussed in Sec~\ref{sec:intro}). 

Let us introduce the most general Gaussian Operator-Valued Hamiltonian. Specifically, the ``operator-valued" trait refers to the operational dependence of the Hamiltonian only on a single hermitian DV operator, such that 
$\hat{H}=H (\hat{r}, \hat{J}_z)$, implying that $\hat{H}$ is diagonal in the DV eigenbasis of $\hat{J}_z=\sum_{j}\ket{j}\bra{j}$. In this sense, these are the interaction Hamiltonians that commute with the local Hamiltonian of the DV, i.e. the DV energy is conserved during these dynamics and no transitions occur. Further, we will make the assumption that the Hamiltonian is second order polynomials in $\hat{r}$, giving rise to Gaussian dynamics within each branch of a DV superposition. Thus, the most general operator valued Gaussian Hamiltonian~is 
\begin{equation}
\label{eq:Hamiltonian_max}
    \hat{H}_g (\hat{r}, \hat{J}_z) = \hbar \omega \sum_{j  = 0}^{N-1}\left(  \frac{1}{2} \hat{r}^{\rm T} H_{j} \hat{r} - r_{j} ^{\rm T} \hat{r} +  \frac{1}{2} H_j^0 \right)\otimes \ket{j} \bra{j}     \;,
\end{equation}
where $H_j$ are $N$ real symmetric matrices of dimension $2n \times 2n$, $r_{j}$ are $N$ real $2n$-dimensional vectors, and $H_j^0$ are scalars, all classically labelled by the eigenvalues $j$.

In the specific case of a single qubit, by choosing the qubit operational basis as ${\mathbbm 1}$ and $\sigma_z$, we adopt a different parametrisation such that the general qubit operator-valued Gaussian Hamiltonian $H_g (\hat{r}, \hat{\sigma}_z)$ can be written as
\begin{align}
\label{eq:Hamiltonian_max_qubit}
    \hspace{-0.3cm}\frac{\hat{H}_g (\hat{r}, \hat{\sigma}_z)}{\hbar \omega} &=  \left(  \frac{1}{2} \hat{r}^{\rm T} H_m \hat{r} - r_m^{\rm T} \hat{r}  \right) \otimes \mathds{1}\nonumber \\
    &\hspace{1cm} + \left( \frac{1}{2} \hat{r}^{\rm T} H_q \hat{r} - r_q^{\rm T} \hat{r} + \frac{1}{2} H_q^0 \right) \otimes \hat{\sigma}_z
\end{align}
where $H_m$ and $H_q$ are two $2 n \times 2n $ symmetric matrices describing the modes and operator-valued quadratic potentials, $r_m$ and $r_q$ are $2n$ dimensional vectors, representing forces which can be ``unconditional" or operator-valued by the qubit, and $H_q^0$ is the zero energy splitting of the qubit. This recovers the general form of Eq.~(\ref{eq:Hamiltonian_max}), with $j = \pm 1$ via $H_m  = (H_{+ 1} + H_{-1})/ 2$ and $ H_q  = (H_{+ 1} - H_{-1})/ 2$; and similarly for $r_m$ and $r_q $. 

Let us now lay out how we shall solve the quantum dynamics via GBCS, and thus understand the limitation imposed by  the required assumptions of the Hamiltonian. The diagonality assumption allows us to derive the time evolution of the branched characteristic functions (Eq.~\ref{eq:chara_jk_def}) by mapping the Lindblad-type equations to sets of $2^N$ uncoupled PDEs of $\chi_{jk}$ (which would be coupled if the Hamiltonian were not DV-diagonal). The Gaussianity assumption implies that the PDEs are quadratic in the phase-space variables and their partial derivatives. We finally show that the solution to the uncoupled Gaussian PDEs are the branched characteristic functions of a time-dependent GBCS (Eq.~\ref{eq:chara_jk}).  If either the assumption of diagonality for the qubits or Gaussianity for the CV system is dropped, the dynamics does not have a general analytical solution, as the algebra does not close~\cite{wei_james_1963,legget_dynamics_1987,circuit_niemczyk_2010,superconducting_yoshihara_2017,ralph_dynamical_2018,roda_macroscopic_2024}. In Sec.~\ref{sec:perturbation}, we overcome this limitation by discussing perturbation theory with GBCS.

Specifically, we find that a GBCS is solution of a particular quantum dynamics if and only if the time evolution of its phase-spaces parameters, $\sigma_{jk}(\tau)$, $\bar{r}_{jk}(\tau)$,  and $r^{(0)}_{jk} (\tau)  $,  are solutions of a set of coupled ODEs of the form
\begin{align}
   \dot{\sigma}_{jk} &= f(\sigma_{jk}), \;\;
   \dot{\bar{r}}_{jk} = g(\sigma_{jk}, \bar{r}_{jk}), \;\;
   \dot{r}_{jk}^{(0)} = h(\sigma_{jk}, \bar{r}_{jk}),  \nonumber
\end{align}
 where $f$, $g$ and $h$ are functions depending on the considered type of dynamics. The explicit forms of the sets of ODEs and the integral form of their solutions (when generically derivable) represent one of the central results of this work, presented in Sec.~\ref{sec:unitary} for unitary dynamics and in Sec.~\ref{sec:open} for open dynamics. In Appendices \ref{app:map_phase_space}, \ref{app:unitary_general}, and~\ref{app:noise}, the mapping from Hilbert space dynamics to phase space parameters is employed to derive our dynamical equations.

The set of initial conditions is needed to find the particular solution to describe the evolution of a specific initial quantum state. Let us first consider an initial product state of the form $\hat{\varrho}(0) = \hat{\varrho}_m^{(g)} (0) \otimes \hat{\varrho}^q (0) $, where $\hat{\varrho}_m^{(g)} (0) $ is a general Gaussian state of the $n$ modes and $\hat{\varrho}^q (0)$ is a general qubit density matrix. The initial first and second moments ($r_0$ and $\sigma_0$) of $\hat{\varrho}_m^{(g)} (0) $, with the elements of the initial QRDM ($\varrho^q_{jk} (0)$), form a complete set of initial conditions for a particular solution of the time-evolution ODEs governing the time evolution of the phase space parameters. This choice allows one to describe a large set of experiments, as often the initial state of the modes is indeed Gaussian -- for such are coherent, thermal and squeezed states. In the following, we shall give concrete examples of exact dynamics for product initial states.  Yet, let us point out that the phase-space formulation of the dynamics allows one to consider more general states. If a GBCS is the initial state of a quantum dynamics, its phase-space parameters are still valid as initial conditions ($\sigma_{jk}(0)$, $\bar{r}_{jk}(0)$ and  $r_{jk}^{(0)}(0)$). This allows repetitive use of the dynamics, with the inclusion of intermediate operations and measurements (see Sec.~\ref{sec:operations} and~\ref{sec:measure}).


 %Furthermore, this analysis holds for time-dependent, open, monitored, and with feedback dynamics (see Sec. \ref{sec:open}).

\begin{table*}[!]
    \centering
    \renewcommand{\arraystretch}{1.7}
    \begin{tabular}{ |p{4.0cm}||P{6cm}|P{6cm}|}
        \hline
        Phase Spaces Parameters  & \multicolumn{2}{c|}{ODEs for Unitary Dynamics with Operator-valued Gaussian Interactions} \\ 
        \hline
        Second Moments  & \multicolumn{2}{c|}{\ensuremath{ \dot{\sigma}_{jk} =   \frac{1}{2} \left[  \Omega  ( H_j + H_k )  \sigma_{jk} - \sigma_{jk}  ( H_j + H_k ) \Omega     
  -  i   \sigma_{jk} (H_j - H_k) \sigma_{jk}  + \Omega (H_j - H_k) \Omega \right]}}\\
       First Moments & \multicolumn{2}{c|}{
       \ensuremath{
       \dot{\bar{r}}_{jk} =\frac{1}{2}\left[ \left(  \Omega ( H_j + H_k )  - i   \sigma_{jk} (H_j - H_k) \right)  \bar{r}_{jk} 
     - \Omega (r_j + r_k) + i  \sigma_{jk} (r_j - r_k) \right] }} \\
     QRDM Exponent & \multicolumn{2}{c|}{\ensuremath{ \dot{r}_{jk}^{(0)} =   - \frac{i}{2}  \bar{r}_{jk}^{\rm T} (H_j - H_k)  \bar{r}_{jk}  + i (r_j - r_k)^{\rm T}  \bar{r}_{jk} 
   -  \frac{i}{4} \text{Tr} [(H_j - H_k)\sigma_{jk}] -   \frac{i}{2} (H_j^0 -H_k^0}) } \\
    \hline
    %\hline
    % von Neumann Equation  
    % & \multicolumn{2}{c|}{\ensuremath{
    %   \partial_\tau \hat{\varrho} = i  [ \hat{\varrho}, H (\hat{r},\hat{\sigma}_z) ] }}
    %\\ 
%\hline 
    \end{tabular}
    \caption{ODEs governing the time evolution of the phase space parameters of a GBCS (Eq. \ref{eq:chara_jk}) undergoing a operator-valued Gaussian unitary dynamics with Hamiltonian matrices $H_j$ and force vectors $r_j$ (Eq. \ref{eq:Hamiltonian_max}).}
    \label{tab:unitary}
\end{table*}

\subsection{Unitary Dynamics\label{sec:unitary}}

The unitary evolution of the system is given by the von Neumann equation \begin{equation}
\label{eq:vonNeumann_master}
    \frac{\partial \hat{\varrho}}{\partial \tau} =  \frac{i}{\hbar \omega} \left[ \hat{\varrho} , \hat{H}_g (\hat{r}, \hat{J}_z)  \right] \, .
    \end{equation}
The evolution of the branched $\hat{\varrho}_{jk}$ under this globally unitary dynamics is not trace-preserving nor hermitian (see Appendix~\ref{app:unitary_general}), leading to a vast, rich structure~\cite{Real_Bender_1998,ashida_nonhermitian_2020,third_mcdonald_2023}, analogous to that of generalized master equations~\cite{RevModPhys.70.101,RevModPhys.81.1665,PRXQuantum.5.020201}. Among the class of operator-valued Gaussian Hamiltonians, a main distinction can be made between quadratic and linear qubit-modes interactions, which will be analysed separately in the two sections that follow. 
%The latter are the subset of interaction for which $H_j = H_m \forall j$, which find a close integral form.

%These dynamics, being quadratic and diagonal, are integrable, and, in fact, it is known that if either of these assumptions are not met, the arising quantum dynamics is not analytically solvable like, for instance, in spin-boson models~\cite{legget_dynamics_1987}, ultra-strong coupling~\cite{circuit_niemczyk_2010,superconducting_yoshihara_2017}, and nonlinear CV dynamics \cite{ralph_dynamical_2018, roda_macroscopic_2024}\textcolor{red}{[again non gaussian]}.}

\subsubsection{Quadratic Interactions \label{subsec:quadratic}}

The former is the general case of Eq.~(\ref{eq:Hamiltonian_max})and, by performing the map from Hilbert to superposed phase space (Appendix \ref{app:unitary_general}), we prove that the GBCS with characteristic functions of Eq.~(\ref{eq:chara_jk}) is solution of the the von Neumann equation with Hamiltonian of Eq.~(\ref{eq:Hamiltonian_max}) if and only if $\sigma_{jk} (\tau)$, $\bar{r}_{jk} (\tau)$, and $  \bar{r}_{jk}^{(0)}(\tau) $ are solutions of the ODEs in Table \ref{tab:unitary}. The quadratic coupling induces a complex Riccati term (i.e. $i \sigma (H_j - H_k) \sigma$) in the evolution of the covariance matrix, which does not allows for a general close integral solution. In Sec.~\ref{subsec:riccati}, we provide a systematic way to compute specific analytical solutions of the Riccati equation -- and thus of the dynamics -- via a linear system (at the cost of doubling the variables). It is worth noticing that the ODEs of Table~\ref{tab:unitary} hold for both time-independent and time-dependent Hamiltonians. In fact, this treatment can be trivially extended to time-dependent Hamiltonians by considering $H_j (\tau)$, $r_j (\tau)$, and $H_j^{0} (\tau)$, which are matrices and vectors of time-dependent couplings, while keeping the same form of ODEs. We have shown an application of GBCS to time-dependent Gaussian dynamics in Ref.~\cite{braccini_exponential_2024}. 

Furthermore, one can study steady-state solutions of (stable) hybrid systems: by setting $\dot{\sigma}_{jk}= 0$ and $\dot{\bar{r}}_{jk}= 0$  and solving for the resulting algebraic equation, the covariance matrices and first moments of the steady state can be computed and used to find the form of $\dot{r}^{(0)}_{jk}$. \textcolor{blue}{In this case (when the Hamiltonian and state parameters of the CV sector are not time-dependent), the QRDM is trivially computed as $r^{(0)}_{jk}(\tau) \sim \tau \dot{r}^{(0)}_{jk}$ such that the right-hand-side of $\dot{r}^{(0)}_{jk}$ can be interpreted as encoding rates and decoherence and phases.}

%\begin{align}
%\label{eq:ODE_sigma_jk}
%    \dot{\sigma}_{jk} &=   \frac{1}{2} \left[ \sigma_{jk}  ( H_j + H_k ) \Omega  -   \Omega  ( H_j + H_k )  \sigma_{jk} \right]  \\
%    &\hspace{1cm} -  i  \left[ \sigma_{jk} (H_j - H_k) \sigma_{jk}  + \Omega (H_j - H_k) \Omega \right] \nonumber
%\end{align}
%\begin{align}
%\label{eq:ODE_r_jk}
%    2 \dot{\bar{r}}_{jk} &=  \left[ - \Omega ( H_j + H_k )  + i   \sigma_{jk} (H_j - H_k) \right]  \bar{r}_{jk} \\
%    &\hspace{1cm} +  \Omega (r_j + r_k) + i  \sigma_{jk} (r_j - r_k) \nonumber
%\end{align}
%\begin{align}
%\label{eq:ODE_r_0_jk}
%     \dot{\bar{r}}_{jk}^{(0)} &=   \frac{i}{2}  \bar{r}_{jk}^{\rm T} (H_j - H_k)  \bar{r}_{jk}  - i (r_j - r_k)^{\rm T}  \bar{r}_{jk} \\
%    &\hspace{1cm} +  \frac{i}{2} \text{Tr} [(H_j - H_k)\sigma_{jk}]+  \frac{i}{2} H_q^0(j-k) \nonumber
%\end{align}
 




%When restricting to interactions on the $n$ modes and a qubit which is Gaussian and diagonal in the computational basis -- up to any rotation of the basis of the qubit-- each PDE can be solved. For not diagonal interaction, for instance the Boson model, it is known that the dynamics is generically not integrable, and truncations are needed. As we shall see, the quantum dynamics can be rewritten as a set of ordinary differential equations governing the time evolution of two parameters which acts as a generalisation of the first and second moments, as well as the reduced density matrix of the qubit.


\begin{table*}[!]
    \centering
    \renewcommand{\arraystretch}{1.7}
    \begin{tabular}{ |p{4.0cm}||P{6cm}|P{6cm}|}
    \hline
        Phase Spaces Parameters & \multicolumn{2}{c|}{Solution of Unitary Dynamics with Operator Valued Linear Interactions} \\ 
        \hline
        \hline
        Covariance Matrices   & \multicolumn{2}{c|}{\ensuremath{
   \sigma (\tau) := \sigma_{jk}(\tau)  = S_m (\tau) \sigma_0 S_m^{\rm T} (\tau)}}\\
       First Moments & \multicolumn{2}{c|}{
       \ensuremath{\bar{r}_{jk} (\tau) = r_0(\tau) - \frac{1}{2}  \left( \tilde{r}_j(\tau) - \tilde{r}_j + \tilde{r}_k(\tau) - \tilde{r}_k  \right) -  \frac{i}{2} \sigma (\tau)   \Omega  \left[ (\tilde{r}_j (\tau) - \tilde{r}_j)  - (\tilde{r}_k (\tau) - \tilde{r}_k ) \right]   }} \\
     QRDM Contrasts & \multicolumn{2}{c|}{
       \ensuremath{\mathcal{C}_{jk} (\tau) = \frac{1}{4} \left[ (\tilde{r}_j (\tau) - \tilde{r}_j)  - (\tilde{r}_k (\tau) - \tilde{r}_k ) \right]^{\rm T} \Omega^{\rm T} \sigma (\tau)   \Omega  \left[ (\tilde{r}_j (\tau) - \tilde{r}_j)  - (\tilde{r}_k (\tau) - \tilde{r}_k ) \right]  }}\\
    QRDM Phases & \multicolumn{2}{c|}{\ensuremath{ \phi_{jk} (\tau) = - (\tilde{r}_j - \tilde{r}_k)^{\rm T} \Omega \left[  ( \tilde{r}_0 (\tau) - \tilde{r}_0) -  \frac{1}{2} \left( \tilde{r}_j (\tau) - \tilde{r}_j + \tilde{r}_k (\tau) -  \tilde{r}_k  \right) \right]  }} \\
 & \multicolumn{2}{c|}{\ensuremath{ \hspace{2cm}  + \frac{\tau}{2} \left[ (\tilde{r}_j - \tilde{r}_k)^{\rm T} H_m (\tilde{r}_j + \tilde{r}_k) - (H_j^{0}-H_k^{0})\right] }} \\
    \hline
    \end{tabular}
    \caption{Time evolution of the phase space parameters of a GBCS (Eq.~\ref{eq:chara_jk}) undergoing unitary evolution with operator-valued forces (Eq.~\ref{eq:Hamiltonian_force_only}), where $\tilde{r}_a = H^{-1}_m r_a$, $S_m (\tau) = \exp (  \Omega H_m \tau)$, and $\tilde{r}_a (\tau) = S_m (\tau) \tilde{r}_a$, with $a \in 0,j,k$. The initial state is a $n$-modes Gaussian state with first and second moments $r_0$ and $\sigma_0$ and a general QRDM $\varrho^{q}_{jk}(0)$.}
    \label{tab:force}
\end{table*}

For clarity, let us explicitly write the ODEs (and when possible their solution) for the on- and off-diagonal terms in the single qubit case, thus providing one with the general solution of the dynamics in terms of the state given by Eq.~(\ref{eq:density_one_qubit}). On the diagonal elements, i.e. $\ket{\pm 1}\bra{\pm 1}$,  the ODEs take the form 
\begin{align}
\label{eq:general_unitary_ODE_on}
    \dot{\sigma}^{\text{on}}_{\pm} &=  \Omega  (H_m \pm H_q)  \sigma^{\text{on}}_{\pm} -  \sigma^{\text{on}}_{\pm} ( H_m \pm H_q) \Omega   \nonumber \;,  \\ 
    \dot{\bar{r}}^{\text{on}}_{\pm} &=  \Omega  (H_m \pm H_q)   \bar{r}^{\text{on}}_{\pm} - \Omega (r_m \pm r_q) \;, 
\end{align} 
and the diagonal elements of the QRDM are unchanged. Eq.~(\ref{eq:general_unitary_ODE_on}) represents two Gaussian dynamics of the same initial Gaussian state under two Gaussian Hamiltonians labelled with $\pm$ (as per Eq.~\ref{eq:ODE_gaussian}). Analogously to Sec.~\ref{sec:intro_phase}, the general solution to these equations (see Appendix~\ref{app:only_force}) can be compactly written by defining the symplectic transformations
\begin{equation}
\label{eq:symplectic_pm}
    S_{\pm}(\tau) = \exp( \tau \Omega (H_m \pm H_q)  ) \;,
\end{equation}
as well as the four vectors $r_i^{\pm} = (H_m \pm H_q)^{-1} r_i$ and their time evolution $ r_i^{\pm} (\tau) =S_{\pm}(\tau)  r_i^{\pm}$, with $i = \{m, q\}$. Then, a Gaussian state with initial moments $r_0$ and $\sigma_0$ and a qubit with density matrix element $\varrho_{jk}^q(0)$ evolve to 
\begin{align}
\label{eq:time_diagonal_general}
    \sigma^{\text{on}}_{\pm} (\tau) &= S_\pm (\tau) \sigma_0 S^{\rm T}_\pm (\tau), \\
     \bar{r}^{\text{on}}_{\pm 1}(\tau) &=   S_\pm (\tau) r_0 - \left[ (r_m^{\pm}(\tau) - r_m^{\pm}) \pm (r_q^{\pm}(\tau) - r_q^{\pm}) \right], \nonumber 
\end{align}
while the QRDM remains constant. As expected, the on-diagonal terms of a GBCS evolution are described by a statistical mixture of Gaussian processes which preserve the Gaussianity and normalization of the states associated with the diagonal element of the density matrix. 

For the off diagonal terms, the ODEs read
\begin{align}
\label{eq:general_unitary_ODE_off}
     \dot{\sigma}^{\text{off}} &=   \Omega  H_m  \sigma^{\text{off}} -  \sigma^{\text{off}}    H_m   \Omega   -   i  \left(  \sigma^{\text{off}} H_q \sigma^{\text{off}}  + \Omega H_q  \Omega \right) \;, \nonumber \\
     \dot{\bar{r}}^{\text{off}} &=  \left( \Omega  H_m  -  i   \sigma^{\text{off}} H_q \right) \bar{r}^{\text{off}}   -  \Omega r_m + i \sigma^{\text{off}} r_q \;,   \\
     \dot{\mathcal{C}} &= \,\text{Im}\!\left[ (\bar{r}^{\text{off}})^{\rm T} H_q \bar{r}^{\text{off}} + \frac{1}{2}\text{Tr}\!\left(H_q \sigma^{\text{off}}\right)  - 2 r_q^{\rm T} \bar{r}^{\text{off}} \right] \;, \nonumber \\
     \dot{\phi} &= -\,\text{Re}\!\left[ (\bar{r}^{\text{off}})^{\rm T} H_q \bar{r}^{\text{off}} + \frac{1}{2}\text{Tr}\!\left(H_q \sigma^{\text{off}}\right)  - 2 r_q^{\rm T} \bar{r}^{\text{off}} \right] - H_q^0 \;. \nonumber
\end{align}
\textcolor{red}{removed minus sign from $\dot{C}$.}

Let us note that the off-diagonal phase-space parameters indeed evolve to complex values: the covariance matrix via the nonlinear Riccati term ($\sigma^{\text{off}} H_q \sigma^{\text{off}}$) and the first moments via the coupling of the covariance matrix (as well as the QRDM). These features are also proper to the general case for $j\neq k$, as per Table~\ref{tab:unitary}.  %\textcolor{red}{In Sec. XXX, an example of this dynamics is numerically solved to study entanglement between two qubits interacting via a resonator, usually how, in many technology, CNOTs are implemented.}



\subsubsection{Linear Interactions \label{subsec:force}}

Let us now consider the most general Hamiltonian involving a purely linear coupling between CVs and DVs:
\begin{equation}
\label{eq:Hamiltonian_force_only}
    \frac{\hat{H}_l (\hat{r}, \hat{J}_z)}{\hbar \omega} = \frac{1}{2} \hat{r}^{\rm T} H_{m} \hat{r} - \sum_{j = 0}^{N -1}\left(  r_{j} ^{\rm T} \hat{r} -  \frac{1}{2} H_j^0 \right)\otimes \ket{j} \bra{j}     \;,
\end{equation}
whose associated ODEs have analytical solutions that can always be written in integral form by a single transformation governing the quadratic evolution of all superposed phase spaces. From Table \ref{tab:unitary}, one finds the simplified ODEs,
\begin{align}
    \dot{\sigma}_{jk} &=  \Omega  H_m  \sigma_{jk} -  \sigma_{jk}  H_m \Omega \;, \nonumber  \\
    \dot{\bar{r}}_{jk} &=     \Omega  H_m  \bar{r}_{jk}  -   \frac{1}{2} \Omega (r_j + r_k) - \frac{i}{2}  \sigma_{jk} (r_j - r_k) \;, \nonumber  \\
    \dot{\bar{r}}_{jk}^{(0)} &=   i  (r_j - r_k)^{\rm T}   \bar{r}_{jk}   -  \frac{i}{2} (H_j^0-H_k^0) \;.
\end{align}
As shown in Appendix~\ref{app:only_force}, by defining and computing the single symplectic transformation $S_m (\tau) = \exp(  \Omega H_m \tau)$, the solution of these ODEs can be expressed in integral form as in Table \ref{tab:force}, for an initial Gaussian state with first and second moments $r_0$ and $\sigma_0$ and an initial QRDM $\varrho_{jk}^q (0)$. 
Observe that all covariance matrices of the density matrix are trivially the same , i.e. $\sigma_{jk}(\tau)  = \sigma (\tau) = S_m (\tau) \sigma_0 S_m^{\rm T} (\tau) \; \forall j,k$: given that the coupling is only linear, their evolution is given by a Gaussian dynamics of the form (\ref{eq:cov_gaussian}), as expected. 

Let us define the vectors $\tilde{r}_i =  H_m ^{-1} r_i$ and their time evolution $ \tilde{r}_i (\tau) =S_m(\tau)  \tilde{r}_i$,  with $i = \{j,k,m, q, 0\}$. For the single-qubit case of Eq.~(\ref{eq:density_one_qubit}), the first moments evolve according to
\begin{align}
\label{eq:r_on_force_sol}
     \bar{r}^{\text{on}}_{\pm}(\tau) &=   r_0 (\tau) -  (\tilde{r}_m(\tau) - \tilde{r}_m) \mp (\tilde{r}_q(\tau) - \tilde{r}_q), \\ 
     \bar{r}^{\text{off}} (\tau) &= r_0(\tau) -  \left( \tilde{r}_m(\tau) - \tilde{r}_m \right) - i \sigma (\tau)   \Omega  \left( \tilde{r}_q (\tau) - \tilde{r}_q \right)  , \nonumber
\end{align}
where the former represents the two Gaussian evolutions with different forces, while the latter is a complex quantity, which includes the covariance matrix itself. The time evolution of the QRDM is then given by 
\begin{equation}
\label{eq:QRDM_1}
    \hspace{-0.2cm}\varrho_q (\tau) = \begin{pmatrix}
        p_+ & q {\rm e}^{ - \mathcal{C} (\tau) + i  \phi (\tau) } \\
       q^* {\rm e}^{ - \mathcal{C} (\tau) - i  \phi(\tau)}   & p_- \\ 
    \end{pmatrix} \;.
\end{equation}
where the contrasts and phases at time $\tau$ read
\begin{align} 
\label{eq:phi_c_force_sol}
     \mathcal{C} (\tau) &=   \left(\tilde{r}_q (\tau) - \tilde{r}_q \right)^{\rm T} \Omega^{\rm T}  \sigma (\tau)   \Omega  \left(\tilde{r}_q (\tau) - \tilde{r}_q \right) \nonumber \;,\\
      \phi (\tau) &=  \tau (2 \tilde{r}_q H_m \tilde{r}_m - H_q^0) \\
      & \hspace{0.5cm}- 2 r_q^{\rm T} \Omega \left[ \tilde{r}_0 (\tau) - \tilde{r}_0 - ( \tilde{r}_m (\tau) - \tilde{r}_m ) \right]   \nonumber \;,
\end{align}
respectively. Since $\sigma(\tau)>0$ one may infer that $ \mathcal{C} (\tau) =0$ if, and only if, $\tilde{r}_q (\tau) = \tilde{r}_q$, i.e., if the first moments go back to the initial value, implying a full revival of qubit coherence when $S_m(\tau) = \mathds{1}$~\cite{braccini_exponential_2024}. Furthermore, note that the coherent phase is independent of the covariance matrix and only senses the forces. 
This solves in integral form the most general unitary dynamics of a Gaussian state evolving in Gaussian potentials and linearly coupled to a qubit.


\begin{table*}[!]
    \centering
    \renewcommand{\arraystretch}{2}
    \begin{tabular}{ |p{3.8cm}|P{7cm}|P{7cm}|}
        \hline
        Phase Spaces Parameters  & \multicolumn{2}{c|}{Solution of Open Dynamics with Operator-valued Gaussian Interactions} \\ 
        \hline
        \hline
   \multicolumn{3}{|c|}{(a) Gaussian Interactions ODEs}  \\
    \hline
        Covariance Matrices   & \multicolumn{2}{c|}{\ensuremath{ \dot{\sigma}_{jk} =   \frac{1}{2} \left[  \Omega  ( H_j + H_k +2 E)  \sigma_{jk} - \sigma_{jk}  ( H_j + H_k +2 {\rm E}^{\rm T}) \Omega  -  i   \sigma_{jk} (H_j - H_k) \sigma_{jk}  + \Omega (H_j - H_k) \Omega \right]  + D}}\\
       First Moments & \multicolumn{2}{c|}{
       \ensuremath{
       \dot{\bar{r}}_{jk} =\frac{1}{2}\left[ \left(  \Omega ( H_j + H_k +2 E)  - i   \sigma_{jk} (H_j - H_k) \right)  \bar{r}_{jk} + 
      \Omega (2 d - r_j - r_k) + i  \sigma_{jk} (r_j - r_k) \right] }} \\
     QRDM Exponent & \multicolumn{2}{c|}{\ensuremath{ \dot{r}_{jk}^{(0)} =   - \frac{i}{2}  \bar{r}_{jk}^{\rm T} (H_j - H_k)  \bar{r}_{jk}  + i (r_j - r_k)^{\rm T}  \bar{r}_{jk} 
   -  \frac{i}{4} \text{Tr} [(H_j - H_k)\sigma_{jk}]   -   \frac{i}{2} (H_j^0 -H_k^0)  - \frac{1}{2} (\Gamma_z^{jj} + \Gamma_z^{kk} ) + \Gamma_z^{jk}} } \\
   \hline 
   \multicolumn{3}{|c|}{(b) Linear Interactions Integral Forms}  \\
    \hline
    \hline
        Covariance Matrix   & \multicolumn{2}{c|}{\ensuremath{ \sigma(\tau) = S_A(\tau) \sigma_0 S_A^{\rm T}(\tau) +  \int_{0}^{\tau} \text{d}t S_A(\tau - t) D S^{\rm T}_A(\tau - t) }}\\
       Vectors or First Moments & \multicolumn{2}{c|}{
       \ensuremath{\; \bar{r}_{jk} (\tau) =  S_A (\tau) r_0 +   \left(S_A(\tau) - \mathds{1} \right)  \left[ d -  \frac{1}{2}    (r_j + r_k)\right] - 
    \frac{i}{2} S_A (\tau) \sigma_0 \left( S_A^{\rm T} (\tau) - \mathds{1}  \right) \Omega (A^{\rm T})^{-1}  (r_j - r_k) }} \\
    & \multicolumn{2}{c|}{
       \ensuremath{\hspace{2cm} - \frac{i}{2}  \int_0^{\tau} \text{d}t'  S_A(\tau- t') D \left( S^{\rm T}_A(\tau - t') - \mathds{1} \right) \Omega (A^{\rm T})^{-1} (r_j - r_k)}} \\
     QRDM Contrasts  & \multicolumn{2}{c|}{ \hspace{-0.4cm} \ensuremath{ \mathcal{C}_{jk} (\tau)  =  \frac{1}{4} (r_j - r_k)^{\rm T}  A^{-1} \Omega^{\rm T} \left( S_A(\tau) - \mathds{1} \right) \sigma_0 \left(S^{\rm T}_A(\tau) - \mathds{1} \right) \Omega  (A^{-1} )^{\rm T} (r_j - r_k) + \tau \left[\frac{1}{2} (\Gamma_z^{jj} + \Gamma_z^{kk} ) - \Gamma_z^{jk} \right] }} \\
 & \multicolumn{2}{c|}{\ensuremath{ \hspace{1.4cm}  + \frac{1}{4}(r_j - r_k)^{\rm T}  A^{-1}    \Omega^{\rm T} \left[  \int_0^\tau \text{d} t'    \left(S_A(\tau - t') - \mathds{1} \right) D \left(S_A^{\rm T}(\tau - t') - \mathds{1} \right)   \right]  \Omega  (A^{-1} )^{\rm T} (r_j - r_k) }} \\ 
     QRDM Phases  & \multicolumn{2}{c|}{ \hspace{-3.6cm} \ensuremath{ \;  \phi_{jk} (\tau) =  - (r_j - r_k)^{\rm T} A^{-1} \Omega \left( S_A (\tau)  - \mathds{1} \right) \left[ r_0  + A^{-1} \left( d  - \frac{1}{2} (r_j + r_k) \right) 
    \right]}} \\
 & \multicolumn{2}{c|}{\ensuremath{ \hspace{1.1cm}  + \tau \left[ (r_j - r_k)^{\rm T} A^{-1} \left(d - \frac{1}{2} (r_j + r_k) \right) -  \frac{1}{2} (H_j^0 -H_k^0) \right] }} \\
 \hline%   \hline
%   Master Equation  
%     & \multicolumn{2}{c|}{\ensuremath{
%       \partial_\tau \hat{\varrho} =  i  [ \hat{\varrho}, H_g (\hat{r}, \hat{\sigma}_z)] + i [ \hat{\varrho}, d^{\rm T} \hat{r}]    + \sum_{m, n}  \left(\frac{1}{2} \Omega^{\rm T} D \Omega - i  E \right)^{mn} \left(  \hat{r}^m \hat{\varrho} \hat{r}^n  - \frac{1}{2} \{ \hat{\varrho}, \hat{r}^m  \hat{r}^n \} \right)  +  \frac{\Gamma_z}{2} \left(\hat{\sigma}_z \hat{\varrho} \hat{\sigma}_z  - \hat{\varrho} \right) }} \\ 
%\hline 
    \end{tabular}
    \caption{(a) ODEs governing the time evolution of the phase space parameters of a GBCS (Eq. \ref{eq:chara_jk}) undergoing a operator-valued Gaussian open dynamics with Hamiltonian matrices $H_j$, force vectors $r_j$, driving term $d$, decay matrix $E$, diffusion matrix $D$, and dephasing rates $\Gamma_z^{jk}$ (Eq. \ref{eq:linbladian}). The additional terms from the unitary case Table \ref{tab:unitary} are in blue. (b) time evolution of the phase space parameters of a GBCS (Eq.~\ref{eq:chara_jk}) undergoing open Gaussian evolution with operator-valued forces: where $A = H_m + E$ and $S_A (\tau) = \exp (  \Omega A \tau)$. The initial state is a $n$-modes Gaussian state with first and second moments $r_0$ and $\sigma_0$ and a general QRDM $\varrho^{q}_{jk}(0)$. $\delta_{jk}$ is the Kronecker delta.}
    \label{tab:open}
\end{table*}
 
\subsubsection{No Interactions: Gaussian Dynamics of GBCS}

In experimental practice and designing protocols, after the creation of a GBCS via, for instance, a conditional operation (discussed in Sec.~\ref{sec:operations}), it is often convenient to evolve the superposed state via a Gaussian dynamics. This subset of GBCS dynamics, which can be trivially included in the formalism, is the one with no interaction between the CV and DV subsystems and Hamiltonian
\begin{equation}
\label{eq:Hamiltonian_no_int}
    \frac{\hat{H}_g (\hat{r}) \otimes \hat{H}_q(\hat{J}_z)}{\hbar \omega} = \frac{1}{2} \hat{r}^{\rm T} H_{m} \hat{r} - r_{m} ^{\rm T} \hat{r}  + \frac{1}{2} \sum_{j = 0}^{N - 1}   H_j^0  \otimes \ket{j} \bra{j} \nonumber    \;,
\end{equation}
In this case, each branch of the superposition evolves according to a single Gaussian dynamics described by Eq.~(\ref{eq:cov_gaussian}) under different initial conditions $\sigma_{jk}^{0}$ and $\bar{r}_{jk}^{0}$, depending on the branch. The QRDM builds the phases $\phi_{jk} = i \tau/2 (H_j^{0} - H_k^{0})$.

%Although only tangential to the subject matter of this paper, notice also that the solutions of Table~\ref{tab:force} and, for open systems, \ref{tab:open}(b), also yield the evolutions of pure CV cat-like states (i.e., superpositions of Gaussian states), obtained by setting ${r}_j={r}$ for all $j$.


\subsection{Open Quantum Dynamics \label{sec:open}}

Our formalism requires little adjustment to include effects of open quantum dynamics, presented in this section. The most general Markovian and Gaussian operator-valued open dynamics with diagonal decoherence in the qubit coupling basis is described by the Lindblad equation  
\begin{align}\label{eq:linbladian}
    \hspace{-0.1cm} \frac{\partial \hat{\varrho}}{\partial \tau} &=  \frac{i}{\hbar \omega} [ \hat{\varrho}, \hat{H}_g (\hat{r}, \hat{J}_z)] +  i [ \hat{\varrho}, d^{\rm T} \hat{r}]   + \mathcal{D}[\hat{r}] (\hat{\varrho})  + \mathcal{D}[\hat{J}_z] (\hat{\varrho}) ,
\end{align}
where 
\begin{align}
    \mathcal{D}[\hat{r}] (\hat{\varrho})  &= \sum_{a, b = 1}^{2 n}  B^{a b} \left(  \hat{r}^a \hat{\varrho} \hat{r}^b  - \frac{1}{2} \{ \hat{\varrho}, \hat{r}^a  \hat{r}^b \} \right), \label{lindo}  \\
    \mathcal{D}[\hat{J}_z] (\hat{\varrho}) &= \sum_{j, k  = 0}^{N-1} \Gamma_z^{jk} \left(\hat{P}_{j} \hat{\varrho} \hat{P}_{k}   - \frac{1}{2} \{ \hat{\varrho}, \hat{P}_{j} \hat{P}_{k} \}  \right),\label{linda} 
\end{align}
where $B$ is a $2n \times 2n$ dimensionless matrix of coefficients representing the strength of the noise in the $n$-modes~\cite{serafini_quantum_2017}, $d$ is the driving, $\Gamma_z^{jk}$ are the unitless DV decay rates and $\hat{P}_{j} = \ket{j}\bra{j}$~\cite{PhysRevLett.114.010501}?. The additional Lindbladian terms of Eq.~(\ref{eq:linbladian}) encode all the noises that are diagonal in the DV subsystem and which have jump-operators linear in the CV canonical operators ($\hat{r}$). They often represent experimental leading-order noises in CV-DV systems, such as cavity decay, CV diffusion, and qubit dephasing. 

To ensure that $\hat{\varrho}$ is an Hermitian operator under such a dynamics, it follows that also $B$ has to be Hermitian, and can be, in general, rewritten as 
\begin{equation}
    B = \frac{1}{2} \Omega^{\rm T} D \Omega - i  E \; ,
\end{equation}
where $D$ and $E$ are symmetric and antisymmetric $2n \times 2n$ matrices representing, respectively, the diffusion and drift of the dynamics. We explicitly note that if the Lindbladian terms are known in the ladder operator basis (with noise matrix $B_a$ and driving $d_a$), one can derive the noise matrix in the canonical operators basis as $B = U^{\dag} B_a U$ and $d = U^{\dag} d_a$, where $U$ is given in Eq.~(\ref{eq:change_basis}).  At the fundamental and microscopic level, these terms arise via the linear interaction between the system's modes and a Markovian bath of a large number of modes with first and second moments given by $r_{\text{in}}$ and $\sigma_{\text{in}}$. Given the coupling Hamiltonian between the system and the bath $\hat{r}_s H_c \hat{r}_{\text{in}}$, from the Hamiltonian matrix $H_C$, it is possible to derive $D = \Omega H_C \sigma_{\text{in}} H_C^{\rm T} \Omega^{\rm T}$, $E =H_C \Omega H_C^{\rm T}/2$, and $d = H_C r_{\text{in}}$, which imply $D \geq 0$, since $\sigma_{\text{in}} \geq 0$~\cite{serafini_quantum_2017}. 

In Appendix~\ref{app:noise}, the additional noisy terms are added to the ODEs for the unitary dynamics (Table \ref{tab:unitary}), to find that the phase-space parameters of a GBCS undergoing a general open Markovian operator-valued Gaussian dynamics evolve according to Table \ref{tab:open}. First of all, we note that, from the form of the coupling between the sets of ODEs, the DV noise does not affect the dynamics of the CV system. Yet, the noise in the modes changes the time evolution of the QRDM, as $\dot{\bar{r}}_{jk}^{(0)}$ has an explicit dependence on $\sigma_{jk}$ and $\bar{r}_{jk}$. Specifically, the differences with respect to the unitary case due to CV noise are: (a) the inclusion of an asymmetric part in the Hamiltonian (typical of decays, embodied by $E$); (b) a linear increase in the covariance matrices (due to the diffusive term including $D$); and (c) the linear term in the first moments (the driving term $d$). 

The DV noise leads to a decay with rate $\gamma_{jk} := \frac{1}{2} (\Gamma_z^{jj} + \Gamma_z^{kk} ) - \Gamma_z^{jk} $. It may be insightful to state the explicit form of DV noise for physical systems:
(i)~for a single qubit dephasing, described by the jump operator $\sqrt{\Gamma_z/2} \hat{\sigma}_z$, has $\Gamma_z^{jk} = \Gamma_z/2 (2j-1) (2k- 1)$ and $\gamma_{jk}  = \Gamma_z (j - k)^2$, where we recall that the unitless $\Gamma_z  = 1/(\omega T_2)$, i.e., the inverse of the dephasing time in frequency units; (ii)~for a large spin $J$ dephasing, with the jump operator $\sqrt{2 \Gamma_z} \hat{J}_z$, which implies $\Gamma_z^{jk} = 2 \Gamma_z (j-J) (k- J)$ and $\gamma_{jk}  = \Gamma_z (j - k)^2$; and (iii)~for $N_q$ qubits, with $N_q$ jumps operators $\sqrt{\Gamma_{z}^{i}/2} \hat{\sigma}^{i}_z$ one finds the update to the QRDM via 
\begin{equation}
    e^{-\gamma_{jk}\tau}
=
\left[
\bigotimes_{i=1}^{N_q}
\begin{pmatrix}
1 & e^{-\Gamma_z^{(i)}\tau} \\[4pt]
e^{-\Gamma_z^{(i)}\tau} & 1
\end{pmatrix}
\right]_{jk} \; .
\end{equation}


Remarkably, when restricting to operator-valued linear interactions (Eq.~\ref{eq:Hamiltonian_force_only}), the solution to the open dynamics can be expressed in integral form as in the case of unitary dynamics. In fact, being the coupling only linear, one can define a single, non-symmetric matrix $A = H_m + E$, which governs the quadratic evolutions of all phase spaces, yielding under exponentiation a (non symplectic) transformation $S_A(\tau) = \exp( \Omega A \tau)$~\cite{serafini_quantum_2017}. Thus, as in the case of unitary evolutions, the covariance matrices $\sigma_{jk}$ evolve all according to the Lyapunov equation
\begin{equation}
\label{eq:Lyapunov_sol}
    \dot{\sigma} =  \Omega A \sigma - \sigma A^{\rm T} \Omega + D 
\end{equation}
The general integral form of its solution, along side $\bar{r}_{jk}(\tau)$, $\mathcal{C} (\tau)$, and $\phi (\tau)$ are given in Table~\ref{tab:open}, such that the branched characteristic functions of Eq.~(\ref{eq:chara_jk}) are solutions of the open quantum dynamics under any linear operator-valued interaction. Furthermore, for $E=0$, another simplification can be found in  Tab~\ref{tab:open_app} (and derived Appendix~\ref{tab:open_app}), where we note that the noise does not affect the phases (as they do not depend on $D$). 

%The on and off diagonal first moments are given by
%\begin{align}
%\label{eq:first_diffused}
%     \bar{r}^{\text{on}}_{\pm}(\tau) &=   \tilde{r}_0 (\tau)  - \left[ (\tilde{r}_m(\tau) - \tilde{r}_m) \pm (\tilde{r}_q(\tau) - \tilde{r}_q) \right]  \\
%     \bar{r}^{\text{off}} (\tau) &= r_0(\tau) -  \left( \tilde{r}_m(\tau) - \tilde{r}_m \right) + i \sigma (\tau)   \Omega  \left( \tilde{r}_q (\tau) - \tilde{r}_q \right)  \nonumber \\
%     & \hspace{0.5cm} + i  \Omega \int_0^{\tau} \text{d}t  D (t - \tau) \left[ \tilde{r}_q (\tau - t) - \tilde{r}_q(\tau) \right]
%\end{align}
%where 
%\begin{equation*}
%    D (\tau) =  S_A(\tau) D S_A^{\rm T}(\tau) \;, \hspace{0.5cm} 
%    \tilde{r}_i (\tau) = S_A (\tau) (A^{-1})^{\rm T} r_i \;,
%\end{equation*}
%with $i = m,q$ \textcolor{red}{I am so scared of the transpose}. The %expression for the phase is unchanged (since the real part of $\bar{r}^{\text{off}} %(\tau)$ is the same as the unitary dynamics), but with the replacement $H_m %\to A$. Instead, the contrasts includes an additional term
%\begin{align*}
%     \mathcal{C} (\tau)  &= r_q^{\rm T}   \Omega^{\rm T} \left( S_A(\tau) - \mathds{1} \right) \sigma_0 \left(S^{\rm T}_A(\tau) - \mathds{1} \right) \Omega   r_q  \\
%     &\;\;\; - r_q^{\rm T}    \Omega^{\rm T} \left[  \int_0^\tau \text{d} t'    \left(S_A(\tau - t') - \mathds{1} \right) D \left(S_A^{\rm T}(\tau - t') - \mathds{1} \right)   \right]  \Omega r_q
%\end{align*}
%This solves exactly any operator-valued force dynamics with Gaussian noise. 
%\begin{align}
%    &= r_0(\tau) - \frac{1}{2}     \left(S_A(\tau) - \mathds{1} \right) A^{-1}  (r_j + r_k) - 
%    \frac{i}{2} S_A (\tau) \sigma_0 \left( S_A^{\rm T} (\tau) - \mathds{1}  \right) \Omega (A^{\rm T})^{-1}  (r_j - r_k)   \\
%    &\hspace{1cm} - \frac{i}{2}  \int_0^{\tau} \text{d}t'  S_A(\tau- t') D \left( S^{\rm T}_A(\tau - t') - \mathds{1} \right) \Omega (A^{\rm T})^{-1} (r_j - r_k)
%\end{align}


%\subsection{Gaussian Dynacmics of Non-Gaussian States}

%One may be interested in solving the gaussian dynamics on the modes of an initial non gaussian state entangled with a qubit without direct qubit-modes interactions. In this case, it is

\subsection{Linear Solutions to Riccati Equations \label{subsec:riccati}}

The Riccati term present in quadratic coupling is the only term that, in this work, does not allow for a general solution in integral form for all the phase space parameters. Thankfully, Riccati equations are quite common, especially in control theory of classical~\cite{Kalman_new_1960,AbouKandil2003} and quantum systems~\cite{wiseman_optimal_2005,wiseman_quantum_2010,serafini_quantum_2017} (as well as in time-dependent and dissipative quantum systems~\cite{CRUZ201544,CRUZ2016609,Schuch2018}),
% as well as theatrical analysis of quantum systems\textcolor{red}{\cite{}}; thus, providing a 
so that a few available approaches to their solution exist~\cite{Bini2012,blanco_interplay_2019}. Here, we discuss systematic paths to find solutions of the Riccati equation appearing in Tab~\ref{tab:unitary}~and~\ref{tab:open}, by solving an equivalent linear system for twice the number of variables (see Appendix~\ref{app:doubling} for proof and details). 

The map from a Riccati equation to a linear system of ODE is achieved by consider the ansatz for the covariance matrices of the form $\sigma_{jk} = \mathds{Y}_{jk} \mathds{X}_{jk}^{-1}$, where $\mathds{X}_{jk} $ and $ \mathds{Y}_{jk}$ are two time dependent $2n$-dimensional matrices. One then finds that $\sigma_{jk}$ is solution of the Riccati equation if 
\begin{align}
    \frac{\text{d}}{\text{d} \tau}\begin{pmatrix}
        \mathds{X}_{jk} \\
        \mathds{Y}_{jk}
    \end{pmatrix} =  \mathds{A}_{jk}
    \begin{pmatrix}
        \mathds{X}_{jk} \\
        \mathds{Y}_{jk}
    \end{pmatrix} \nonumber
\end{align}
where 
\begin{equation}
    \mathds{A}_{jk}  = \frac{1}{2} \begin{pmatrix}
           ( H_j + H_k) \Omega  &  i ( H_j - H_k) \\
          i \Omega ( H_j - H_k)  \Omega^{\rm T }   & \Omega  ( H_j + H_k   )
          \end{pmatrix}
\end{equation}
under the initial conditions $\mathds{Y}_{jk}(0) = \sigma_0$ and $\mathds{X}_{jk}(0) = \mathds{1}$, to enforce that  $\sigma (0) = \sigma_0$ and that the solution is the physical one (or similarly $\mathds{Y}_{jk} = \sigma^0_{jk}$ for an initial GBCS). Being a linear system, the solution can be expressed via the transformation $\mathds{S}_{jk} (\tau) = \exp \left( \mathds{A}_{jk}  \tau \right) $ and is given by $(\mathds{X}_{jk}, \mathds{Y}_{jk})^{\rm T} (\tau) = \mathds{S}_{jk} (\tau) ( \mathds{1}, \sigma_0)^{\rm T}$. Hence, the analytical solution exists and can be found if it is possible to compute $\mathds{S}_{jk}$ and $\mathds{X}_{jk}$ is invertible.


In the specific case of a single qubit (with only $\sigma^{\text{off}} = \mathds{Y} \mathds{X}^{-1}$ having the Riccati terms and Hamiltonian of Eq.~\ref{eq:Hamiltonian_max_qubit}), the linear system reads 
\begin{equation}
    \frac{\text{d}}{\text{d} \tau}\begin{pmatrix}
        \mathds{X} \\
        \mathds{Y}
    \end{pmatrix} =
    \begin{pmatrix}
           H_m \Omega &     i  H_q \\
          i \Omega H_q \Omega^{\rm T}  & \Omega H_m 
    \end{pmatrix}\begin{pmatrix}
        \mathds{X} \\
        \mathds{Y}
    \end{pmatrix}  \;.
\end{equation}
This can in turn be used to solve the other phase space parameters coupled to the time evolution of the covariance matrix (see Appendix~\ref{app:doubling}). This procedure can be used to compute specific analytical solutions (as we shall see in the examples of Sec.~\ref{sec:operations} and Sec.~\ref{sec:physical_examples}) and semi-analytical in terms of a restrict number of linear systems. 

\section{Operations \& Deterministic Maps \label{sec:operations}}

In practice, it is often useful to describe operations that may not be dynamical and are instead applied on a state at a given time (while they may arise from fundamental dynamics). Alongside the class of local operations on the DV (such as a rotation) and Gaussian operations on the CV (such as a beam-splitter), GBCS can also be used to describe the Gaussian operations conditional on the qubit states: this is achieved as maps of the phase space parameters characterizing a GBCS. Two key operations of the latter class are conditional displacements and rotations, and they will be used as pedagogical examples in Sec.~\ref{sec:physical_examples} and whose dynamical and experimental implementations are discussed in Sec.~\ref{sec:physical_examples}. The interplay between these two conditional operations, allowing one to prepare GBCS, along with Gaussian CV processes and DV operations is key to achieving sophisticated quantum states and tasks. In Appendix~\ref{app:operation}, we provide the reader with proofs concerning the following sections.

\subsection{Gaussian Operations \label{subsec:gaussian_operations}}

A vast toolbox of experimentally useful operations are those that preserve the Gaussian structure of states, which have been extensively studied~\cite{serafini_quantum_2017}, and are immediately applicable on GBCS, given their phase space representation. As introduced in Sec.~\ref{sec:intro_phase}, a general Gaussian completely positive and trace preserving map $\Phi_g$ is uniquely determined by two matrices $X$ and $Y$, such that $Y+ i \Omega \geq i X \Omega X^{\rm T}$. As in the case of Gaussian states, by applying the action of the dual Gaussian map on the displacement operators, used as the CV basis, one finds that a GBCS evolves to the new parameters $\sigma_{jk} \to X  \sigma_{jk} X^{ {\rm T}} + Y$ and $
    \bar{r}_{jk}  \to X \bar{r}_{jk}$, while the QRDM is unchanged. 
Examples of these maps are: (a) classical mixing with $X= \mathds{1}$ and $Y >0 $, (b) attenuators with $X = \cos \theta \mathds{1}$ and $Y = (\sin \theta)^2 (2 N_p + 1) \mathds{1}$, and (c) amplifiers with $X = \cosh r \mathds{1}$ and $Y = (\sinh r)^2 (2 N_p + 1) \mathds{1}$, where $N_p\ge0$ parametrizes the noise induced by the operations. The specific instance of Gaussian unitary operations is captured by $X = S \in Sp_{2, \mathds{R}}$ and $Y = 0$, including, for instance, CV rotations, beam splitters, local squeezing, and two mode squeezing~\cite{serafini_quantum_2017} (see Appendix~\ref{app:operation} for their explicit forms). 

\subsection{Qubit Operations}

Often, in order to perform protocols, operations are independently implemented on the DV, for instance, rotations of qubits. A Completely Positive DV operation is defined by a set of Kraus operators $K^i = K^i_{jk} \mathds{1} \otimes \ket{j}\bra{k}$ , labelled by $ i \in [1, r]$ with $r \leq N^2$ (where we also include the trivial identical action on the CV subsystem). The specific instance of a quantum channel, which is also trace preserving, requires $\sum_{i = 1}^{r}  K^{(i) \;\dag}_{jk} K^{(i)}_{kl} = \mathds{1}$. The update of GBCS can be trivially computed by its definition and the action of the channel $\Phi_{q} [\hat{\varrho}] = \sum_{i = 1}^{r}  K^{i \; \dag} \hat{\varrho} K^i$, such that the characteristic functions are updated according to 
\begin{equation}
    \chi_{jk} (\tilde{\xi})\to \sum_{i = 1}^{r}  K^{(i) \;\dag}_{jm}  \chi_{m n}(\tilde{\xi}) K^{(i)}_{n k} 
\end{equation}
where we used the Einstein notation on rotated indexes. As apparent, this mixes the original branched characteristic functions.

For instance, the rotation around a the $x$-axis of a single qubits has only one Kraus operator 
\begin{equation}
    R_x(\theta) = {\rm e}^{-i\theta\sigma_x/2} = \begin{pmatrix} \cos(\theta/2) & -i\sin(\theta/2) \\ -i\sin(\theta/2) & \cos(\theta/2) \end{pmatrix} ,
\end{equation}
updating  the characteristic functions to 
   $ \chi_{jk} (\tilde{\xi})\to R^{jm}_x(\theta) \chi_{m n}(\tilde{\xi}) R^{nk}_x(\theta)  $. Hence, each branched characteristic function becomes a sum of Gaussians. By virtue of the fact that the PDEs are linear, the dynamics of the resulting state can still be analysed in terms of the ODEs and their solutions, feeding in $N^2$ initial conditions for each~ODE.

\subsection{Conditional Operations}

The set of quadratic operations conditional on the state of a DV subsystem -- including, for instance, conditional rotations, beam splitters, and displacements -- require utmost care, especially due to the required handling of unnormalized operators. A general conditional unitary quadratic operation may be written as
\begin{align}
    \hspace{-0.2cm}\Phi^{\mathcal{C}}[\hat{\varrho}] &= \hspace*{-0.2cm} \sum_{j, k  = 0}^{N-1} \int_{\mathds{R}^{2n}} \frac{\text{d} \xi}{(2 \pi)^n} \chi_{jk} (\xi) \hat{D}_{d_j} \hat{S}_{\theta_j} \hat{D}_{\xi} \hat{S}_{\theta_k}^\dag \hat{D}_{d_k}^\dag  \otimes \ket{j} \bra{k} , \nonumber
\end{align}
where $\hat{S}_{\theta_j} = \exp(i  \hat{r}^{\rm T} C_j  \hat{r}/2)$ and $\hat{D}_{d_j} = \exp(i  d_j^{\rm T} \Omega \hat{r}^{\rm T} )$, are two sets of $N$ exponential operator quadratic in the canonical operators, parametrized by some matrix $C_j$ and $d_j$. For linear operations, i.e., conditional displacement with $C_j = 0$, one can proceed by computing the update rule of the characteristic function $\chi_{jk}$ from its trace definition and compositions of displacement operators; allowing us to find the general updated rule (Sec.~\ref{subsubsec:constional_displacement}). However, for quadratic operations, there is not a general composition rule for $\hat{S}_{\theta_j}$ operators, and (also in the specific instances where compositions are possible) the trace is not computable as these operators are not trace norm. Thankfully, by knowing the generators of the transformations (i.e. the matrix $C_j$), one can use the ODEs of Table~\ref{tab:unitary}  (under $H_j \to C_j$, $r_j \to 0$ and $\tau = 1$) to obtain the update rule in differential form and compute specific instances for a given state.  

\begin{figure*}[!t]
    \centering
    \includegraphics[width=1\linewidth]{displacement_wigner.pdf}
    \caption{Representation of (a) a conditional displacement with vectors $d_\pm = \left( \bar{x} \pm \delta \bar{x} , \bar{p} \right)$ of the ground state creating a superposition in position and (b) a conditional rotation of angle $\theta_\Delta$ creating the same GBCS from a coherent state. (c) The Classical and Quantum component Wigner function given of the GBCS with arbitrary (asymmetric) displacement. \textcolor{red}{change wigner and angle on roff rot}} %of an initial squeezed thermal state undergoing diffusive Stern-Gerlach dynamics
    \label{fig:conditional_dis}
\end{figure*}

\subsubsection{Conditional Rotations}

The previous point is better illustrated by a specific example: the conditional rotation of a single CV mode by angles $\theta_{\pm}$, written as $\hat{S}_{\theta_\pm}  = \exp ( i \theta_\pm/2 (\hat{x}^2 + \hat{p}^2))$. This is the case of $H_{\pm} = \theta_{\pm} \mathds{1} $. On the diagonal elements, the two Gaussian unitary transformations give the update rules: $\sigma^{\text{on}}_\pm = \to R_{\theta_\pm}\sigma^{\text{on}}_\pm  R_{\theta_\pm}^{\rm T}$ and  $\bar{r}^{\text{on}}_\pm  \to  R_{\theta_\pm} \bar{r}^{\text{on}}_\pm $ with $R_{\theta_\pm} = \exp(\theta_\pm \Omega)$, both phase space rotations (Sec.~\ref{subsec:gaussian_operations}). The off diagonal ODEs read
\begin{align}
\label{eq:ODE_conditional_rot}
    \dot{\sigma}^{\text{off}} &= \bar{\theta} \left[\Omega \sigma - \sigma \Omega \right] + i \theta_{\Delta} \left( \mathds{1} - \sigma^{\text{off}} \sigma^{\text{off}} \right) ,\nonumber \\ 
    \dot{\bar{r}}^{\text{off}} &= \left[ \bar{\theta} \Omega -  i \theta_{\Delta} \sigma \right]\bar{r}^{\text{off}} , \\ 
    \dot{r}^{\text{(0) off}} & = - i \theta_{\Delta} \left[ {r}^{\text{off} \; \rm T}{r}^{\text{off}} +  \text{Tr}[\sigma^{\text{off}} ]/2\right] , \nonumber
\end{align}
where $\dot{r}^{\text{(0) off}}$ is the off diagonal exponent of the QRDM and we defined $\theta_{\Delta} = (\theta_+ - \theta_-)/2$ and $\bar{ \theta} = (\theta_+ + \theta_-)/2$. For an initial coherent state, with $\sigma_0 = \mathds{1}$ and $r_0 = (x_0 , p_0)^{\rm T}$, the evolution of the covariance matrix is trivial, i.e. $\sigma^{\text{off}} (\tau) = \mathds{1} \forall \tau$, and one recovers the known results  $r^{\text{off}} (\tau) = e^{-i \theta_\Delta } r_0$ and $\rho_q^{\text{off}} = \exp \left(  - |r_0|^2 \left( 1 - e^{-i 2 \theta_\Delta} \right)/2 -  i \theta_\Delta \right)$. Looking for a non-trivial evolution of the covariance matrix, one can compute the conditional rotation of a squeezed state -- with $\sigma_0 = \text{diag} (1/s, s)$ -- by solving the Riccati equation via the doubling method (Sec.~\ref{subsec:riccati}), finding the update rule
\begin{align}
\label{eq:cov_rotation}
     \sigma_0 &\to    R_{\bar{\theta}} \begin{pmatrix}
 {\rm e}^{i \psi} |\zeta | & 0 \\
 0 &  {\rm e}^{-i \psi}/|\zeta|
    \end{pmatrix} R_{\bar{\theta}}^{\rm T}  
\end{align}
where $ R_{\bar{\theta}}$ is the rotation of the mean angle and, by using the polar decomposition of the complex covariance matrix, one finds the squeezing strength $\zeta^2 = 
 \frac{\cos ^2(\theta_{\Delta} )+s^2 \sin ^2(\theta_{\Delta} )}{\sin ^2(\theta_{\Delta} )+s^2 \cos ^2(\theta_{\Delta} )} $ and the complex phase $\psi = \arctan \left( \frac{1}{2} (s -\frac{1}{s}) \sin(2 \theta_\Delta) \right) $.  In Appendix~\ref{app:operation}, we derive the update for the first moment and elements of the QRDM, which can be visualized as in  Fig.~\ref{fig:conditional_rot} over complex evolution.

\subsubsection{Conditional Displacements\label{subsubsec:constional_displacement}}

To create GBCS, displacements conditional on the qubit states are often performed (for instance, in trapped ion via the photon absorption recoil~\textcolor{red}{\cite{ }}). A general conditional displacement is given by $D_{d_1,...,d_N} = \sum_{i}^{N}  D_{d_j}  \ket{j} \bra{j}$, where the $d_j$'s are $N$ vectors parametrizing the displacements. 
The conditional structure of this operation allows one to write its action on a GBCS as 
\begin{align}
\label{eq:displacement_channel}
    \Phi^{D}[\hat{\varrho}] &= \frac{1}{(2 \pi)^n} \sum_{j, k  = 0}^{N-1}  \int_{\mathds{R}^{2n}} \text{d} \xi \; \chi_{jk} (\xi)  \hat{D}_{d_j} \hat{D}_{\xi} \hat{D}_{d_k}^\dag  \ket{j} \bra{k} .
\end{align}
By applying the composition rule of displacements ($ \hat{D}_{r_1} \hat{D}_{r_2} =  {\rm e}^{\frac{i}{2} r_2^{\rm T} \Omega r_1} \hat{D}_{r_1+r_2}$) and redefining the variable of integration, one can recast Eq.~(\ref{eq:displacement_channel}) into Eq.~(\ref{eq:chara_jk_def}), thus finding the map of the GBCS parameters under general conditional displacements:  
\begin{align}
\label{eq:cond_dis_update}
    \bar{r}_{jk} &\to \bar{r}_{jk} - \frac{1}{2} (d_j +d_k) - \frac{i}{2} \sigma_{jk} \Omega  (d_j - d_k), \nonumber\\
    \mathcal{C}_{jk} &\to \mathcal{C}_{jk} + \frac{1}{4} \left( d_j - d_k\right)^{\rm T} \Omega^{\rm T} \sigma_{jk}  \Omega \left( d_j - d_k\right),\nonumber \\
    &\hspace{1.5cm} +  (d_j - d_k)^{\rm T} \Omega \; \text{Im} [\bar{r}_{jk}]  \nonumber \\
    \phi_{jk} &\to \phi_{jk} - \frac{1}{2}  d_j^{\rm T}  \Omega d_k  +  (d_j - d_k)^{\rm T} \Omega \text{Re} [ \bar{r}_{jk}] ,
\end{align}
where we separated the real and complex parts of the first moments $\bar{r}_{jk} = \text{Re} [ \bar{r}_{jk}]  + i \text{Im} [\bar{r}_{jk}]$. As expected, conditional displacements do not affect the covariance matrix ($\sigma_{jk} \to \sigma_{jk}$). 

For the case of a single qubit,  $\hat{D}_{d_1, d_{-1}} = \hat{D}_{d_+} \ket{1}\bra{1} + 
    \hat{D}_{d_{-}} \ket{-1}\bra{-1} $ for the two displacement vectors $d_\pm$. One finds that $\bar{r}^{\text{on}}_{\pm} \to \bar{r}^{\text{on}}_{\pm} - d_\pm$ and 
\begin{align}
\label{eq:cond_dis_update_qubit}
    \bar{r}^{\text{off}} &\to \bar{r}^{\text{off}} - \bar{d} - i  \sigma^{\text{off}} \Omega  d_\Delta, \nonumber\\
    \mathcal{C} &\to \mathcal{C} +  d_{\Delta}^{ \rm T} \Omega^{\rm T} \sigma^{\text{off}}  \Omega d_\Delta  \nonumber + 2 d_\Delta^{\rm T} \Omega  \text{Im} [\bar{r}^{\text{off}}] \nonumber \\
    \phi &\to \phi - \frac{1}{2}  d_+^{\rm T}  \Omega d_{-} + 2  d_\Delta^{\rm T} \Omega\text{Re} [\bar{r}^{\text{off}}] ,
\end{align}
where $\bar{d} = (d_+ + d_-)/2$ and  $d_\Delta = (d_+ - d_-)/2$ are the average and difference in the displacements. The visualization of this the complex update rules is in Fig.~\ref{fig:conditional_dis}.

\section{Geometry of Wigner Functions \& Interpretation for Complex Moments \label{sec:wigner}}
\begin{figure*}[!t]
    \centering
    \includegraphics[width=1\linewidth]{rotation_wigner.pdf}
    \caption{Representation of (a) a conditional displacement of the ground state of two vactors $d_\pm$ and (b) a conditional rotation of angle $\theta_\Delta$ creating the same GBCS. (c) The Classical and Quantum component Wigner function given of the GBCS. \textcolor{red}{add Gaussians on the side, change color of centered gaussian, add map in fig C}} %of an initial squeezed thermal state undergoing diffusive Stern-Gerlach dynamics
    \label{fig:conditional_rot}
\end{figure*}
The interplay between the Wigner function of GBCS and the phase-space parametrization allows us to provide the reader with a geometrical description of the Wigner negativity common in cat-states and, hence,  with an interpretation for the real and imaginary components of the first and second moments. By computing the Fourier transform of Eq.~(\ref{eq:chara_jk}), the branched Wigner functions of a general GBCS are given by
\begin{equation}
\label{eq:wigner_jk}
    \mathcal{W}_{jk} (\xi) = \frac{2^n \varrho_{jk}^q}{\pi^n \sqrt{\text{det} \sigma_{jk} }}\exp \left( -\left(\xi - \bar{r}_{jk} \right)^{\rm T}  \sigma_{jk} ^{-1} \left(\xi - \bar{r}_{jk} \right) \right).
\end{equation} 
On the diagonal terms, they are real and positive. On the off diagonal terms, $\mathcal{W}_{jk}$ are complex which, when summed over $j$ and $k$, interfere with each other; giving raise to the fringes typical of GBCS.

To illustrate this, let us consider the case of $n = 1$ and $N=2$ (i.e. the state of Eq.~\ref{eq:density_one_qubit}). For the phase-space picturization of GBCS, two quantities are relevant: the  Wigner functions of the classical components, $\mathcal{W}_c := \mathcal{W}_+  + \mathcal{W}_{-}$, and the `quantum interference', given by $\mathcal{W}_q :=   \mathcal{W}^{\rm off} + \mathcal{W}^{\rm off\; *} = 2 \text{Re} [\mathcal{W}^{\rm off}]$. The former, known since~\textcolor{red}{~\cite{ }}, is the statistical mixture of the Wigner functions of two Gaussian states with moments $\sigma_{\pm}$ and $ \bar{r}_{\pm}$, weighted according to the on-diagonal probabilities ($p_{\pm}$). This contribution to the Wigner function can be used to build physical intuition about the processes that the system undergoes, representing the two CV Gaussian processes: the first moments are the \textit{semi-classical} interferometric paths taken with uncertainties captured by the corresponding covariance matrix. However, such representation, which is restricted to the real plane, does not give a full description of the quantum state (for instance, omitting the relative phase between the branches). 

The quantum interference component -- which encodes the truly quantum features of GBCS -- find its geometrical interpretation with the off diagonal first and second moments evolving in the complex plane. The form of the quantum component can always be expressed as
\begin{align}
\mathcal{W}_q (\xi) &= \mathcal{N} \mathcal{G} (\xi) \cos \left[ \Phi_W(\xi) + \Phi_0 \right], 
\end{align}
where $\mathcal{G}(\xi)$ is a smearing Gaussian distribution, $\Phi_W(\xi) $ is the Wigner phases of interference (function of $\xi$),  $\mathcal{N}$ is a normalization constant, and $\Phi_0$ is an additional constant phase. Let us define the real and complex components of the first moments and the inverse of the covariance matrix: namely, $\bar{r}^{\text{off}} = \bar{r}_R + i \bar{r}_{I}$ and $\sigma^{\text{off} -1} = \sigma_R^{-1} + i \sigma_{I}^{-1}$\footnote{Where we note that, for instance, $\sigma_R^{-1}$ is \textit{not} the inverse of the real part of the covariance matrix, but is the real part of the inverse of the covariance matrix. The are equivalent only is $\sigma^{\text{off}}$ is diagonal.}. Then, one finds the general 
\begin{align}
\mathcal{G} (\xi) &= \exp \left( {-(\xi -\bar r_R)^{\mathsf T} \sigma_R^{-1} (\xi -\bar r_R)- 2 \tilde \xi^{\mathsf T} \sigma_{I}^{-1}  \bar r_I} \right) \nonumber \\
\Phi_W (\xi) &= (\xi-\bar r_R)^{\mathsf T} \sigma_{I}^{-1} (\xi-\bar r_R)-2(\xi -\bar r_R)^{\mathsf T} \sigma_R^{-1}  \bar r_I \nonumber \\
\mathcal{N} &= \frac{4|q|}{\pi\sqrt{|\det(\sigma)|}}\,\exp \left( - \mathcal{C} + \bar r_I^{\mathsf T}\sigma_R^{-1}\bar r_I + \bar r_R^{\mathsf T} \sigma_{I}^{-1}  \bar r_I \right) \nonumber \\
\Phi_0 &=  \bar{r}_I^{\mathsf T} \sigma_{I}^{-1} \bar{r}_I - \frac{1}{2} \arg ( \det(\sigma)) + \phi
\end{align}

First, let us consider the case when the superposition is only present in the first moments and $ \sigma^{\text{off}} \in \mathds{R}$. Then, one can understand that the real components of the off-diagonal first moments ($\bar r^{R}$) are the centre of the Gaussian envelope that smear the fringes. The complex components $\bar r^{I}$ of the first moments are what causes the interference fringes, which in this case, are linear in $\xi$. The fringes are fully characterized by the Wigner wave-vector $k_W = 2 \sigma^{-1}\bar{r}_{I}$ with the zero at $\bar{r}_R$. Furthermore, the imaginary part of the first moments induces an increment in the normalization (which often, in ideal cases, cancels the DV decay $\mathcal{C}$). Conditional displacements and rotations presented in the Section~\ref{sec:operations} can be used as examples:  (a) a conditional displacement of the vacuum (Eq.~\ref{eq:cond_dis_update} with $r_0 = (0,0)$) has fringes $k_W =  2 \Omega d_{\Delta}$ and the envelope Gaussian is centred at $\bar{d}$, with $\mathcal{N} = 2/\pi$; or (b) by a conditional rotation from a coherent state with $r_0 = (x_0, p_0)$ has $k_W =  - 2\sin\!\left(\frac{\Delta\theta}{2}\right) r_0$ smeared by a normalize Gaussian ($\mathcal{N} = 2/\pi$) of minimal uncertainty centred at $\cos\!\left(\frac{\Delta\theta}{2}\right) r_0$. This allows the interpretation and picutrization of Fig.~\ref{fig:conditional_dis}. \textcolor{red}{change in the figure the vectors for $k_w$ and }


Secondly, when $\sigma^{\text{off}} \in \mathds{C}$, the imaginary part of the inverse of the covariance matrix ($\sigma_{I}^{-1}$), induces non-linear fringes that are quadratic in the phase space variables, via the term $(\xi-\bar r_R)^{\mathsf T} \sigma_{I}^{-1} (\xi-\bar r_R)$.
It further includes a shift of the Gaussian envelope (of the vector $\sigma_{I}^{-1}  \bar r_I$) and a change in the normalization ($\exp( \bar r_R^{\mathsf T} \sigma_{I}^{-1}  \bar r_I )$). 
For example, in the case of the conditional rotation of a squeezed state (Sec.~\ref{sec:operations}, with $\bar{\theta} = 0$), the resulting covariance matrix is the one of a squeezed state tilted in the complex plane of an angle $\psi$ (Eq.~\ref{eq:cov_rotation}) as shown in Fig.~\ref{fig:conditional_rot}. This implies that $\sigma^{-1}_R = \cos \psi \text{diag} \left(1/\zeta , \zeta  \right) $ and $\sigma^{-1}_I =  \sin \psi \text{diag} \left( - 1/\zeta , \zeta \right) $. This implies that the quadratic term is hyperbolic, as $\Phi_W $ contains a term $\sin \psi  \left( - \tilde{x}^2 /\zeta + \tilde{p}^2 \zeta \right) $~\cite{nicacio_phase_2010}.


\section{Measurements \& Expectation Values \label{sec:measure}}

Having solved a large class of deterministic dynamics and operations, one may be interested in the behaviour of GBCS under quantum measurements. As in the case of dynamics and operations, the phase-space description of the system allows for easy computations and for the compact expression of expectation values, measurement probabilities, and post measurement quantum states in terms of the phase-space parameters $\sigma_{jk}$, $\bar{r}_{jk}$, and $\varrho_{jk}^q$. 

Let us start by giving an overview of computing expectation values of GBCSs, by considering the general observables of the joint system $  \hat{\mathcal{O}}^m\otimes\hat{\mathcal{O}}^q$. The CV observables can be canonical operators $\hat{r}^k$ (where $k$ labels the $k$-th element of the $\hat{r}$ vector), products of these (e.g., the product $\hat{r}^{k_1} \hat{r}^{k_2}$ or higher order terms), or observables in the ladder operator basis (such as the phonon number $a^\dag a$, which can be always computed via the transformation of Eq.~\ref{eq:change_basis}). Any DV observable can be written in the coupling basis as $\hat{\mathcal{O}}_q = \sum_{j, k  = 0}^{N-1} O^{q}_{jk} \ket{j} \bra{k}$. 

In order to evaluate the joint system expectation value one can first perform the trace on the modes degrees of freedom. Generically, the product of $n$ canonical observables $\mathcal{O}^m = \hat{r}^{k_1} \hat{r}^{k_2}... \hat{r}^{k_n} $ can be computed from the repetitive use of the generating function
\begin{equation}
\label{eq:expectation_characterstic}
    \braket{\hat{\mathcal{O}}^m_{jk}} = (-i )^n \left. \frac{\partial^{n} \chi_{jk}(\tilde{\xi})}{\partial \tilde{\xi}^{k_1} \partial \tilde{\xi}^{k_2} ... \partial \tilde{\xi}^{k_n}} \right|_{r=0} \;,
\end{equation}
via the partial derivative of the corresponding phase space variable. Eq.~(\ref{eq:expectation_characterstic}) can always be evaluated in terms of $\sigma_{jk}$, $\bar{r}_{jk}$, and $\bar{r}_{jk}^{(0)}$. Specifically, for a general CV observables, one finds $\braket{\hat{r}}_{jk} := \text{Tr} (\hat{\varrho}_{jk} \hat{r}) = \varrho_{jk}^q \bar{r}_{jk}$ and similarly $\braket{\hat{r} \hat{r}^{\rm T}}_{jk} := \text{Tr} (\hat{\varrho}_{jk}\hat{r} \hat{r}^{\rm T}) = \varrho_{jk}^q  \left( \sigma_{jk} - \bar{r}_{jk} \bar{r}_{jk}^{\rm T} \right)$, representing indeed the first and second moments of each element of the density matrix, respectively. Finally, one can compute the joint observable by tracing over the DV Hilbert space, such that $\braket{\hat{\mathcal{O}}^m\otimes\hat{\mathcal{O}}^q} = \sum_{jk} \braket{\hat{\mathcal{O}}^m_{jk}} \mathcal{O}_{jk}^q$. Two specific cases can be derived from this general treatment. The expectation value of any observable of the qubit subsystem can be computed via the QRDM, as Eq.~(\ref{eq:expectation_characterstic}) simplifies to $\chi_{jk} (0) = \varrho_{jk}^q$, implying $\braket{\mathds{1} \otimes \hat{\mathcal{O}}_q } = \sum_{jk} O_{jk} \varrho^{q}_{jk}  $. Any observables of only the CV subsystem can be computed as the statistical mixture (weighted according to the QRDM diagonal terms, i.e. probabilities) of the expectation values computed with the Gaussian state of each branch: for example, the vectors of expansion values reads $\braket{\hat{r} \otimes \mathds{1}} = \sum_{i} \varrho_{ii}^q r_{ii} $, while the matrix of correlations is $\braket{\hat{r}\hat{r}^{\rm T} \otimes \mathds{1}} = \frac{1}{2}\sum_{i} \varrho_{ii}^q \left( \sigma_{ii} - r_{ii} r_{ii}^{\rm T} \right)$. 



Let us now move on to describing the granular outcomes of measurements and their associated conditional states, starting from single-subsystem measurements and concluding with joint measurements.

An ideal projective measurement of an observable $\mathcal{O}^q$ has $2^N$ outcomes (which may be degenerate) given by the eigenvalues ${\rm e}^i$. From the associated eigenvectors $\ket{{\rm e}^i}_q$ it is possible to derive the positive operator valued measure (POVM), i.e., in ths case, the set of projectors $\{M^i =\ket{{\rm e}^i}_q \bra{{\rm e}^i}_q= M^i_{jk} \ket{j}\bra{k} \}$, where $\ket{j}$ and $\ket{k}$ refer to the coupling basis, and such that $\sum_i M^i = \mathds{1}$. The probabilities of the $i$-th  outcome can be computed as the expectation value using the QRDM as previously described for observables of the qubit's subsystem, i.e., as $P(i) = \sum_{jk} M^i_{jk} \varrho_{jk}^q$. The post-measurement state after the outcome $i$ has the characteristic function 
\begin{equation}
    \chi^{\text{post}}_i (\xi) = \frac{1}{P(i)} \sum_{j, k  = 0}^{N-1}  M^{i}_{jk} \chi_{jk} (\xi) \;,
\end{equation}
 which may well show fringes of GBCS and negativity in the Wigner function (see Sec.~\ref{sec:physical_examples}). 

The class of Gaussian measurements of the CV subsystem is described by the ``general-dyne" measurements (with the POVM of Eq.~\ref{eq:generadyne}). These measurements are described by sets of projectors $\Pi_{r_o}$ labelled by the continuous labels of possible outcomes $r_o$, which represent the first moments of the post-measurement Gaussian state $\hat{\varrho}_o$ of covariance $\sigma_o$, i.e., $\Pi_{r_o} = \hat{D}_{-r_o} \hat{\varrho}_g \hat{D}_{r_o}/(2 \pi)^n$, where $\varrho_{g}$ is the Gaussian state with null first moments and covariance matrix $\sigma_{o}$. As in the case of expectation values, it is convenient to first trace over the CV subsystem, which gives the unnormalized post-measurement QRDM associated with the outcome $r_o$, i.e.,  $ \varrho^{q,\; \text{post}} (r_o) = \text{Tr}_m [\Pi_{r_o} \hat{\varrho}_{jk} ]$. Similarly to Eq.~(\ref{eq:prob_generadyne_gaussian}), its elements can be easily computed as the overlap between a GBCS and the measurement Gaussian state,\footnote{This follows from the orthogonality of Weyl operators under trace, $\text{Tr}_m [\hat{D}_{r} \hat{D}_{-s}] = (2 \pi)^{n} \delta^n (r - s)$ and from the Fourier--Weyl representation of the involved states.} such that
\begin{equation}
\label{eq:qubit_PMST}
    \varrho_{jk}^{q,\; \text{post}} (r_o) = \varrho^q_{jk}  \frac{{\rm e}^{ - (\bar{r}_{jk} - r_o)^{\rm T} ( \sigma_{jk} + \sigma_o )^{-1}  (\bar{r}_{jk} - r_o) }}{ \pi^n \sqrt{ \text{det} [ \sigma_{jk} + \sigma_o ] }} \;, \nonumber
\end{equation}
while the CV state has collapsed into the Gaussian state that has been measured. The probability distribution of outcomes $r_o$ can be found by tracing out the unnormalized post-measurement QRDM 
\begin{equation}
\label{eq:prob_generadyne}
    P(r_o) = \sum_{i} \varrho^q_{ii}  \frac{{\rm e}^{ - (r_{ii} - r_o)^{\rm T} ( \sigma_{ii} + \sigma_o )^{-1}  (r_{ii} - r_o) }}{\pi^n  \sqrt{ \text{det} [ \sigma_{ii} + \sigma_o ] }} \;,
\end{equation}
which is the statistical mixture of the general-dyne measurement of $N$ Gaussian states, weighted according to the QRDM diagonal elements (see Fig.~\ref{fig:example_2}).

This treatment encompasses ideal and noisy homodyne and heterodyne detection schemes, which we will describe for a single mode only but are trivially generalisable to $n$ modes. Specifically, homodyne detections measure the observable $\hat{x}_{\phi} = \cos (\phi) \hat{x} + \sin(\phi) \hat{p}$, by mixing the system state with a probing coherent state of high intensity at the same frequency ($\ket{\alpha} = |\alpha| {\rm e}^{i \phi}$ with $\alpha \gg 1$). For finite detection efficiency $0 \leq \eta \leq 1$, homodyne detection is described by the general-dyne POVM with covariance matrix 
\begin{equation}
\label{eq:homodyne}
    \sigma_o = \lim_{z \to 0} R^{\rm T} (\phi) \begin{pmatrix}
        z^{2} & 0  \\
        0 & 1/z^2 
    \end{pmatrix} R (\phi) + (\tan \theta)^2 \mathds{1} \;,
\end{equation}
where we define $\theta = \cos^{-1}(\sqrt{\eta})$. The heterodyne detection schemes is similar to the homodyne one but the reference signal of the probe has a different frequency from the measured mode. In this case, $\sigma_o = (1 + 2( \tan(\theta))^2) \mathds{1}$. From Eq.~(\ref{eq:prob_generadyne}), it is possible to calculate the probability distribution of these two specific measurements, which are  apt to detect statistical mixtures of Gaussian states (see Sec.~\ref{sec:physical_examples}). The ideal case can be recovered by setting $\eta = 1$ (i.e., $\theta =0$).


Finally, the measurement of the joint system described by a POVM $\Pi_{r_m} \otimes M^i$ has $2^N$ probabilities distributions of $r_m$, labeled by the discrete outcome $i$, which can be expressed as
\begin{equation}
    P^{i}(r_o) = \sum_{j, k  = 0}^{N-1} M^{i}_{jk} \frac{\varrho^q_{jk} {\rm e}^{ - (\bar{r}_{jk} - r_o)^{\rm T} ( \sigma_{jk} + \sigma_o )^{-1}  (\bar{r}_{jk} - r_o) }}{\pi^n  \sqrt{ \text{det} [ \sigma_{jk} + \sigma_o ] }}\; . \nonumber
\end{equation}
Such probabilities distributions are the fringe pattern often related to Schr\"odinger cat states, which in the DV-CV systems arises under joint measurement (see Fig.~\ref{fig:example_1}). 




\section{Perturbation Theory \label{sec:perturbation}}


While the operator-valued interactions are very common in practice, often the Hamiltonians and noises of hybrid systems have additional uncoupled non-diagonal DV terms, which one would like to incorporate into the presented formalism. In the following (and Appendix~\ref{app:perturbation}), we develop a perturbation theory that exploits the analytical solution of GBCS to account for Hamiltonian and Lindbladian terms involving DV operators that do not commute with $\hat{J}_z$ (of the form $\hat{V} = \mathds{1}_m \otimes \sum_{jk} V_{jk} \ket{j} \bra{k}$). The perturbative analysis is indeed required in these cases, as the algebra of operators does not close if one incudes such non-commuting DV operators. 

For the following, it may be useful to vectorize the DV Hilbert space (to dimension $N^2$), according to the known practices~\cite{fundamentals_gyamfi_2020}: for instance, the state of Eq.~(\ref{rhojk}) is $\dket{\varrho} = \sum_{j, k  = 0}^{N-1} \hat{\varrho}_{jk} \ket{j} \otimes \ket{k}$, with the specific case of a qubit $\dket{\varrho} = (\hat{\rho}_{11}, \hat{\rho}_{01}, \hat{\rho}_{10}, \hat{\rho}_{00})^{\rm T}$. 
The solutions of quantum dynamics presented in Sec.~\ref{sec:dynamics} relies on solving $N^2$ uncoupled homogeneous PDE for the characteristic functions: $\left( \partial_\tau - \mathcal{L} \right) \dket{\chi} = 0$, where $\dket{\chi}$ is the vectorial form of branched characteristic functions and $\mathcal{L}$ is a diagonal superoperator (a matrix $N^2 \times N^2$). Each element of $\mathcal{L}$ is a complex differential operator in the phase space, quadratic in the variables $\xi$ and its partial derivatives. 

If some additional non-diagonal DV operators are present in the quantum dynamics, they result in an additional inhomogeneous terms that couple the PDEs (as a time-dependent source term~\cite{Antonakakis_close_2013}): the phase-space evolution reads $\left( \partial_\tau - \mathcal{L} \right) \dket{\chi} = \mathds{V}\dket{\chi}$, 
where  $\mathds{V}$ is a Liouvillian DV superoperator independent from~$\xi$. As shown in Appendix~\ref{app:perturbation}, the solutions presented in Sec.~\ref{sec:dynamics} can be formulated in terms of the diagonal matrix of spatial Green's function  $\mathcal{G}(\tau,  \tilde{\xi}, \tilde{\xi}')$, for instance, for a qubit, $\mathcal{G}= \text{diag} (\mathcal{G}_{00}, \mathcal{G}_{01},\mathcal{G}_{10},\mathcal{G}_{11})$, such that one finds the the solution of the homogenous equations presented in the Sec.~\ref{sec:dynamics} as $\dket{\chi^{h} (\tau, \tilde{\xi})} = \int_{\mathds{R}^{2n}} \text{d} \tilde{\xi}' \mathcal{G}(\tau,  \tilde{\xi}, \tilde{\xi}') \dket{\chi_0(\tilde{\xi}')}$, with the vectorized initial conditions $\dket{\chi_0}$. Then, the solution of the coupled PDEs can be formulated as a recursive integral~\cite{Gripenberg1990}
\begin{equation}
\label{eq:volterra}
    \dket{\chi(\tau, \tilde{\xi})} = \dket{\chi^{h} (\tau, \tilde{\xi})}  + \int_0^\tau\hspace*{-2mm} \text{d} t \hspace*{-1mm}\int_{\mathds{R}^{2n}}\hspace*{-4mm} \text{d} \tilde{\xi}' \mathcal{G}(\tau - t ,  \tilde{\xi}, \tilde{\xi}') \mathds{V} \dket{\chi(t,\tilde{\xi}')},  \nonumber
\end{equation}
where the additional term is a particular solution of the inhomogenous PDE. 
Define the kernel $\mathds{k}(\tau, \tilde{\xi}, \tilde{\xi}') \equiv \mathcal{G}(\tau,\tilde{\xi},\tilde{\xi}')\,\mathds{V}\ $, the phase-space convolution as 
\begin{equation}
    (\mathds{k}*\mathds{k})(\tau, \tilde{\xi}, \tilde{\xi}') = \int_0^\tau \text{d} t \int_{\mathds{R}^{2n}} \text{d} \tilde{\xi}'' \mathds{k}(\tau- t, \tilde{\xi}, \tilde{\xi}'')\mathds{k}(t, \tilde{\xi}'', \tilde{\xi}') \nonumber
\end{equation}
and its $k$-th iteration as $\mathds{k}^{* k}.$
Under these definitions, the integrand reads $\dket{\chi(\tau,\tilde \xi)}=\dket{\chi^{h}(\tau,\tilde \xi)}+(\mathds{k}*\dket{\chi})(\tau,\tilde \xi)$, where $*$ also labels convolution between a superoperator and a vector. This types of integrals can be solved via an iterative series~\cite{Gripenberg1990}
\begin{equation}
    \label{eq:series}
    \dket{\chi(\tau,\tilde{\xi})} = \dket{\chi^{h} (\tau,\tilde{\xi})} + \sum_{k=1}^{\infty} \big[\mathds{k}^{*k}*\dket{\chi^{h}}\big](\tau,\tilde{\xi})\ .
\end{equation}
which is the phase-space representation of the Dyson series for superoperators and the nested convolutions respect time-ordering. In the power of $\mathds{V}$, the zeroth order is the homogenous solution, the first order is explicitly given by
\begin{equation}
    \label{eq:order1}
    \dket{\chi^{(1)}(\tau,\tilde{\xi})} = \int_0^\tau {\rm d}t \int_{\mathds{R}^{2n}} {\rm d}\tilde{\xi}'\ \mathcal{G}(\tau-t,\tilde{\xi},\tilde{\xi}')\,\mathds{V}\,\dket{\chi^{h}(t,\tilde{\xi}')} \;.  \nonumber
\end{equation}
%and the second is $\dket{\chi^{(2)}(\tau,\tilde{\xi})} = \big[\mathds{k}*\mathds{k}*\dket{\chi^{h}}\big](\tau,\tilde{\xi}) $  and the $N$-th is given in Appendix~\ref{app:perturbation}. As it will be explicit in the 

\subsection{DV Decay Noise: $T_1$}

Decay in the DV subsystem has to be considered to give a realistic analysis of quantum systems in which the decay time ($T_1$) is similar to the dephasing one  ($T_2$), for instance, in some trap ion and superconducting qubit architectures with large coherence time~\cite{Wang2021,bland20252dtransmonslifetimescoherence}. The additional Lindbladian term which, in the most general case of an $N$-level system, describes decay is  
\begin{equation}
\label{eq:Nlevel_lindblad}
    D_{T_1} [\rho] = \sum_{j=0}^{N-1} \hat{L}_j \rho\, \hat{L}_j^\dagger - \tfrac{1}{2} \{ \hat{L}_j^\dagger \hat{L}_j, \rho \}, 
\end{equation}
with jump operators $\hat{L}_{j} = \sqrt{\Gamma_{j}} \ket{j-1}\bra{j}$, for $j \in [0,N-1]$. To illustrate how the perturbation theory can be applied to this context, let us consider the single qubit state of Eq.~\ref{eq:density_one_qubit} under qubit decay, i.e.  $\hat{L}_{1} = \sqrt{\Gamma_{1}} \hat{\sigma}_-$. Observe that, in the case of decay, the truncation of the series occurs exactly as soon as all the excitations in the DV sector have decayed to the lowest excitation state.  In the case of the qubit,  such that the evolving state is described by the vector of characteristic functions
\begin{equation}
    \dket{\chi^{T_1} (\tau, \tilde{\xi})} = \begin{pmatrix}
        {\rm e}^{-\Gamma_{1}\tau  } \; \chi^{\text{on}}_+ (\tau, \tilde{\xi})    \\
{\rm e}^{-\Gamma_{1}\tau  /2 }\; \chi^{\text{off}}  (\tau, \tilde{\xi})  \\
{\rm e}^{-\Gamma_{1}\tau /2 } \;\chi^{* \text{off}} (\tau, \tilde{\xi})    \\
\chi^{\text{on}}_- (\tau, \tilde{\xi})   + \chi^{(1)}  (\tau, \tilde{\xi})
    \end{pmatrix}  ,
\end{equation} 
where the perturbation arises only in the DV ground state  (the only state where the excited state can decay) and it has the additional characteristic function
\begin{equation}
\label{eq:decay_qubit}
    \chi^{(1) } (\tau, \tilde{\xi})  =  p_+ \Gamma_1 \int_0^{\tau}  {\rm d} t {\rm e}^{-\frac{1}{4} \tilde{\xi}^{\rm T} \sigma^{\Gamma_1} ( \tau, t) \tilde{\xi}  +  i  \tilde{\xi}^{\rm T} \bar{r}^{\Gamma_1} ( \tau, t) - \Gamma_1 t } ,
\end{equation}
with the ``two times" covariance matrix and first moments given by
\begin{align}
\label{eq:two_times_parameters}
    \sigma^{\Gamma_1} ( \tau, t) &= S_-(\tau - t) S_+(t) \sigma_0 S^{\rm T}_+(t) S^{\rm T}_-(\tau - t) \\
    \bar{r}^{\Gamma_1} ( \tau, t) &= S_-(\tau - t) \left[ S_+(t) r_0 - \left(  S_+(t) - \mathds{1} \right) r_+ \right] \nonumber \\
    &\hspace{2cm} - \left(  S_-(\tau - t) - \mathds{1} \right) r_-   \nonumber
\end{align}
where, to illustrate the point, we considered the unitary case in the CV ($S_{\pm}$ are given in Eq.~\ref{eq:symplectic_pm}, and $r_{\pm} = r_m \pm r_q$), while the open dynamics is discussed in Appendix~\ref{app:perturbation}. The change in the qubit populations can be computed from the diagonal elements of the QRDM, which undergoes decays at rate $\Gamma_1$ for the excited states ($\rho^q_{11} = p_+ e^{-\Gamma_1 \tau}$), while the ground state increases
\begin{align}
    \rho^q_{00} = \chi^{h}_- (0,\tau) + \chi^{(1)}(0,\tau) &  = p_- + p_+ \Gamma_1 \int_0^{\tau} dt e^{-\Gamma_1 t} \nonumber\\
    &= p_- + p_+ (1 - e^{-\Gamma_1 \tau}) \;. \nonumber
\end{align}
as expected. The presented solution finds a clear interpretation as a quantum jump: in the interval of time $[t, t + dt]$, there is a probability $P(t) - P(t+dt) = p_+  e^{-\Gamma_1 (t+dt)} - p_+  e^{-\Gamma_1 t} \approx p_+  \Gamma_1 e^{-\Gamma_1 t} dt $ that a jump occurs. When the jump does not occur, the branched characteristic functions $\dket{ \chi}$ evolve according to their given dynamics, and are described by their homogenous solution (with the inclusion of the diagonal part of $T_1$ noise which gives the decays). When the jump occurs, $\chi^{\text{on}}_+ (\tilde{\xi},t)$ is the new initial condition for the dynamics of the $-$ branch (given by Eq.~\ref{eq:two_times_parameters}), and all other branches are annihilated. After integrating over all possible jump times, one finds the total state as, indeed, the one of Eq.~(\ref{eq:decay_qubit}).  An example of the Wigner function arising from $\chi^{(1) } (\tau, \tilde{\xi})$ is given in Sec.~\ref{sec:physical_examples}. In a similar way, one can describe the decay of an $N$ level system as $N$ possible jumps; as this noise only induces change in populations on the diagonal term and additional decay in the off diagonal one.


\begin{figure*}[!t]
    \centering
    \includegraphics[width=1\linewidth]{Example_Figure_Beggining.pdf}
    \caption{Schematic Representations of the two presented examples along with the Wigner functions of the related statistical mixture of Gaussian processes: (a) Resonator entangled via frequency shift $\chi$ to two qubits undergoing measurement-based entanglement ($\eta$, homodyne efficiency and $\kappa$, decay rate) and (b) a Stern-Gerlach matter-wave interferometer of a mass $m$ trapped at frequency $\omega$ with operator valued force $f_q$ in diffusive environment ($\Gamma_x$).} %of an initial squeezed thermal state undergoing diffusive Stern-Gerlach dynamics
    \label{fig:example_initial}
\end{figure*}


\section{Physical Examples \\ 
\& Comparative Analysis \label{sec:physical_examples}}

In this Section we present two examples that showcase our method's capability of producing predictions for realistic experiments. We shall cover (a) a qubit (or atom) dispersively interacting with a resonator undergoing dynamical squeezing induced by a parametric amplifier; (b) a levitated nanoparticle undergoing Stern-Gerlach non-Gaussian matter-wave interferometry. \textcolor{red}{This is schematically represented in Fig.~\ref{fig:example_initial}, along with the Wigner function of the statistical mixture of the Gaussian processes analysed in this section. }
The former example represents an instance of quadratic operator-valued DV-CV interaction, reproducing, in the steady state, experimentally detected quantities~\cite{eddins_high_2019}. The latter showcases the results of linear DV-CV interactions in open dynamics: providing novel analytical predictions for sensing and generation of Wigner negativity with levitated systems in the open non-Gaussian regime. Both these examples will allow us to illustrate the physical roles of the phase-space parameters and how they relate to dynamical properties of realistic quantum systems (via the element of Table~\ref{tab:open}).




%In the process, we will present novel results, to the best of our knowledge: such as Wigner functions fringes and probability measurements -- both for the specific examples, which may be known -- and the general form. 


\subsection{Transient Behaviour or
 Parametric Amplified Dispersive Qubit Readout\label{subsec:atom}}



The coherent absorption and emission of a cavity photon by and to a qubit is central to quantum computation. Examples of architectures featuring such a behaviour are found in superconducting Josephson junctions~\cite{strong_wallraff_2004,coherent_chiorescu_2004,coupling_majer_2007,blais_quantum_2007}, as well as dopant spins with magnetic~\cite{strong_kubo_2010,strong_chen_2018,realization_liu_2023} and electric~\cite{strong_samkharadze_2018,strong_samkharadze_2018,coherent_mi_2018,coherent_landig_2018} coupling.  Resonators, on the other hand, range from optical cavities to LC superconducting circuits and nanomechanical resonators~\cite{xiang_hybrid_2013}.
In these schemes, the dispersive coupling between the qubit and the cavity creates a distinguishable superposition of the CV state that can be readout in order to perform a non-demolition measurement of the qubit.  Squeezing the resonator is known to provide quantum enhancement ~\cite{Entanglement_Wahyu_2008,Sete2013,Szorkovszky_2014,didier_fast_2015,walter_rapid_2017,kam_fast_2024} and it has been experimentally demonstrated by embedding a qubit in a driven parametric amplifier that dynamically squeezes the state of the resonator, thus increasing the steady-state detection efficiency~\cite{eddins_high_2019}. Here, we discuss the dispersive dynamics and the transitional behaviour underlying qubit reading-out processes which exploits dynamical squeezing.



%On the other hand, measurements of the state of the resonators in a coherent state have allowed for the probabilistic generation of entanglement via measurement based protocols~\cite{Ruskov2003,Lalumiere2010,Roch2014,delva_measurement_2024}. To the best of our knowledge, early work has been performed which considers an initial squeezed state~\cite{matsuoka_entanglement_2013,Matsuoka2016}, a rigorous analysis at the level of entanglement -- involving the off diagonal element of the density matrix -- in measurement-based protocols with the inclusion of dynamical parametric squeezing and noise is still missing. We shall address this in the following.

Let us consider a realistic qubit coupled to a resonator. The creation and annihilation operators of the cavity of frequency $\omega$ are $a$ and $a^\dag$, $\sigma_j$ with $j\in\{x,y,z,\pm\}$ are the Pauli operators of the qubit, $\omega_q$ is the frequency of the qubit. Under the rotating-wave approximation, the interaction between qubit and resonator is described by the Jaynes-Cummings model with $g$ being the vacuum Rabi splitting. %\begin{equation}
%    \hat{H}_{\text{JC}} = \omega  \hat{a}^\dag \hat{a} + \sum_{i \in {1,2}} \left[  \omega_q \hat{\sigma}_z^{(i)}  + g \left( \hat{a}^\dag \hat{\sigma}_-^{(i)} + \hat{a} \hat{\sigma}_{+}^{(i)} \right) \right] \;, \nonumber
%\end{equation}
In the off-resonant regime (also known as the dispersive regime), i.e.,~for $\omega \gg \Delta \gg g$, where $\Delta = |\omega - \omega_q|$ is the detuning frequency, the effective Hamiltonian of the system reads
%\begin{align}
%\label{eq:example_1_hamitonian}
%    \frac{\hat{H}_1}{\hbar \omega}  =  &   \frac{\chi}{2}  \hat{\sigma}_z \hat{a}^\dag \hat{a} + \frac{i \lambda}{2} \left( \hat{a}^{\dag 2} - \hat{a}^2 \right) \\ &  \hspace{0.5cm} + \sqrt{\kappa} \left( \alpha_{in} \hat{a}^{\dag} + \alpha^*_{in} \hat{a} \right) \nonumber    + \frac{\chi}{2} \left( \hat{\sigma}^{(1)}_x \hat{\sigma}^{(2)}_x + \hat{\sigma}^{(1)}_y \hat{\sigma}^{(2)}_y\right) , 
%\end{align} 
\begin{align}
\label{eq:example_1_hamitonian}
    \frac{\hat{H}_1}{\hbar \omega}  =  &  \left(  \Delta + \frac{\mu}{2} \left( \hat{a}^\dag \hat{a} +\frac{1}{2} \right) \right) \otimes \hat{\sigma}_z  + \frac{i \lambda}{2} \left( \hat{a}^{\dag 2} - \hat{a}^2 \right),
\end{align} 
where we considered the rotating frame at frequency $\omega$. The first term represents a frequency shift conditional on the qubit state of strength given by the dispersive coupling $\mu = 4 g^2 / \Delta$ and the second is the squeezing of strength $\lambda$. % The third represents a driving with a coherent state $\ket{\alpha_{\text{in}}}$ of the environment which couples to the cavity at rate $\kappa$. The last is a direct qubit-qubit interaction, also known as XY-Hamiltonian, representing the exchange of virtual photons between the two qubits. 
The inclusion of realistic noise is done through the Master equation 
\begin{align}
\label{eq:example_1_master}
    \hspace{-0.1cm}\frac{\partial \hat{\varrho}}{\partial \tau} =&  \frac{i}{\hbar \omega} [ \hat{\varrho}, \hat{H}_1] +  i \sqrt{\kappa}  \left[ \hat{\varrho}, \alpha_{in} \hat{a}^{\dag} + \alpha^*_{in} \hat{a} \right]   \\&  \nonumber + \kappa \mathcal{D}[\hat{a}] (\hat{\varrho})  + \frac{1}{T_1}\mathcal{D}[\hat{\sigma}_-] (\hat{\varrho}) + \frac{1}{T_2}\mathcal{D}[\hat{\sigma}_z] (\hat{\varrho}) ,
\end{align}
which describes qubit dephasing and decay, cavity losses and a driving term~\cite{eddins_high_2019}. It is possible to recast this dynamics to our formalism via the Hamiltonians $H_q^{(0)} = 2 \mu j $,  $r_j = (0,0)^{\rm T} $, and $H_j= \lambda \sigma_x + j \mu \mathds{1}$ for $j \in \{-1,0,1\}$; for the noise terms, via $D = \kappa \mathds{1}$, $E = \kappa/2 \Omega_2$ and 
\begin{equation}\gamma_z = 
    \begin{pmatrix}
        \frac{1}{T}_1 & \frac{1}{2 T}_1+\frac{1}{T}_2 \\
        \frac{1}{2 T}_1+\frac{1}{T}_2 & - \frac{1}{T}_1
    \end{pmatrix}
\end{equation}
where we included the diagonal component of DV decay and take care to include the inhomogeneous source term induced by the decay in the DV. The driving term is, in canonical coordinates, $d = |\alpha_{in}| \sqrt{2 \kappa}(\cos \phi_d, \sin \phi_d )$, where $\alpha_{in}=|\alpha_{in}|{\rm e}^{i\phi_d}$. The initial state is a tensor product of the cavity ground state and $\ket{+}  = (\ket{1}+\ket{-1})/2$ DV state.

%While not explicitly diagonal in the qubit basis due to the XY-Hamiltonian term, the permutation invariance of the two qubits makes the system effectively behaving as a single spin with $J = 1$~\cite{dicke_coherence_1954} equipped with the $\hat{J}_z = \sum_{m\in \{ \pm 1, 0 \}} \ket{m} \bra{m}$ operator for which the dynamics (without DV decay) can be solved with our method. Specifically, the eigenstates are $\ket{m=-1} =  \ket{00}$,  $ \ket{m=1} = \ket{11}$ and $ \ket{m=0} = (\ket{01} + \ket{10})/\sqrt{2}$. Interestingly, the state degeneracy in the odd subspace ($\ket{m=0}$), while the $\ket{m = \pm 1}$ states show the frequency shifts of $\pm \mu$, have allowed, by measuring the cavity mode to probabilistically collapse the state into the odd subspace and, conditionally on specific outcomes, to generate an entangled state of the qubits~\cite{ }. 
 


%~\footnote{It is worth mentioning that the direct qubit-qubit interaction does not entangle this initial state.}. 

For the off-diagonal element, the dynamics can be solved via the doubling method of the Riccati equation as outlined in Sec.~\ref{sec:dynamics} in terms of the phase space parameters which are given in Appendix~\ref{app:examples}. The frequency shift induced by the qubit rotates the initial coherent state of the cavity at two frequencies conditionally on the qubit state. The squeezing and decaying of the cavity creates a GBCS with a stable steady-state, both superposed in the first moments as well as the squeezing parameter. This leads, on the diagonal elements, to distinct $\sigma^{\text{on}}_{\pm}$ and $\bar{r}^{\text{on}}_{\pm}$. By varying the driving quadrature, one can tune between squeezed ($\phi_d = \pi/2$) and amplified ($\phi_d = 0$) regimes, creating a superposition of first moments along a phase space direction either orthogonal or parallel to the squeezing direction~\cite{eddins_high_2019}. In Fig.~\ref{fig:},  we present the GBCS in the steady state for these two regimes. 

The dynamical creation of a superposition in the CV induces dephasing in the qubit state via two processes with rates immediately apparent from inspection of the QRDM ODE of Eq~(\ref{eq:general_unitary_ODE_off}): on the one hand, the difference in squeezing between the branches induces and effective dephasing at the time-dependent rate 
\begin{equation}
\Gamma_{\sigma}(\tau) =  \,\text{Im}\!\left[ \frac{1}{2}\text{Tr}\!\left[H_q \sigma^{\text{off}} \right] \right] = \frac{\mu}{2} \text{Tr} \left[ \sigma_\text{I} (\tau) \right]  \; ; \end{equation} 
on the other hand, the superposition in first moments gives a dephasing rate 
\begin{equation}
\Gamma_{\bar{r}} (\tau)= \text{Im}\left[ (\bar{r}^{\text{off}})^{\rm T}  H_q \bar{r}^{\text{off}} \right] = 2 \mu \bar{r}_{\text{R}}^{\rm T}(\tau) \bar{r}_{\text{I}}(\tau).
\end{equation}
The former is unavoidable as it stems from the non-commuting nature of squeezing and conditional rotations, while the latter depends on the driving, and can be minimized by choosing $\phi_d = 0$. The total decoherence rate of the off-diagonal element of the QRDM is then given by $\Gamma_{\text{tot}} (\tau) = \Gamma_{\sigma}(\tau) + \Gamma_{\bar{r}} (\tau) + \frac{1}{2 T}_1+\frac{1}{T}_2$ (see Fig.~\ref{fig:}), from which one can compute the total decay $\mathcal{C}_{\text{tot}} (\tau) = \int_{0}^\tau \text{d}t \Gamma_{\text{tot}} (t)$.  As $\tau \to \infty$, we recover the steady state decoherence rates, measured in experimental setting, of~\cite{eddins_high_2019}. Thus, our formalism naturally captures these transient processes and provides physical intuition of dynamical decoherence in terms of the time-dependent generalised complex moments: the imaginary part of the off-diagonal covariance matrix captures the not-overlapping features of the superposition at the level of the uncertainties, while the coupling between the real and imaginary component at the level of the expectation values. 


%We shall consider the homodyne detection of the quadrature $\hat{Q}$ of the cavity mode orthogonally to the driving mode, i.e. $\hat{Q} = -\sin(\phi_d) \hat{x} + \cos(\phi_d) \hat{p}$, which is the optimal for information gain. This can be described in our formalism by homodyne measurement with covariance of Eq~\ref{eq:homodyne} under $\phi = \phi_d +\pi/2$. 



%Via the use of the completeness of Pauli matrices, the last term can be rewritten as $\chi \left(\frac{1}{2}\hat{J}^2 - \frac{1}{2} \mathds{1} - \hat{\sigma}_z^{(1)} \hat{\sigma}_z^{(2)}\right)$, where $\hat{J}^2 $ is the total spin of the system. Hence, this additional term is diagonal in the qubit basis, and can thus be solved according to our methodology, leading to an overall phase ($\hat{J}^2$) and relative phases ($\hat{\sigma}_z^{(1)} \hat{\sigma}_z^{(2)})$. In other words, the eigmbasis of the 



 
\subsection{Stern-Gerlach Matter-Wave Interferometry: \\Squeezing, Diffusion, and Wigner Negativity \label{subsec:mass}}


The creation of non-classical states of levitated nanoparticles (NPs) is an active area of research, aiming to detect forces and entanglement for both applied and fundamental objectives~\cite{Romero_Isart_quantum_2011,bose_spin_2017,marletto_gravitationally-induced_2017}. In non-Gaussian matter-wave interferometry, large superpositions in the CV may be experimentally achieved via the linear coupling -- known as Stern-Gerlach interaction or spin-dependent force -- between a qubit and the NP, forming and recombining a GBCS. This interaction may be used to generate non-classical states in levitated masses with application to sensing forces and entanglement~\cite{bose_qubit_2006,Romero_Isart_quantum_2011,scala_matter_2013,geraci_sensing_2015,qvarfort_gravimetry_2018,marshman2020mesoscopic,eva_requirements_2023,barker2022entanglement,toros_relative_2021,bose2016matter,bose_spin_2017, marletto_gravitationally-induced_2017}.  The systems that are candidate for this type of couplings are dopant spins (such as Nitrogen-Vacancy centers) in diamagnetic NPs~\cite{scala_matter_2013,yin_large_2013,wan_free_2016,marshman_constructing_2022}, superconducting flux qubits with magnetic NPs~\cite{romero_isart_quantum_2012,nair_massive_2023} and atoms coupled to NPs via optical fields~\cite{toros_creating_2021}. While great advances have been made in the noise analysis of these systems~\cite{scala_matter_2013, pedernales_motional_2020,toros_loss_2020,toros_relative_2021,rijavec_decoherence_2021,tilly_qudits_2021,ma_torque_2021,schut_improving_2022,henkel_internal_2022,henkel_universal_2023,fragolino_decoherence_2023}, a full dynamical description of the experiments, with the inclusion of state preparation and diffusion noise in the CV degrees of freedom (one of the main sources of noise) is, to the best of our knowledge, still lacking in the literature. 

Here, the CV observables $\hat{X}$ and $\hat{P}$ describe the position and momentum of the centre of mass of the levitated NP, respectively. This degree of freedom is cooled to a thermal state (with $\sigma_T = (1+2N_p) = \coth({\hbar} \omega_t/ (2 k_B T_m))$ -- where $N_p$ is the number of phonons, $k_B$ is the Boltzmann constant and $T_m$ the temperature~\cite{serafini_quantum_2017,schlosshauer_quantum_2007}) of a \textit{tight} harmonic trap -- for instance, an optical tweezer of frequency $\omega_t$~\cite{delic_cooling_2020,magri_real_2021}. At $\tau = 0$, the cooling trap is rapidly switched off and the mass undergoes Stern-Gerlach interferometry in a dark trap, with frequency $\omega < \omega_t$, whose ground state spread is $x_0= \sqrt{\hbar/(2 M \omega)}$.  A qubit (of natural frequency $ \omega_q$) applies an operator-valued force of strength $F_q$ to the mass. The total Hamiltonian is 
\begin{equation}
\label{eq:hamiltonian_1}
    \hat{H}_{2} = \frac{\hbar \omega_q}{2} \hat{\sigma}_z + \frac{\hat{P}^2}{2M} + \frac{1}{2} M \omega^2 \hat{X}^2    -  F_q \hat{X}  \otimes \hat{\sigma}_z - F_u \hat{X},
\end{equation}
where the first term is the free evolution of the qubit, the second and third terms are the kinetic energy and harmonic trapping of the mass, the last two are the operator-valued and unknown forces ($F_u$), respectively. 

  The centre of mass is initially centred with respect to the dark trap and, given the change of frequency and imperfect cooling, the initial state is a thermal squeezed state of $N_p$ phonons and squeezing parameter $s = \omega/\omega_t$, i.e., its initial first and second moments are $r_0 = (0,\; 0)^{\rm T}$ and $\sigma_0 = (1+2N_p) \;\text{diag}(s,1/s)$. The qubit is initialised in the state $\varrho_q (0) = \ket{+}\bra{+}$.
By moving to dimensionless canonical observable, we can define the unitless forces $f_i = \sqrt{2} x_0 F_i/(\hbar \omega)$ with $i \in u, q$. The Hamiltonian of Eq.~(\ref{eq:hamiltonian_1}) can be written as  Eq.~(\ref{eq:Hamiltonian_force_only})  by setting $H_m = \mathds{1}$, $H_q = 0$, and $r_j = r_m + j r_q$, with $r_m=(f_u,\;0)^{\rm T}$ and $r_q = (f_q, \; 0)^{\rm T}$. During the evolution,  the qubit decays at rate $\Gamma_1$ and dephases at a rate $\Gamma_2$, while the CV undergoes diffusive dynamics in an Ohmic bath of thermal photons at high temperature $T_{p} \gg \omega/k_B$, such that the rate is at rate $\Gamma_x = 2 \gamma_0 k_B T_p/(\hbar \omega^2)$, where $\gamma_0$ is customarily measured in experimental settings~\cite{schlosshauer_quantum_2007}. Diffusive noise is equivalent to averaging over stochastic random forces experineced by the NP~\cite{weiss_large_2021}. The dynamics is then described by the Lindbladian 
\begin{align}
    \frac{\partial \hat{\varrho}}{\partial \tau} =  \frac{i}{\hbar \omega} [ \hat{\varrho}, \hat{H}_{2}] &+ \Gamma_x \left(  \hat{x} \hat{\varrho} \hat{x}  - \frac{1}{2} \{ \hat{\varrho}, \hat{x}^2 \} \right)  \\
    &\;\;\;\; + \frac{\Gamma_2}{2} \left(\hat{\sigma}_z \hat{\varrho} \hat{\sigma}_z  - \hat{\varrho} \right) + \frac{\Gamma_1}{2} \left(\hat{\sigma}_- \hat{\varrho} \hat{\sigma}_-  - \hat{\varrho} \right). \nonumber
\end{align}
\textcolor{red}{$2 \pi$?}
At first, let us solve the homogenous case (Tab~\ref{tab:open}) and consider the qubit decay as a perturbation theory (Sec.~\ref{sec:perturbation}). In this case, the dynamics is a specific case of Eq.~(\ref{eq:linbladian}), with 
\begin{equation}
    B = \begin{pmatrix}
         \Gamma_x &0\\
        0&0
    \end{pmatrix} \implies  D =  \begin{pmatrix}
        0 & 0\\
        0 &  2 \Gamma_x
    \end{pmatrix} .
\end{equation}
Given that the Hamiltonian has only linear interaction between the center of mass and the qubit ($H_q = 0$) and that $E =0$ (i.e., the noise has only the symmetric diffusive term $D$), the dynamics can be analytically solved via the single sympletic transformation $S_m(\tau) = {\rm e}^{\tau \Omega H_m}$ (as detailed in Sec.~\ref{sec:open}), which represents a phase space rotation of angle~$\tau$ (since $ H_m = \mathds{1}$). The time evolutions of the phase-space parameters are given in the Appendix \ref{app:examples}. The single covariance matrix evolves as a squeezed thermal state according to a Gaussian diffusive process, while the superposed first moments undergo the trajectories of \textcolor{red}{Fig.~\ref{fig:example_2}}. Furthermore, as discussed in Appendix \ref{app:examples}, the $T_1$ decay noise can also be treated exactly since the corresponding perturbation theory naturally truncates at first order.

\begin{figure*}[!t]
    \centering
    \includegraphics[width=1\linewidth]{Example_1.pdf}
    \caption{$\sigma_x$ measurement of a Stern-Gerlach Interferometry with a hypothetically unknown force ($f_u = 2$, $f_q = 0.5$): (a) Dynamical decay in contrast ($\mathcal{C}$) as a function of time ($\tau$) for unitary and diffusive ($\Gamma_x = 0.1$) dynamics with initial coherent state, thermal state ($N_p = 0.8$) and thermal squeezed state ($N_p = 0.8$, $s=2$); (b) Wigner functions of the two post-measurement states ($\tau = \pi$,  $N_p = 0.8$, $s=2$, $\Gamma_x = 0.02$) . \textcolor{red}{diffuive}} %of an initial squeezed thermal state undergoing diffusive Stern-Gerlach dynamics
    \label{fig:example_1}
\end{figure*}

Let us now investigate a $\sigma_x$-measurement, which can be used to detect the unknown force in terms of a qubit phase estimation protocol or generate Wigner negativities in the CV. The measurement outcomes $\pm1$ have associated projectors $\Pi_\pm = \ket{\pm}\bra{\pm}$, with $ \ket{\pm} = (\ket{1}\pm\ket{-1})/\sqrt{2}$ being the eigenstates of $\sigma_x$. The time evolution of the outcomes' probabilities is given by
\begin{align}
\label{eq:probability_SG}
    P_x(\pm, \tau) = \text{Tr}_q\left[\Pi_\pm \hat{\varrho}^q(\tau)\right] = \frac{1}{2} \left[ 1 \pm {e\rm}^{-\mathcal{C} (\tau)} \cos  \phi (\tau)  \right], \nonumber
\end{align}
where $\phi(\tau) = 2 f_q f_u (\tau +\sin \tau ) + \frac{1}{2} \omega_q \tau $ and
\begin{equation}
    \mathcal{C} (\tau) = \int_{0}^{\tau} \text{d}t \left( - 2 \text{Im} [r_q^{\rm T} r^{\text{off}} (t) ] + \frac{1}{2 T_1}  + \frac{1}{T_2}  \right) \;,
\end{equation}
whose analytical form (which can computed in general as in Table~\ref{tab:open}) is given in Appendix \ref{app:examples} and is plotted in Fig.~\ref{fig:example_1} for different initial states and dynamics. The time evolution of $\mathcal{C}$ shows the apparent decoherence of the qubit due to the entanglement with the CV degree of freedom, its value increasing at early times as the center of mass is superposed at two distinct locations. However, as shown in Fig.~\ref{fig:example_1} and mentioned in Sec.~\ref{sec:sup_phase_space}, the interference contrast shows a revival of coherence, when $\mathcal{C}$ decreases in value, reaching a local minimum at $\tau = 2 \pi$. At this time, as in the case of unitary dynamics of initial thermal states~\cite{scala_matter_2013}, the wave-matter interferometer performs one loop, also for the case of open dynamics with a squeezed thermal state. In fact,  the mass is recombined and centered in the trap again with $\bar{r}^{\text{on}}_{\pm} (2 \pi) = (0,0)^{\rm T}$, thus representing a full recombination of the interferometer without loss of visibility due to the so-called ``Humpty-Dumpty effect''~\cite{englert_is_1988, japha_quantum_2023, xiang2024phononinducedcontrastmatter, zhou2024gyroscopicstabilitynanoparticlessterngerlach}. However, due to the inclusion of the diffusion term, the first moments of the off-diagonal terms are $\bar{r}^{\text{off}} (2 \pi) = (0,2 i \pi f_q \Gamma_x)^{\rm T}$. Being not zero and complex, this leads to an exponential decay of the visibility of the qubit measurement, caused by the diffusion process, with $\mathcal{C}(2 \pi) = 6 \pi  f_q^2 \Gamma_x + \pi \Gamma_1 + 2 \pi \Gamma_2 $.  Interestingly, $\mathcal{C}(2 \pi)$ is found to be independent of the squeezing parameter =$s$ and the initial number of phonons $N_p$, but linear in the diffusion and dephasing rates $\Gamma_i$ (with $i\in x,z$) and quadratic in the superposition size $f_q$.  It is remarkable how our formalism allows one to pinpoint the origin of such effects affecting the dynamics: similarly to the previous example, for linear couplings, the imaginary part of the off-diagonal first moments gives the time-dependent decoherence rate for the qubit $\Gamma (\tau) = -2 r_q^{\rm T} r^{\text{off}} (\tau)$ \textcolor{red}{-?}

The maximum superposition is achieved at $\tau = \pi$, given by $\delta X = \bar{r}^{\text{on}}_{+}(\pi) -  \bar{r}^{\text{on}}_{-}(\pi) = 4 f_q$ in ground spread units. At this moment, if a $\sigma_x$ measurement occurs, the loss of visibility is $ \mathcal{C} (\pi) = f_q^2 \left(4 (1 +2 N_p)/s + 3 \pi \Gamma_x \right) +  \pi \Gamma_2 $, which is still present even in the ideal case ($N_p = \Gamma_x = \Gamma_2 = 0$). This is due to the fact that part of the information is retained within the CV quantum state and the qubit has decohered with $P_x (\pm, \pi) \approx 1/2$.  Summing over elements of Eq.~(\ref{eq:wigner_jk}), the post-measurement states of the mode have Wigner functions $\mathcal{W}_{\pm} = \mathcal{W}_c +  \mathcal{W}_\Gamma \pm \mathcal{W}_q $. $\mathcal{W}_c$ and $\mathcal{W}_q$ are the classical and quantum component discussed in Sec.~\ref{sec:wigner} and analytically computed in Appendix~\ref{app:examples} for this dynamics.   Finally, $\mathcal{W}_\Gamma$ is the component of Wigner function due to the DV decay, and it is the Fourier transform of the inhomogeneous characteristic function $\chi^{(1)} (\tau = \pi, \tilde{\xi})$  (given in Appendix~\ref{app:examples} as a series of Bessel functions). In Fig.~\ref{fig:example_1}, we plot the different contributions of the Wigner functions and the total Wigner functions for the post measurements states.  The interesting quantity that should be noted is the Wigner phase
\begin{equation}
    \phi_W = 4 f_q \left(2 \tilde{p} - \frac{4 \Gamma_x  \left(\tilde{x}-2 f_u\right)}{\pi  \Gamma_x+  (1 +2 N_p)s}- \frac{\pi  \Gamma_x  \tilde{p} }{\pi  \Gamma_x + (1 + 2 N_p)/s }\right) \nonumber,
\end{equation}
illustrating how the fringes' spacing decreases with the superposition size (proportional to $f_q$). One then can envision the application of this analytical solution to identify a coupling $f_q$, which is function of experimentally accessible quantities, that creates superpositions large enough to resolve the two-Gaussian statistical mixture $\mathcal{W}_c$ from the fringes, but also small enough to keep the fringes resolvable under the considered measurement apparatus (for instance, homodyne detection of scattered light).


\subsection{Comparative Analysis with Known Methods}


Having showcased the vast appearance of GBCS among quantum protocols and experiments, we shall discuss GBCS's parametrization in the context of existing techniques for the modelling of DV-CV hybrid system. Numerical solutions rely on truncating the CV Hilbert space in the $n$-basis or on other appropriate spectral decompositions~\cite{jaynes_comparison_1963,Braak_Integrability_2011,chen_exact_2012,Batchelor_integrability_2015}, making each mode, effectively, a DV variable subsystem of dimensionality $N_t$ (where $N_t$ is the truncation number). The numerical nature of these methods equips them with great flexibility: being able to solve, at least in principle, any hybrid dynamics with any initial-state. However, when either the states involved approaches high quantum numbers (as it is always the case with significant squeezing) or the Hamiltonian approaches instabilities or large differences between couplings, the required $N_t$ to accurately simulate the system increases. To approximately solve an $n$-mode-qu$d$it open quantum system, the state parameters scales as $\mathcal{O}( N_t^{2 n} \times d^2)$ for full-density matrix treatments, $\mathcal{O}( N_t^{n} \times d^2)$ for trajectories methods (which however requires sampling) and $\mathcal{O}( N_t^{2} \times d^2)$ for tensor network. By restricting at the class of operator-valued Gaussian Hamiltonian, which often are the non-linearities common in experimental practice, the formalism presented in this work shows that it is possible to parametrize the state with $\mathcal{O}( n^2 \times N^2)$ parameters in this restricted regime, and hence, showing an exponential improvement in efficiency, compounded with the advantage of an exact, or semi-analytical solution (by the latter, we refer to case where the Riccati equation for CM may need numerical solutions: these are however much more reliable, and allow for a much more accurate quantification of errors, than the outcome of Hilbert space truncated dynamics). Notice that, in general, Hilbert space truncations are prone to issues and complications~\cite{ashhab26}, which our methods avoid altogether. 

The types of interactions and states presented in this work have been extensively studied in analytical treatments present the literature, but limited to specific examples. General statements have been achieved for linear-coupling whose resulting unitary dynamics can be decomposed in a Gaussian unitary and conditional displacement (equivalent to Tab~\ref{tab:force}). This is due to the fact that in the unitary case of linear coupling, expansions like to Magnus one, or Lie algebraic approaches, naturally truncate. Yet, these analyses fail to account for Gaussian noise in the general setting (Tab~\ref{tab:open}).  Similar is the case for the quadratic coupling, with many instances treating either unitary dynamics or limited to initial coherent states.

Let us also mention that, after the appearance of the first version of this study, analogous equations were derived in the context of generalised master equations, without, however drawing connections with hybrid systems and dynamics.
%~\footnote{Yet, the presence of these recent works subsequent of the first version of this one, enlights the relevance of our general treatment}. 
Our state parametrization encompasses all these specific dynamics, by mapping the dynamics to general complex Riccati equation (and providing a systematic procedure to solve them), alongside the action of a vast range of experimentally accessible measurements and operations. 
\textcolor{blue}{}Cite Marco/Francesco here,
~\cite{serafini_minimum_2004,Entanglement_Wahyu_2008}
~\cite{utami_entanglement_2008}



\section{Conclusions\label{sec:conclusion}}


Quantum experiments and computational tasks are often achieved via a successions of dynamics, measurements, and operations, which creates and manipulates a superposition of CV Gaussian state entangled with a DV subsystem. Our work presents a formalism to analytical treat these processes and discuss some of the many instances where these states arises in the quantum world. The presented state parametrization of GBCS is a complex generalization of first moments and covariance matrices. These quantities are operationally meaningful, in the sense that, the state is defined via expectation values and uncertainties often measured in experimental practice. Their real and imaginary parts are directly related to coherence and phases of the DV state and to the defining properties of Wigner functions of superposed Gaussian states. This allows visualization of the CV parameters, as in Gaussian quantum mechanics, by extending the phase space to a complex space. We have illustrated the applicability to modelling experimental setting through two case studies: (a) dynamical effects in dispersive qubit red-out experiments via a parametric amplified resonator, and (b) the evolution of a levitated nanoparticle in a Stern-Gerlach interferometer with diffusive environment for phase estimation and the certification of Wigner negativity.

Given its flexibility and clear operational grounding, our formalism lends itself to bespoke experimental predictions, spanning a vast range of systems relevant to quantum technologies and fundamental physics, which calls for further investigations. Our results opens a systematic avenue to analytically or semi-analytically treat processes which where before computed numerically. For instance, the open-driven dynamics may allow for optimization procedure to minimize unwanted noise via pulse design (for instance, superconducting or ion-based quantum systems). The linear scaling of the parameters in the number of resonator, while having the same scaling for the number of qubit, could lead to large simulation of analogue implementations of multi-qubit gates. As shown in the first example. the dynamical analysis may shine light on transitional behaviour of quantum system which have been, until now, described at steady state. On the fundamental applicability, quantum field theories with superposed Gaussian states could be described via our methodology via a truncation to an high number of modes (as it has been routinely done for Gaussian QFTs). 

To summarise, our work represents a generalisation of Gaussian techniques to non-Gaussian processes, i.e. superpositions of Gaussian processes, by pointing out that cats have non-linear complex moments. 


\textcolor{blue}{
A dedicated Python software library for the solutions of all the dynamics presented (all tables) and measurements, together with the examples presented, is shared on \href{https://github.com/LorenzoBracciniUCL/GBCS}{GitHub}. Furthermore, the analytical results of Table~\ref{tab:force} and ~\ref{tab:open_force} are also provided as Mathematica Notebooks. }


%The presented formalism is a general framework for treating dynamics involving operator-valued Gaussian interaction Hamiltonians between qubits and CV systems, hinging on a phase space, covariance matrix approach. Our method allows for the exact treatment of the arising dynamics, open or unitary, as well as of measurements. These non-linear dynamics and joint measurements lead to superpositions of Gaussian processes and, thus, to cat-like entangled states, which we dubbed Gaussian-Branched Cat States. 






\section{Acknowledgments}

\textcolor{red}{We thank Harv }L.B. would like to acknowledge Engineering and Physical Sciences Research Council (EPSRC) grants (EP/R513143/1 and EP/W524335/1). The research of
The work of S.B. is funded in part by the Gordon and Betty Moore Foundation through
Grant GBMF12328, DOI 10.37807/GBMF12328, and by the Alfred P. Sloan Foundation under Grant
No. G-2023-21130. A.S. acknowledges financial support from the Leverhulme Trust through
Research Project Grant RPG-2024-287.
 
\bibliography{apssamp}% Produces the bibliography via BibTeX.

\appendix

\onecolumngrid

\section{Map from Hilbert to Phase Space \label{app:map_phase_space}}

By repetitive use of Eqs.~(\ref{eq:phase_space_corresp1},\ref{eq:phase_space_corresp2}), one finds 
\begin{equation}
\label{eq:map_chara_first}
   i \hat{r}^m \hat{\varrho}  \longleftrightarrow  \left(  \partial_{\tilde{\xi}^m} + \frac{i}{2} \Omega^{mn} \tilde{\xi}^n \right) \chi (\tilde{\xi}) \hspace{1cm}  i  \hat{\varrho} \hat{r}^m \longleftrightarrow  \left(  \partial_{\tilde{\xi}^m} - \frac{i}{2} \Omega^{mn} \tilde{\xi}^n \right) \chi (\tilde{\xi})
\end{equation}
which will be useful for the Hamiltonians linear terms. Instead, for the quadratic ones, one can note that  
\begin{align*}
     \partial_{\tilde{\xi}^n} \partial_{\tilde{\xi}^m} \chi (\tilde{\xi}) &\longleftrightarrow  - \frac{1}{4} \left(\hat{r}^n ( \hat{r}^m  \hat{\varrho} + \hat{\varrho} \hat{r}^m  ) +  (\hat{r}^m  \hat{\varrho} + \hat{\varrho} \hat{r}^m ) \hat{r}^n  \right)  
\end{align*}
\begin{align*}
    \tilde{\xi}^n \partial_{\tilde{\xi}^m} \chi (\tilde{\xi}) &\longleftrightarrow  \frac{i}{2} \Omega^{on} \left(\hat{r}^o ( \hat{r}^m  \hat{\varrho} + \hat{\varrho} \hat{r}^m  ) -  (\hat{r}^m  \hat{\varrho} + \hat{\varrho} \hat{r}^m ) \hat{r}^o  \right)  
\end{align*}
\begin{align*}
    \tilde{\xi}^n \tilde{\xi}^m \chi (\tilde{\xi}) &\longleftrightarrow  \Omega^{on} \Omega^{pm} \left(\hat{r}^o ( \hat{r}^p  \hat{\varrho} - \hat{\varrho} \hat{r}^p  ) -  (\hat{r}^p  \hat{\varrho} - \hat{\varrho} \hat{r}^p ) \hat{r}^o  \right) 
\end{align*}
%Furthermore, one may note that
%\begin{align*}
%     \partial_{\tilde{\xi}^m} \left( \tilde{\xi}^n \chi (\tilde{\xi}) \right) = \delta^{nm} \chi +  \tilde{\xi}^n \partial_{\tilde{\xi}^m}  \chi (\tilde{\xi}) &\longleftrightarrow \delta^{nm} \hat{\varrho} +  \frac{i}{2} \Omega^{on} \left(\hat{r}^o ( \hat{r}^m  \hat{\varrho} + \hat{\varrho} \hat{r}^m  ) -  (\hat{r}^m  \hat{\varrho} + \hat{\varrho} \hat{r}^m ) \hat{r}^o  \right)  
%\end{align*}
such that inverting 
\begin{align*}
     (a) & \;\; \hat{r}^n  \hat{r}^m  \hat{\varrho} + \hat{\varrho} \hat{r}^m  \hat{r}^n + \hat{r}^n  \hat{\varrho}  \hat{r}^m  +  \hat{r}^m  \hat{\varrho} \hat{r}^n \longleftrightarrow   - 4 \partial_{\tilde{\xi}^n} \partial_{\tilde{\xi}^m} \chi (\tilde{\xi})  \\
      (b) & \;\; \hat{r}^n  \hat{r}^m  \hat{\varrho} - \hat{\varrho} \hat{r}^m  \hat{r}^n + \hat{r}^n  \hat{\varrho}  \hat{r}^m  -  \hat{r}^m  \hat{\varrho} \hat{r}^n   \longleftrightarrow   - 2 i \Omega^{n o} \tilde{\xi}^o  \partial_{\tilde{\xi}^m} \chi (\tilde{\xi}) \\
       (c) &\;\; \hat{r}^n  \hat{r}^m  \hat{\varrho} + \hat{\varrho} \hat{r}^m  \hat{r}^n -  \hat{r}^n  \hat{\varrho}  \hat{r}^m  -  \hat{r}^m  \hat{\varrho} \hat{r}^n    \longleftrightarrow \Omega^{n o} \Omega^{m p} \tilde{\xi}^o \tilde{\xi}^p \chi (\tilde{\xi}) 
\end{align*}
Multiplying each expression for a symmetric matrix $A$, i.e. contracting with  $A^{mn}$:
\begin{align}
\label{eq:sim_phase_master}
     A \cdot &(a) \implies \;\; A^{mn} \left( \hat{r}^n  \hat{r}^m  \hat{\varrho} + \hat{\varrho} \hat{r}^m  \hat{r}^n + 2 \hat{r}^n  \hat{\varrho}  \hat{r}^m \right) \longleftrightarrow   - 4 A^{mn} \partial_{\tilde{\xi}^n} \partial_{\tilde{\xi}^m} \chi (\tilde{\xi}) \nonumber \\ 
      A \cdot &(b) \implies \;\;  A^{mn} \left( \hat{r}^n  \hat{r}^m  \hat{\varrho} - \hat{\varrho} \hat{r}^m  \hat{r}^n  \right) \longleftrightarrow   - 2 i  A^{mn} \Omega^{n o} \tilde{\xi}^o  \partial_{\tilde{\xi}^m} \chi (\tilde{\xi})  \nonumber\\
       A \cdot &(c) \implies \;\;  A^{mn} \left( \hat{r}^n  \hat{r}^m  \hat{\varrho} + \hat{\varrho} \hat{r}^m  \hat{r}^n - 2 \hat{r}^n  \hat{\varrho}  \hat{r}^m   \right)  \longleftrightarrow A^{mn} \Omega^{n o} \Omega^{m p} \tilde{\xi}^o \tilde{\xi}^p \chi (\tilde{\xi}) 
\end{align}
We arrive to the equality needed to describe operator-valued Hamiltonian terms:
%\begin{equation}
%    A \cdot (a - c) \implies A^{mn} \hat{r}^n  \hat{\varrho}  \hat{r}^m \leftrightarrow   A^{mn} \left[ - \partial_{\tilde{\xi}^n}\partial_{ \tilde{\xi}^m}  + \frac{1}{4} \Omega^{n o} \Omega^{m p} \tilde{\xi}^o \tilde{\xi}^p  \right] \chi (\tilde{\xi}) 
%\end{equation}
\begin{equation}
\label{eq:map_chara_second_1}
    A \cdot (a + 2 b + c) \implies A^{mn} \hat{r}^n  \hat{r}^m  \hat{\varrho} \leftrightarrow   A^{mn} \left[ -\partial_{\tilde{\xi}^n}\partial_{\tilde{\xi}^m}   -  i   \Omega^{n o} \tilde{\xi}^o  \partial_{\tilde{\xi}^m}  + \frac{1}{4} \Omega^{n o} \Omega^{m p} \tilde{\xi}^o \tilde{\xi}^p  \right] \chi (\tilde{\xi}) 
\end{equation}
\begin{equation}
\label{eq:map_chara_second_2}
    A \cdot (a - 2 b + c) \implies A^{mn}  \hat{\varrho} \hat{r}^m  \hat{r}^n   \leftrightarrow   A^{mn} \left[ -\partial_{\tilde{\xi}^n}\partial_{\tilde{\xi}^m}   +  i   \Omega^{n o} \tilde{\xi}^o  \partial_{\tilde{\xi}^m}  + \frac{1}{4} \Omega^{n o} \Omega^{m p} \tilde{\xi}^o \tilde{\xi}^p  \right] \chi (\tilde{\xi}) 
\end{equation}
For diffusive dynamics, it will be uesfull to multiply for a constant skey-symmetric matrix $B$ with elements $B^{mn}$ 
\begin{align}
\label{eq:anti_phase_master}
     B \cdot (a) \implies \;\; B^{mn} \left( \hat{r}^n  \hat{r}^m  \hat{\varrho} + \hat{\varrho} \hat{r}^m  \hat{r}^n  \right) \longleftrightarrow   - 4 B^{mn} \partial_{\tilde{\xi}^n} \partial_{\tilde{\xi}^m} \chi (\tilde{\xi}) 
\end{align}
\begin{align}
\label{eq:antsim_phase_master}
      B \cdot (b) \implies \;\;  B^{mn} \left( \hat{r}^n  \hat{r}^m  \hat{\varrho} - \hat{\varrho} \hat{r}^m  \hat{r}^n  + 2 \hat{r}^n  \hat{\varrho}  \hat{r}^m \right) \longleftrightarrow   - 2 i  B^{mn} \Omega^{n o} \tilde{\xi}^o  \partial_{\tilde{\xi}^m} \chi (\tilde{\xi}) 
\end{align}
\begin{align*}
       B \cdot (c) \implies \;\;  B^{mn} \left( \hat{r}^n  \hat{r}^m  \hat{\varrho} + \hat{\varrho} \hat{r}^m  \hat{r}^n  \right)  \longleftrightarrow B^{mn} \Omega^{n o} \Omega^{m p} \tilde{\xi}^o \tilde{\xi}^p \chi (\tilde{\xi}) 
\end{align*}


\section{Derivation of ODEs for Unitary Evolution \label{app:unitary_general}}

The unitary dynamics in frequency-unit time ($\tau$) is given by the von Neumann equation which, recalling the definitions the state in Eq.~(\ref{rhojk}) and the Hamiltonian of Eq.~(\ref{eq:Hamiltonian_max}), can be written in the qubit basis as
\begin{equation}
\label{eq:vonNeumann_master}
    \frac{\partial \hat{\varrho}}{\partial \tau} =  \frac{i}{\hbar \omega} \left[ \hat{\varrho} , \hat{H} \right] \Leftrightarrow  \frac{\partial \hat{\varrho}_{jk}}{\partial \tau} = \frac{i}{2} \left(  H_k^{mn} \hat{\varrho}_{jk} \hat{r}^m \hat{r}^n  - H_j^{mn} \hat{r}^m \hat{r}^n \hat{\varrho}_{jk}  \right) + i r_j^m \hat{r}^m \hat{\varrho}_{jk} - i r_k^m  \hat{\varrho}_{jk}  \hat{r}^m - \frac{i}{2} (H_j^0-H_k^0)\hat{\varrho}_{jk} \;.
\end{equation}
As claimed in the main text, we seek a solution of the form 
\begin{equation*}
    \hat{\varrho}(\tau) = \frac{1}{(2\pi)^n}\sum_{j, k  = 0}^{N-1}  \int_{\mathds{R}^{2n}} \text{d} \xi \; \chi_{jk} (\xi ,\tau)  D_{r} \otimes \ket{j} \bra{k}
\end{equation*}
where $\chi_{jk} (\xi ,\tau) = \text{Tr} [ D_{- r}  \hat{\varrho}_{jk} (\tau)] $ are four characteristic functions. Define $\tilde{\xi} = \Omega r$. From Eq.~(\ref{eq:map_chara_second_1}) and~(\ref{eq:map_chara_second_2}), one finds that 
\begin{equation*}
    \frac{i}{2} H^{mn}_j \hat{r}^n  \hat{r}^m  \hat{\varrho} \longleftrightarrow  \frac{1}{2} H^{mn}_j \left( - i \partial_{\tilde{\xi}^n} \partial_{\tilde{\xi}^m} + \Omega^{n o} \tilde{\xi}^o  \partial_{\tilde{\xi}^m} + \frac{i}{4} \Omega^{n o} \Omega^{m p} \tilde{\xi}^o \tilde{\xi}^p \right)  \chi (\tilde{\xi})
\end{equation*}
\begin{equation*}
    \frac{i}{2} H^{mn}_k \hat{\varrho}  \hat{r}^n  \hat{r}^m  \longleftrightarrow \frac{1}{2}  H^{mn}_k \left( - i \partial_{\tilde{\xi}^n} \partial_{\tilde{\xi}^m} -  \Omega^{n o} \tilde{\xi}^o  \partial_{\tilde{\xi}^m} + \frac{i}{4} \Omega^{n o} \Omega^{m p} \tilde{\xi}^o \tilde{\xi}^p \right)  \chi (\tilde{\xi})
\end{equation*}
and, from Eq.~(\ref{eq:map_chara_first}), 
\begin{equation*}
   i r_j^m \hat{r}^m \hat{\varrho}  \longleftrightarrow  r_j^m \left(  \partial_{\tilde{\xi}^m} + \frac{i}{2} \Omega^{mn} \tilde{\xi}^n \right) \chi (\tilde{\xi}) \hspace{2cm}  i  r^m_k \hat{\varrho} \hat{r}^m \longleftrightarrow  r^m_k \left(  \partial_{\tilde{\xi}^m} - \frac{i}{2} \Omega^{mn} \tilde{\xi}^n \right) \chi (\tilde{\xi})
\end{equation*}
Thus, from Eq.~(\ref{eq:vonNeumann_master}), the time evolution of the characteristic functions is described by the  PDE
 \begin{align}
 \label{eq:PDE_chara_general}
     \dot{\chi}_{jk} &= \Big[ \frac{i}{2} (H_j - H_k)^{mn} \partial_{\tilde{\xi}^n} \partial_{\tilde{\xi}^m}  - \left(  \frac{1}{2}  ( H_j + H_k )^{mn} \Omega^{n o} \tilde{\xi}^o  - (r_j - r_k)^m \right)  \partial_{\tilde{\xi}^m} \nonumber \\
     \hspace{2cm } &\hspace{3cm}- \frac{i}{8} (H_j - H_k)^{mn} \Omega^{n o} \Omega^{m p} \tilde{\xi}^o \tilde{\xi}^p  + \frac{i}{2} (r_j + r_k)^m  \Omega^{mn} \tilde{\xi}^n - \frac{i}{2} (H_j^0-H_k^0) \Big] \chi_{jk} . 
 \end{align}
Let us consider the ansatz
\begin{equation}
    \chi_{jk} (\tilde{\xi}, \tau) =  \exp \left( -\frac{1}{4} \tilde{\xi}^{\rm T} \sigma_{jk} (\tau) \tilde{\xi}  +  i  \tilde{\xi}^{\rm T} \bar{r}_{jk} (\tau) - \mathcal{C}_{jk} (\tau)  +  i \phi_{jk} (\tau) \right) 
\end{equation}
such that its derivatives are
\begin{equation}
    \partial_{\tilde{\xi}^n} \chi_{jk} (\tilde{\xi}, \tau) = - \left( \frac{1}{2} \sigma_{jk}^{nm} \tilde{\xi}^{m} -  i r^n_{jk}  \right) \chi_{jk} (\tilde{\xi}, \tau)
\end{equation}
\begin{align*}
    \partial_{\tilde{\xi}^n} \partial_{\tilde{\xi}^m} \chi_{jk} (\tilde{\xi}, \tau) &= - \frac{1}{2} \delta^{mn} \sigma_{jk}^{mn} \chi_{jk} (\tilde{\xi}, \tau) + \left( \frac{1}{2} \sigma_{jk}^{pm} \tilde{\xi}^{p} -  i r^m_{jk}  \right) \left( \frac{1}{2} \sigma_{jk}^{on} \tilde{\xi}^{o} -  i r^n_{jk}  \right) \chi_{jk} (\tilde{\xi}, \tau) \\ 
    &= \left[\frac{1}{4} \sigma_{jk}^{pm}  \sigma_{jk}^{on} \tilde{\xi}^{p} \tilde{\xi}^{o} - \frac{i}{2} (\sigma_{jk}^{pm}  r^n_{jk} +  \sigma_{jk}^{pn}  r^m_{jk} ) \tilde{\xi}^{p}  - \left(\bar{r}_{jk}^m \bar{r}_{jk}^n + \frac{1}{2}  \delta^{mn} \sigma_{jk}^{mn} \right) \right]\chi_{jk}  
\end{align*}
\begin{equation*}
    \partial_{\tau} \chi_{jk} = \left( -\frac{1}{4} \tilde{\xi}^{\rm T} \dot{\sigma}_{jk} (\tau) \tilde{\xi}  +  i \tilde{\xi}^{\rm T} \dot{\bar{r}}_{jk} (\tau) + \dot{\bar{r}}_{jk}^{(0)}(\tau) \right) \chi_{jk} (r, \tau)
\end{equation*}
and, substituting, one gets
 \begin{align*}
     \dot{\chi}_{jk} &= \Big[ \frac{i}{2} (H_j - H_k)^{mn} \left[ \frac{1}{4} \sigma_{jk}^{pm}  \sigma_{jk}^{on} \tilde{\xi}^{p} \tilde{\xi}^{o} - \frac{i}{2} (\sigma_{jk}^{pm}  r^n_{jk} +  \sigma_{jk}^{pn}  r^m_{jk} ) \tilde{\xi}^{p}  - \left(\bar{r}_{jk}^m \bar{r}_{jk}^n + \frac{1}{2}  \delta^{mn} \sigma_{jk}^{mn} \right) \right] 
     \\
     &\hspace{1cm}
    + \left(  \frac{1}{2}  ( H_j + H_k )^{mn} \Omega^{n o} \tilde{\xi}^o  - (r_j - r_k)^m \right)  \left( \frac{1}{2} \sigma_{jk}^{pm} \tilde{\xi}^{p} -  i r^m_{jk}  \right)  
     \\
     &\hspace{2cm }
     - \frac{i}{8} (H_j - H_k)^{mn} \Omega^{n o} \Omega^{m p} \tilde{\xi}^o \tilde{\xi}^p  + \frac{i}{2} (r_j + r_k)^m  \Omega^{mn} \tilde{\xi}^n - \frac{i}{2} (H_j^0-H_k^0) \Big] \chi_{jk} 
     \\
\end{align*}
\begin{align*}
     \dot{\chi}_{jk} &= \Big[ +  \frac{1}{4} 
     \left[ ( H_j + H_k )^{mn} \Omega^{no}  \sigma_{jk}^{pm}  +    \frac{i}{2} (H_j - H_k)^{mn} \left( \sigma_{jk}^{pm}  \sigma_{jk}^{on} +  \Omega^{n o} \Omega^{pm}  \right)   \right] \tilde{\xi}^{p} \tilde{\xi}^{o}  
     \\
     &\hspace{1cm} 
     +   i \tilde{\xi}^p  \left[   \frac{1}{2} \Omega^{p  n} ( H_j + H_k )^{nm}  \bar{r}_{jk}^m    - \frac{1}{2}\Omega^{p m}  (r_j + r_k)^m  - i  \left( \frac{1}{2} (H_j - H_k)^{mn}  \sigma_{jk}^{pm}  r^n_{jk}  -   \frac{1}{2} \sigma_{jk}^{pm} (r_j - r_k)^m \right) \right]  
     \\
     &\hspace{2cm} 
     -  i \left(  \frac{1}{2} (H_j - H_k)^{mn} \left(\bar{r}_{jk}^m \bar{r}_{jk}^n + \frac{1}{2}  \delta^{mn} \sigma_{jk}^{mn} \right) - (r_j - r_k)^m  r^m_{jk}  + \frac{1}{2} (H_j^0-H_k^0)  \right) \Big] \chi_{jk} \\
    &= \left( -\frac{1}{4} \tilde{\xi}^{\rm T} \dot{\sigma}_{jk} (\tau) \tilde{\xi}  +  i \tilde{\xi}^{\rm T} \dot{\bar{r}}_{jk} (\tau) + \dot{\bar{r}}_{jk}^{(0)}(\tau) \right) \chi_{jk} 
\end{align*}
where $\dot{\bar{r}}_{jk}^{(0)} = - \dot{\mathcal{C}}_{jk} (\tau) + i \dot{\phi}_{jk} (\tau)$ is understood for the reduced density matrix. The last equality must hold for all $\tilde{\xi}$. Thus, one finds the set of ODE of Table \ref{tab:unitary} by equating the coefficients of each power in $\tilde{\xi}^p$ of the last equality. %, such that
%\begin{align}
%    \dot{\sigma}_{jk} =   \frac{1}{2}   \Omega  ( H_j + H_k )  \sigma_{jk}  - \frac{1}{2} \sigma_{jk}  ( H_j + H_k ) \Omega  - i  \left[ \sigma_{jk} (H_j - H_k) \sigma_{jk}  + \Omega (H_j - H_k) \Omega \right]
%\end{align}
%\begin{align}
%    \dot{\bar{r}}_{jk} =   \frac{1}{2}  \left[   \Omega ( H_j + H_k )  \bar{r}_{jk}   -  i   \sigma_{jk} (H_j - H_k)  \bar{r}_{jk}   -  \Omega (r_j + r_k) +  i  \sigma_{jk} (r_j - r_k)   \right]
%\end{align}
%\begin{align}
%    \dot{\bar{r}}_{jk}^{(0)} = - i \left( \frac{1}{2}  \bar{r}_{jk}^{\rm T} (H_j - H_k)  \bar{r}_{jk} + \frac{1}{4} \text{Tr} [(H_j - H_k)\sigma_{jk}] - (r_j - r_k)^{\rm T}  \bar{r}_{jk}  + \frac{1}{2} (H_j^0-H_k^0) \right)
%\end{align}
%as claimed in the main text. 



\section{Integration of On Diagonal Terms and Linear Operator Valued Interactions: Unitary Dynamics \label{app:only_force}}

In order to integrate the operator valued force, we need to recall some properties: a general symplectic transformation is of the form $S_H(\tau) = \exp( \Omega H \tau)$, such that  $S_H^{-1} (\tau) = \exp (- \Omega H \tau) $ and $S_H^{\rm T} (\tau) = \exp (- H \Omega \tau) $. The symplectic form $\Omega$ respects $\Omega^{-1} = \Omega^{\rm T} = - \Omega$. Furthermore, from the definition of symplectic group (as the set of real matrices $S \in Sp_{2, \mathds{R}}$ such that $S \Omega S^{\rm T} = \Omega$), that $S \Omega  = \Omega (S^{\rm T})^{-1}$, $ \Omega S^{\rm T}  = S^{-1}  \Omega$, and $(S^{-1})^{\rm T}  \Omega = \Omega S $, from which it is possible to derive one of the central identities of the following proofs:
\begin{equation}
\label{eq:identity_usefull}
    \left( S^{\rm T}_H(\tau) - \mathds{1} \right)\Omega =  S^{\rm T}_H(\tau) \left( \mathds{1} - (S^{\rm T}_H)^{-1}(\tau)\right)\Omega = - S^{\rm T}_H(\tau) \Omega (S_H(\tau) - \mathds{1})
\end{equation}
and so its transpose reads $\Omega^{\rm T} \left( S_H(\tau) - \mathds{1} \right) =- (S_H^{\rm T}(\tau) - \mathds{1})\Omega^{\rm T} S_H(\tau) $. Let us also compute the integral
\begin{equation}
    \int_0^\tau \text{d}t S_H(t)  =\int_0^\tau \text{d}t {\rm e}^{ \Omega H t}  =  ( \Omega H)^{-1} \left(S_H(\tau) - \mathds{1} \right) = - H^{-1} \Omega  \left(S_H(\tau) - \mathds{1} \right),
\end{equation}
which implies  
\begin{equation}
    \int_0^\tau S_H^{-1}(t) dt =   \left(S_H^{-1}(\tau) - \mathds{1} \right) H^{-1} \Omega \hspace{1cm} \int_0^\tau S_H^{\rm T}(t) dt =     \left(S_H^{\rm T}(\tau) - \mathds{1} \right) \Omega  H^{-1}
\end{equation}

 Regarding the on diagonal terms, Eq.~(\ref{eq:general_unitary_ODE_on}) is readily integrable:  letting $H \to H_\pm = H_m \pm H_q$, one has
\begin{align*}
    \bar{r}^{\text{on}}_{\pm 1, \pm 1} &= S_{\pm} (\tau)  \bar{r}^{\text{on}} -  S_{\pm} (\tau)  \int_{0}^{\tau} \text{d} t S_{\pm}^{-1} (t) \Omega (r_m \pm r_q) = S_{\pm} (\tau)  \bar{r}^{\text{on}} - S_{\pm} (\tau)  \left( S_{\pm}^{-1} (\tau)  - \mathds{1} \right) H^{-1}_{\pm} \Omega \Omega (r_m \pm r_q) \\
     &= S_{\pm} (\tau)  \bar{r}^{\text{on}} -   \left( S_{\pm} (\tau)  - \mathds{1} \right) H^{-1}_{\pm} (r_m \pm r_q), 
\end{align*}
which is equivalent to Eq.~(\ref{eq:time_diagonal_general}). This results holds also for the on diagonal case of only operator valued force with $H_\pm \to H_m$. 

Let us focus on the case of only opearator valued foces, where the ODEs are integrable also for the off-diagonal case. It will be useful for generalisations (see Sec.~\ref{sec:generalization}), to compute the solution as a function of $r_j$. As stated in the main text, all the covariance matrices evolves under a single symplectic transformation $ S_m (\tau) = \exp(\Omega H_m \tau)$ such that
\begin{equation}
\label{eq:cov_t_force}
   \sigma (\tau) := \sigma_{jk}(\tau)  = S_m (\tau) \sigma_0 S_m^{\rm T} (\tau).
\end{equation}
Then we need to solve 
\begin{align}
    \dot{\bar{r}}_{jk} =     \Omega  H_m  \bar{r}_{jk}  -   \frac{1}{2}\Omega (r_j + r_k) + \frac{i}{2}  \sigma (\tau) (r_j - r_k) ,
\end{align}
obtaining
\begin{align}
\label{eq:app_first_moment_ODE}
    \bar{r}_{jk} (\tau) &= S_m (\tau) r_0  -   S_m(\tau) \frac12 \int_0^{\tau} \text{d} t  S_m^{-1}(t) \left( \Omega (r_j + r_k) - i  \sigma (t) (r_j - r_k)  \right) \nonumber  \\ 
     &= S_m (\tau) r_0 - \frac{1}{2} \left[ (\tilde{r}_j(\tau) - \tilde{r}_j) + (\tilde{r}_k(\tau) - \tilde{r}_k) \right] - \frac{i}{2}  S_m(\tau) \sigma_0 \int_0^{\tau} \text{d} t  S_m^{\rm T}(t)  (r_j - r_k) \nonumber  \\
     &= S_m (\tau) r_0 - \frac{1}{2} \left[ (\tilde{r}_j(\tau) - \tilde{r}_j) + (\tilde{r}_k(\tau) - \tilde{r}_k) \right] + \frac{i}{2}  S_m(\tau) \sigma_0 \left(S_m^{\rm T}(\tau) - \mathds{1} \right)  \Omega  H^{-1} (r_j - r_k) \nonumber \\
      %&= S_m (\tau) r_0- \frac{1}{2} \left[ (\tilde{r}_j(\tau) - \tilde{r}_j) + (\tilde{r}_k(\tau) - \tilde{r}_k) \right] - \frac{i}{2}  S_m(\tau) \sigma_0 S_m^{\rm T}(\tau) \left( \mathds{1} - (S_H^{\rm T})^{-1}(\tau)\right)  \Omega  H^{-1}_m (r_j - r_k) \nonumber \\
     &= S_m (\tau) r_0 - \frac{1}{2} \left[ (\tilde{r}_j(\tau) - \tilde{r}_j) + (\tilde{r}_k(\tau) - \tilde{r}_k )\right] - \frac{i}{2}  S_m(\tau) \sigma_0 S_m^{\rm T}(\tau)   \Omega \left( S_m(\tau)-  \mathds{1}   \right)  H^{-1}_m (r_j - r_k)  \nonumber\\
     &= r_0(\tau) - \frac{1}{2}  \left( \tilde{r}_j(\tau) - \tilde{r}_j + \tilde{r}_k(\tau) - \tilde{r}_k  \right) -  \frac{i}{2} \sigma (\tau)   \Omega  \left[ (\tilde{r}_j (\tau) - \tilde{r}_j)  - (\tilde{r}_k (\tau) - \tilde{r}_k ) \right] 
\end{align}
where the second equality follows from the time evolution of the covariance matrix Eq.~(\ref{eq:cov_t_force}) and the fourth from Eq.~(\ref{eq:identity_usefull}). From this, it is possible to derive Eq.(\ref{eq:r_on_force_sol}). We are left to integrate  
\begin{align*}
    \dot{\bar{r}}_{jk}^{(0)} &= - \dot{\mathcal{C}}_{jk}  + i \dot{\phi}_{jk} =    i  (r_j - r_k)^{\rm T}   \bar{r}_{jk}  - \frac{i}{2} (H_j^0-H_k^0) 
\end{align*}
By substituting Eq.~(\ref{eq:app_first_moment_ODE}) and splitting the real and immaginaty parts, one finds
\begin{align}
\label{eq:ode_C_appendix}
    \dot{\mathcal{C}}_{jk} &=  -  \frac{1}{2}  (r_j - r_k)^{\rm T}   \sigma (\tau)   \Omega  \left( (\tilde{r}_j (\tau) - \tilde{r}_j)  - (\tilde{r}_k (\tau) - \tilde{r}_k ) \right) =   \frac{1}{2}  (r_j - r_k)^{\rm T}  S_m(\tau) \sigma_0 \left(S_m^{\rm T}(\tau) - \mathds{1} \right)  \Omega  H^{-1} (r_j - r_k)  \\
    \label{eq:ode_phi_appendix}
    \dot{\phi}_{jk} &=  (r_j  - r_k )^{\rm T}\left[  r_0(\tau) - \frac{1}{2} \left( \tilde{r}_j(\tau) - \tilde{r}_j + \tilde{r}_k(\tau) - \tilde{r}_k \right)  \right] - \frac{1}{2} (H_j^0-H_k^0) 
\end{align}

In order to integrate the former, one may note that, for a general matrix $A$,  $(r_j - r_k)^{\rm T} A (r_j - r_k)  = \frac{1}{2} (r_j - r_k)^{\rm T} (A+A^{\rm T}) (r_j - r_k)$. Then, define the matrices $M_1(t) = (H^{-1} )^{\rm T}  \Omega^{\rm T} \left(S_m(t) - \mathds{1} \right)  $ and $M_2(t) = \sigma_0 \left(S_m^{\rm T}(t) - \mathds{1} \right)  \Omega H^{-1}$, such that $\dot{M}_1(t)= S_m(t)$ and $ \dot{M}_2(t) =  \sigma_0 S_m^{\rm T} (t)$. Note that $(\dot{M}_1(t) M_2(t))^{\rm T} = M_1(t) \dot{M}_2(t)$, such that 
\begin{align}
    \mathcal{C}_{jk} &=  \frac{1}{2} (r_j - r_k)^{\rm T} \left[ \int_0^{\tau} \text{d} t \dot{M}_1(t) M_2(t)  \right] (r_j - r_k)^{\rm T} =  \frac{1}{4} (r_j - r_k)^{\rm T} \left[ \int_0^{\tau} \text{d} t  \dot{M}_1(t) M_2(t) + (\dot{M}_1(t) M_2(t))^{\rm T}  \right] (r_j - r_k)^{\rm T} \nonumber \\
    &=   \frac{1}{4} (r_j - r_k)^{\rm T} \left[ \int_0^{\tau} \text{d} t  \dot{M}_1(t) M_2(t) + M_1(t) \dot{M}_2(t) \right] (r_j - r_k)^{\rm T}  =   \frac{1}{4} (r_j - r_k)^{\rm T} ( M_1(\tau) M_2(\tau) - M_1(0) M_2(0))  (r_j - r_k)^{\rm T}\nonumber \\
    &=  \frac{1}{4} (r_j - r_k)^{\rm T} H^{-1}  \Omega^{\rm T} \left(S_m(\tau) - \mathds{1} \right)  \sigma_0 \left(S_m^{\rm T}(\tau) - \mathds{1} \right)  \Omega H^{-1} (r_j - r_k) \label{eq:integral_app_contr} \\
    &=  \frac{1}{4} (r_j - r_k)^{\rm T} H^{-1}  \left(S^{\rm T}_m(\tau) - \mathds{1} \right)  \Omega^{\rm T} S_m(\tau) \sigma_0 S^{\rm T}_m(\tau)\Omega  \left(S_m^{\rm T}(\tau) - \mathds{1} \right)  H^{-1} (r_j - r_k)\nonumber \\
    &=  \frac{1}{4} \left[ (\tilde{r}_j (\tau) - \tilde{r}_j)  - (\tilde{r}_k (\tau) - \tilde{r}_k ) \right]^{\rm T} \Omega^{\rm T} \sigma (\tau)   \Omega  \left[ (\tilde{r}_j (\tau) - \tilde{r}_j)  - (\tilde{r}_k (\tau) - \tilde{r}_k ) \right] 
\end{align}
where we used the fact that $M_1(0) M_2(0) = 0$ and Eq.~(\ref{eq:identity_usefull}). Finally, we integrate the phases
\begin{align}
\label{eq:phase_appendix_unitary}
    \phi_{jk} &=  (r_j - r_k)^{\rm T} \left[ \int_0^{\tau} \text{d} t S_m (t) \right] \left( \tilde{r}_0  - \frac{1}{2} (\tilde{r}_j + \tilde{r}_k) \right) + \frac{\tau}{2} \left( (\tilde{r}_j - \tilde{r}_k)^{\rm T} H_m (\tilde{r}_j + \tilde{r}_k) -  (H_j^0-H_k^0) \right) \nonumber \\ 
    &=  - (r_j - r_k)^{\rm T} H^{-1} \Omega \left( S_m (\tau)  - \mathds{1} \right) \left( \tilde{r}_0  - \frac{1}{2} (\tilde{r}_j + \tilde{r}_k) \right) + \frac{\tau}{2} \left( (\tilde{r}_j - \tilde{r}_k)^{\rm T} H_m (\tilde{r}_j + \tilde{r}_k) -  (H_j^0-H_k^0) \right) \\ 
    & =  - (\tilde{r}_j - \tilde{r}_k)^{\rm T} \Omega \left[ ( \bar{r}_0  (\tau) - \bar{r}_0 ) - \frac{1}{2} \left( \tilde{r}_j (\tau) - \tilde{r}_j + \tilde{r}_k (\tau) -  \tilde{r}_k  \right) \right] + \frac{\tau}{2} \left[ (\tilde{r}_j - \tilde{r}_k)^{\rm T} H_m (\tilde{r}_j + \tilde{r}_k) -  (H_j^0-H_k^0) \right] \nonumber
\end{align}
This concludes the derivation in Table \ref{tab:force}. 


%In the on diagonal terms, they are trivially zero. Recall that $r_\pm = r_m \pm r_q$. In the off diagonal term, $\tilde{r}_{\pm} + \tilde{r}_{\mp} =  2 H_m^{-1} r_m = 2 \tilde{r}_m$ and $\tilde{r}_{\pm} - \tilde{r}_{\mp} = \pm 2 H_m^{-1} r_q = \pm 2 \tilde{r}_q$. The same holds for their time evolution. Thus, 
%\begin{align*}
%     \dot{\bar{r}}^{\text{off}}_{\pm 1, \mp 1} &= r_0(\tau) -  \left( \tilde{r}_m(\tau) - \tilde{r}_m \right) + i \sigma (\tau)   \Omega  \left( \tilde{r}_q (\tau) - \tilde{r}_q \right)  \\ 
%     \mathcal{C}^{\text{off}}_{\pm 1, \mp 1} &=   \left(\tilde{r}_q (\tau) - \tilde{r}_q \right)^{\rm T} \Omega^{\rm T} \sigma (\tau)   \Omega  \left(\tilde{r}_q (\tau) - \tilde{r}_q \right) \\
%      \phi^{\text{off}}_{\pm 1, \mp 1}  &=  \pm \tau ( H_q^0 - 2 r_q H_m r_m ) \mp 2 r_q^{\rm T} \Omega \left[ \tilde{r}_0 (\tau) - \tilde{r}_0 - ( \tilde{r}_m (\tau) - \tilde{r}_m ) \right]  
%\end{align*}
%as claimed in the main text \textcolor{red}{it is so wrong}.




\section{Maps of Additional Terms for Open Dynamcis \label{app:noise}}

First of all, let us find out how the ODEs change with the addition of the noise terms of Eq.~(\ref{eq:linbladian}) in the $\hat{r}$ basis. One finds that
\begin{align*}
    \mathcal{L}_B (\hat{r}) &:=  B^{mn} \left(  \hat{r}^m \hat{\varrho} \hat{r}^n  - \frac{1}{2} \{ \hat{\varrho}, \hat{r}^m  \hat{r}^n \} \right) =     \left(\frac{1}{2}  (\Omega^{\rm T} D \Omega)^{mn} -  i  {\rm e}^{mn} \right) \left(  \hat{r}^m \hat{\varrho} \hat{r}^n  - \frac{1}{2} \{ \hat{\varrho}, \hat{r}^m  \hat{r}^n \} \right) 
\end{align*}
where we used the fact that $B$ is hermitian, and we divided the symmetric real part ($\Omega^{\rm T} D \Omega$, where $D$ is symmetric) and antisymmetric complex part ($i  E$, where $E$ is \textit{skey}-symmetric). From Eq.~(\ref{eq:sim_phase_master}) and Eq.~(\ref{eq:antsim_phase_master}), one finds that  
\begin{equation}
    \mathcal{L}_{\hat{r}}  (\hat{\varrho}) \longleftrightarrow - 
    \left[ \frac{1}{4} (\Omega^{\rm T} D \Omega)^{mn}  \Omega^{n o} \Omega^{m p} \tilde{r}^o \tilde{r}^p   -   {\rm E}^{mn} \Omega^{n o} \tilde{r}^o  \partial_{\tilde{r}^m}  \right] \chi (\tilde{r}) =  - 
    \left[ \frac{1}{4} D^{op} \tilde{r}^o \tilde{r}^p  -  {\rm E}^{mn} \Omega^{n o} \tilde{r}^o  \partial_{\tilde{r}^m}  \right] \chi (\tilde{r}) 
\end{equation}
which represents the additional noise terms of the PDE of the time evolution of the branched charactersitic functions (Eq.~\ref{eq:PDE_chara_general}). By taking the ansatz for $\chi_{jk}$ (Eq.~\ref{eq:chara_jk}), one finds 
\begin{align}
    \mathcal{L}_{\hat{r}} (\hat{\varrho}) \longleftrightarrow &-  \left[ \frac{1}{4}  D^{mn} \tilde{r}^m \tilde{r}^n  + {\rm e}^{mn} \Omega^{no}  \left( \frac{1}{2} \sigma_{jk}^{mp} \tilde{r}^{p} -  i r^m_{jk}  \right) \tilde{r}^o  \right] \chi_{jk} (\tilde{r}) \nonumber \\
    &=  \left[ - \frac{1}{4} \left( D^{po} + 2 \sigma_{jk}^{pm} {\rm e}^{mn}  \Omega^{no}     \right) \tilde{r}^p \tilde{r}^o  +  i   {\rm e}^{mn} \Omega^{no}  r^m_{jk}   \tilde{r}^o  \right] \chi_{jk} (\tilde{r}) \nonumber 
\end{align}
By symmetrizing the first term and noticing that $ (\sigma E  \Omega )^{\rm T} = -  \Omega {\rm e}^{\rm T} \sigma $, the correction to the ODE for $\sigma_{jk}$ and $\bar{r}_{jk}$ are
\begin{equation}
\dot{\Sigma}_{jk} = \left( \dot{\sigma}_{jk} \right) + \Omega E  \sigma_{jk}  -  \sigma_{jk} E^{\rm T} \Omega + D \hspace{2cm} \dot{\bar{r}}_{jk} = \left( \dot{\bar{r}}_{jk} \right) + \Omega E \bar{r}_{jk} + \Omega d
\end{equation}
where the inclusion of the driving term ($d$) is easily concluded by noticing that it is just an additional force. We note that noise and driving of the modes do not affect the ODE for $\bar{r}_{jk}^{(0)}$ as there are no constant terms in the PDE of $\chi_{jk}$. 
The addition of these terms, with the inclusion of the DV dephasing, leads to the ODEs of Table \ref{tab:open}. 



\section{Integration of Linear Operator Valued Interactions: Open Dynamics \label{app:only_force_open}}

In the case of operator-valued forces and open dynamics, it is possible to write the general solution of the ODEs (Table~\ref{tab:open}) in integral form (Table~\ref{tab:open}), and by restricting to only diffusive noise ($E=0$), one finds a more compact form (Table \ref{tab:open_app}). In the case of only linear coupling, $H_j = H_m \forall j$, such that it is possible to define a transformation $S_A(\tau) = \exp(  \Omega A \tau)$ where $A =H_m + E$ is no longer symmetric and $S_A$ is not symplectic. We note that $S_A^{-1} (\tau) = \exp ( - \Omega A \tau) $ and $S_A^{\rm T} (\tau) = \exp (- A^{\rm T} \Omega \tau) $, and 
\begin{equation}
    \int_0^\tau \text{d}t S_A(t)  =\int_0^\tau \text{d}t {\rm e}^{ \Omega A \tau}  =  ( \Omega A)^{-1} \left(S_H(\tau) - \mathds{1} \right) = - A^{-1} \Omega  \left(S_H(\tau) - \mathds{1} \right)
\end{equation}
which implies 
\begin{equation}
    \int_0^\tau S_A^{-1}(t) dt =    \left(S_A^{-1}(\tau) - \mathds{1} \right) A^{-1} \Omega \hspace{1cm} \int_0^\tau S_A^{\rm T}(t) dt =     \left(S_A^{\rm T}(\tau) - \mathds{1} \right) \Omega  (A^{\rm T})^{-1}
\end{equation}

\begin{table*}
    \centering
    \renewcommand{\arraystretch}{1.7}
    \begin{tabular}{ |p{4.0cm}||P{7cm}|P{7cm}|}
       \hline
        Phase Spaces parameters & \multicolumn{2}{c|}{Operator Valued Linear Interactions (Symmetric Case)} \\ 
        \hline
        Covariance Matrices   & \multicolumn{2}{c|}{\ensuremath{
   \sigma (\tau) := \sigma_{jk}(\tau) = 
    S_m(\tau) \sigma_0 S_m^{\rm T}(\tau) + \int_{0}^{\tau} \text{d}t S_m(\tau - t) D S^{\rm T}_m(\tau - t) }}\\
    \hline
       Vectors or First Moments & \multicolumn{2}{c|}{
       \ensuremath{\bar{r}_{jk} (\tau) = r_0(\tau)  + d(\tau) - d - \frac{1}{2}  \left( \tilde{r}_j(\tau) - \tilde{r}_j + \tilde{r}_k(\tau) - \tilde{r}_k  \right) - \frac{i}{2}  \sigma (\tau)  \Omega  \left[ (\tilde{r}_j (\tau) - \tilde{r}_j)  - (\tilde{r}_k (\tau) - \tilde{r}_k ) \right] }} \\
           & \multicolumn{2}{c|}{
       \ensuremath{ \hspace{2cm}  - \frac{i}{2}  \int_0^{\tau} \text{d}t  D (\tau-t) \Omega  \left[ (\tilde{r}_j(\tau - t) - \tilde{r}_j(\tau) ) -  (\tilde{r}_k(\tau - t) - \tilde{r}_k(\tau) ) \right]  }} \\
       \hline
     QRDM Contrasts  & \multicolumn{2}{c|}{ \hspace{-0.1cm} \ensuremath{  \mathcal{C}_{jk} (\tau)  =    \frac{1}{4} \left[ (\tilde{r}_j (\tau) - \tilde{r}_j)  - (\tilde{r}_k (\tau) - \tilde{r}_k ) \right]^{\rm T} \Omega^{\rm T} \sigma (\tau)   \Omega  \left[ (\tilde{r}_j (\tau) - \tilde{r}_j)  - (\tilde{r}_k (\tau) - \tilde{r}_k ) \right] + \tau \left[\frac{1}{2} (\Gamma_z^{jj} + \Gamma_z^{kk} ) - \Gamma_z^{jk} \right] }} \\
 & \multicolumn{2}{c|}{\ensuremath{ \hspace{1cm} + \frac{1}{4}     \int_0^\tau \text{d} t   \Big\{ \left[ (\tilde{r}_j (\tau - t) - \tilde{r}_j(\tau))  - (\tilde{r}_k (\tau - t) - \tilde{r}_k(\tau))  \right]^{\rm T}  \Omega^{\rm T}  D(\tau - t) \Omega  }} \\
 & \multicolumn{2}{c|}{\ensuremath{ \hspace{2.5cm} 
 \left[ (\tilde{r}_j (\tau - t)  + \tilde{r}_j(\tau) - 2 r_j)  - (\tilde{r}_k (\tau - t) + \tilde{r}_k(\tau) - 2 r_k)  \right]  \Big\} }} \\
\hline 
    \end{tabular}
    \caption{Open evolution of a GBCS phase space quantity under operator-valued linear Hamiltonian (Eq. \ref{eq:Hamiltonian_force_only}) and Gaussian Diffusive noise ($D$ and $E=0$) of an initial Gaussian state with initial first and second moments $r_0$ and $\sigma_0$, where $\tilde{r}_j = H^{-1}_m r_j$, $S_m (\tau) = \exp (  \Omega H_m \tau)$, and $\tilde{r}_a (\tau) = S_m (\tau) \tilde{r}_a$, with $a \in 0,j,k$. The phases are the same of Table~\ref{tab:force}.}
    \label{tab:open_app}
\end{table*}

Under such definitions, the ODEs for $\sigma_{jk}$ (Table \ref{tab:open}) and their solutions reads,
\begin{align}
    \dot{\sigma}_{jk} &=   \sigma_{jk}  A \Omega  -   \Omega  A^{\rm T}   \sigma_{jk}  + D   \implies \sigma_{jk} = \sigma(\tau) = S_A(\tau) \sigma_0 S_A^{\rm T}(\tau) + \int_{0}^{\tau} \text{d}t S_A(\tau - t) D S^{\rm T}_A(\tau - t)
\end{align}
which is the usual solution of Lyapunov equations, to which all second moments evoleve. Thus 
\begin{align}
     \dot{\bar{r}}_{jk} &=    \Omega A  \bar{r}_{jk}  + \frac{1}{2} \Omega (2 d - r_j - r_k)  +  \frac{i}{2}    \sigma(\tau) (r_j - r_k)   
\end{align}
which can be integrated 
\begin{align*}
    \bar{r}_{jk} (\tau) &= S_A (\tau) r_0  + \frac{1}{2}   S_A(\tau)  \int_0^{\tau} \text{d} t  S_A^{-1}(t) \left( \Omega (2d - r_j - r_k) + i  \sigma (t) (r_j - r_k)  \right)  \\ 
    &= S_A (\tau) r_0  + \frac{1}{2}   S_A(\tau)  \left(S_A^{-1}(\tau) - \mathds{1} \right) A^{-1} \Omega \Omega (2d - r_j - r_k) \\
    &\hspace{1cm} +  \frac{i}{2} S_A(\tau)  \sigma_0 \int_0^{\tau} \text{d} t  S_A^{T}(t)  (r_j - r_k)   + \frac{i}{2} S_A(\tau) \int_0^{\tau} \text{d}t \int_{0}^{t}  \text{d} t'  S_A^{-1}(t) S_A(t-t')D S^{\rm T}_A(t - t')   (r_j - r_k)    \\ 
    &= S_A (\tau) r_0   + \frac{1}{2}     \left(S_A(\tau) - \mathds{1} \right) A^{-1}  (2d - r_j - r_k) + 
    \frac{i}{2} S_A (\tau) \sigma_0 \left( S_A^{\rm T} (\tau) - \mathds{1}  \right) \Omega (A^{\rm T})^{-1}  (r_j - r_k)   \\
    &\hspace{1cm} + \frac{i}{2} \int_0^{\tau} \text{d}t'  S_A(\tau- t') D \int_{t'}^{\tau}  \text{d} t S^{\rm T}_A(t- t')(r_j - r_k)  \\
      &=S_A (\tau) r_0  + \frac{1}{2}     \left(S_A(\tau) - \mathds{1} \right) A^{-1}  ( 2 d - r_j - r_k) + 
    \frac{i}{2} S_A (\tau) \sigma_0 \left( S_A^{\rm T} (\tau) - \mathds{1}  \right) \Omega (A^{\rm T})^{-1}  (r_j - r_k)   \\
    &\hspace{1cm} + \frac{i}{2}  \int_0^{\tau} \text{d}t'  S_A(\tau- t') D \left( S^{\rm T}_A(\tau - t') - \mathds{1} \right) \Omega (A^{\rm T})^{-1} (r_j - r_k) 
\end{align*}
The final parameters for the QRDM can be integrated
\begin{align*}
    \dot{\phi}_{jk} &= - (r_j  - r_k )^{\rm T} \left[ S_A (\tau) r_0  +     \left(S_A(\tau) - \mathds{1} \right) A^{-1} \left[ d - \frac{1}{2}  (r_j + r_k) \right] \right] + \frac{1}{2} (H_j^0-H_k^0) 
\end{align*}
which implies that
\begin{align}
    \phi_{jk} (\tau) =  - (r_j - r_k)^{\rm T} A^{-1} \Omega \left( S_A (\tau)  - \mathds{1} \right) \left[ r_0  + A^{-1} \left( d  - \frac{1}{2} (r_j + r_k) \right) 
    \right] + \tau\left( (r_j - r_k)^{\rm T} A^{-1} \left(d - \frac{1}{2} (r_j + r_k) \right) -  \frac{1}{2} (j-k) H_q^0 \right) \nonumber
\end{align}
from similarity of Eq. \ref{eq:ode_C_appendix} (i.e. the replacement $H_m \to A$ and addition of $d$). For the contrasts one needs to integrate
\begin{align*}
    \dot{\mathcal{C}}_{jk} &=   \frac{1}{2}  (r_j - r_k)^{\rm T}  S_A(\tau) \sigma_0 \left(S_A^{\rm T}(\tau) - \mathds{1} \right)  \Omega  (A^{\rm T})^{-1} (r_j - r_k)  \\
    &\hspace{1cm} +  \frac{1}{2}  (r_j - r_k)^{\rm T} \int_0^{\tau} \text{d}t'  S_A(\tau- t') D \left( S^{\rm T}_A(\tau - t') - \mathds{1} \right) \Omega (A^{\rm T})^{-1} (r_j - r_k)
\end{align*}
such that
\begin{align*}
     \mathcal{C}_{jk} (\tau)  &=  \frac{1}{4} (r_j - r_k)^{\rm T} A^{-1}  \Omega^{\rm T} \left( S_A(\tau) - \mathds{1} \right) \sigma_0 \left(S^{\rm T}_A(\tau) - \mathds{1} \right) \Omega  (A^{\rm T})^{-1} (r_j - r_k)   \nonumber \\
     &\hspace{1cm} + \frac{1}{2}  (r_j - r_k)^{\rm T} \int_0^\tau \text{d} t \int_0^{t} \text{d}t'  S_A(t- t') D \left( S^{\rm T}_A(t - t') - \mathds{1} \right) \Omega (A^{\rm T})^{-1} (r_j - r_k) 
\end{align*}
where the first line comes from Eq. \ref{eq:integral_app_contr}. For the second term one can approach the integration similarly to Appendix \ref{app:only_force}, by noticing that $(r_j - r_k)^{\rm T} M (r_j - r_k)  = \frac{1}{2} (r_j - r_k)^{\rm T} (M+M^{\rm T}) (r_j - r_k)$, for any matrix $M$.  Then, define the matrices $N_1(t, t') = A^{-1}  \Omega^{\rm T} \left(S_A(t - t') - \mathds{1} \right)  $ and $N_2(t, t') =  D \left(S_A^{\rm T}(t - t') - \mathds{1} \right)  \Omega (A^{\rm T})^{-1} $, such that 
$\frac{d N_1}{d t} =   S_A(t - t')$ and $\frac{d N_2}{d t}=  D S_m^{\rm T} (t - t')$. Note that 
\begin{equation*}
    \left(\frac{d N_1}{d t}  N_2 \right)^{\rm T}  =  - \left( S_A(t - t')D \left(S_A^{\rm T}(t - t') - \mathds{1} \right)  \Omega^{\rm T} (A^{\rm T})^{-1} \right)^{\rm T}   =  N_1 \frac{d N_2}{d t}
\end{equation*}
and so the second line reads 
\begin{align*}
    &=  \frac{1}{2} (r_j - r_k)^{\rm T} \int_0^\tau \text{d} t' \int_{t'}^{\tau} \text{d}t  S_A(t- t') D \left( S^{\rm T}_A(t - t') - \mathds{1} \right) \Omega (A^{\rm T})^{-1} (r_j - r_k) \\
     &=   \frac{1}{2} (r_j - r_k)^{\rm T} \int_0^\tau \text{d} t' \int_{t'}^{\tau} \text{d}t   \left[ \frac{d N_1}{d t}  N_2\right]  (r_j - r_k) =  \frac{1}{4} (r_j - r_k)^{\rm T} \int_0^\tau \text{d} t' \int_{t'}^{\tau} \text{d}t  \left[ \frac{d N_1}{d t}  N_2 +  \left(\frac{d N_1}{d t}  N_2 \right)^{\rm T} \right] (r_j - r_k)    \\ 
      &=  \frac{1}{4} (r_j - r_k)^{\rm T} \left[ \int_0^\tau \text{d} t'    N_1(\tau, t') N_2(\tau, t') - N_1(t',t') N_2(t',t') \right] (r_j - r_k)  \\
      &=  \frac{1}{4} (r_j - r_k)^{\rm T}  A^{-1}  \Omega^{\rm T} \left[  \int_0^\tau \text{d} t'    \left(S_A(\tau - t') - \mathds{1} \right) D \left(S_A^{\rm T}(\tau - t') - \mathds{1} \right)   \right]  \Omega (A^{\rm T})^{-1} (r_j - r_k) 
\end{align*}
where we integrated by parts. The entire expression reads
\begin{align*}
     \mathcal{C}_{jk} (\tau)  &=  \frac{1}{4} (r_j - r_k)^{\rm T}  A^{-1} \Omega^{\rm T} \left( S_A(\tau) - \mathds{1} \right) \sigma_0 \left(S^{\rm T}_A(\tau) - \mathds{1} \right) \Omega  (A^{\rm T})^{-1}  (r_j - r_k)   \\
     &\hspace{1cm} + \frac{1}{4} (r_j - r_k)^{\rm T}  A^{-1}   \Omega^{\rm T} \left[  \int_0^\tau \text{d} t'    \left(S_A(\tau - t') - \mathds{1} \right) D \left(S_A^{\rm T}(\tau - t') - \mathds{1} \right)   \right]  \Omega (A^{\rm T})^{-1}  (r_j - r_k) 
\end{align*}

In the case of $E=0$, then  $A^{\rm T} = A = H_m$  and $S_A$ is symplectic. In this case, it is possible to use the identity of Eq.~(\ref{eq:identity_usefull}). By defining $\tilde{r}_j = H^{-1}_m r_j$ and $\tilde{d} = H^{-1}_m d$, and their symplectic evolution $\tilde{r}_j (\tau)= S_A(\tau)\tilde{r}_j$ and $\tilde{d} (\tau)= S_A(\tau)\tilde{d}$, the following further simplification is found
\begin{align*}
\bar{r}_{jk} (\tau) &= r_0(\tau) + d(\tau) - d - \frac{1}{2}  \left( \tilde{r}_j(\tau) - \tilde{r}_j + \tilde{r}_k(\tau) - \tilde{r}_k  \right) - 
    \frac{i}{2} S_A (\tau) \sigma_0 S_A^{\rm T} (\tau) \Omega \left( S_A (\tau) - \mathds{1}  \right) (\tilde{r}_j - \tilde{r}_k)   \\
    &\hspace{1cm}  - \frac{i}{2}  \int_0^{\tau} \text{d}t'  S_A(\tau- t') D S^{\rm T}_A(\tau - t') \Omega \left( S_A(\tau - t') - \mathds{1} \right) (\tilde{r}_j - \tilde{r}_k) \\
     &= r_0(\tau)  + d(\tau) - d - \frac{1}{2}  \left( \tilde{r}_j(\tau) - \tilde{r}_j + \tilde{r}_k(\tau) - \tilde{r}_k  \right) -
    \frac{i}{2} S_A (\tau) \sigma_0 S_A^{\rm T} (\tau) \Omega  \left[ (\tilde{r}_j(\tau) - \tilde{r}_j) - (\tilde{r}_k(\tau) - \tilde{r}_k)  \right]    \\
    &\hspace{1cm}  - \frac{i}{2}  \int_0^{\tau} \text{d}t'  S_A(\tau- t') D S^{\rm T}_A(\tau - t') \Omega \left( S_A(\tau - t') - S_A(\tau) + S_A(\tau) - \mathds{1} \right)  (\tilde{r}_j - \tilde{r}_k) 
 \end{align*}
\begin{align*}
\bar{r}_{jk} (\tau)    &= r_0(\tau)  + d(\tau) - d - \frac{1}{2}  \left( \tilde{r}_j(\tau) - \tilde{r}_j + \tilde{r}_k(\tau) - \tilde{r}_k  \right) \\
    &\hspace{1cm}- 
    \frac{i}{2} \left( S_A (\tau) \sigma_0 S_A^{\rm T} (\tau) + \int_0^{\tau} \text{d}t'  S_A(\tau- t') D S^{\rm T}_A(\tau - t') \right)\Omega  \left[ (\tilde{r}_j(\tau) - \tilde{r}_j) - (\tilde{r}_k(\tau) - \tilde{r}_k)  \right]    \\
    &\hspace{1cm}  - \frac{i}{2}  \int_0^{\tau} \text{d}t'  S_A(\tau- t') D S^{\rm T}_A(\tau - t') \Omega \left( S_A(\tau - t') - S_A(\tau) \right)  (\tilde{r}_j - \tilde{r}_k) \\
      &=  r_0(\tau)  + d(\tau) - d - \frac{1}{2}  \left( \tilde{r}_j(\tau) - \tilde{r}_j + \tilde{r}_k(\tau) - \tilde{r}_k  \right) - \frac{i}{2}  \sigma (\tau)  \Omega  \left[ (\tilde{r}_j (\tau) - \tilde{r}_j)  - (\tilde{r}_k (\tau) - \tilde{r}_k ) \right] \\
      & \hspace{2cm}  - \frac{i}{2}  \int_0^{\tau} \text{d}t  D (\tau-t) \Omega  \left[ (\tilde{r}_j(\tau - t) - \tilde{r}_j(\tau) ) -  (\tilde{r}_k(\tau - t) - \tilde{r}_k(\tau) ) \right] 
\end{align*}
where $D(\tau) = S_A(\tau)D S_A^{\rm T}(\tau)$. Similarly, 
\begin{align*}
C_{jk} (\tau) &=   \frac{1}{4} (\tilde{r}_j - \tilde{r}_k)^{\rm T}   \left(S^{\rm T}_A(\tau) - \mathds{1} \right)  \Omega^{\rm T} S_A(\tau) \sigma_0 S^{\rm T}_A(\tau)\Omega  \left(S_A(\tau) - \mathds{1} \right)   (\tilde{r}_j - \tilde{r}_k)   \nonumber \\ 
& \hspace{1cm}+ \frac{1}{4} (\tilde{r}_j - \tilde{r}_k)^{\rm T}  \left[  \int_0^\tau \text{d} t'    \left(S^{\rm T}_A(\tau - t') - \mathds{1} \right) \Omega^{\rm T} S_A(\tau - t')   D S_A^{\rm T}(\tau - t') \Omega \left(S_A(\tau - t') - \mathds{1} \right)   \right]   (\tilde{r}_j - \tilde{r}_k) \\
&=   \frac{1}{4} (\tilde{r}_j - \tilde{r}_k)^{\rm T}   \left(S^{\rm T}_A(\tau) - \mathds{1} \right)  \Omega^{\rm T} S_A(\tau) \sigma_0 S^{\rm T}_A(\tau)\Omega  \left(S_A(\tau) - \mathds{1} \right)   (\tilde{r}_j - \tilde{r}_k)   \nonumber \\ 
& \hspace{0.4cm} + \frac{1}{4} (\tilde{r}_j - \tilde{r}_k)^{\rm T}   \left[  \int_0^\tau \text{d} t'    \left(S^{\rm T}_A(\tau - t') - S^{\rm T}_A(\tau) +  S^{\rm T}_A(\tau) - \mathds{1} \right) \Omega^{\rm T}  D(\tau - t') \Omega \left(S_A(\tau - t')  - S_A(\tau) +  S_A(\tau) - \mathds{1} \right)   \right]   (\tilde{r}_j - \tilde{r}_k) \\
&=   \frac{1}{4} (\tilde{r}_j - \tilde{r}_k)^{\rm T}   \left(S^{\rm T}_A(\tau) - \mathds{1} \right)  \Omega^{\rm T} \left[  S_A(\tau) \sigma_0 S^{\rm T}_A(\tau) + \int_0^\tau \text{d} t D(\tau - t)  \right] \Omega  \left(S_A(\tau) - \mathds{1} \right)   (\tilde{r}_j - \tilde{r}_k)   \nonumber \\ 
& \hspace{0.4cm} + \frac{1}{4} (\tilde{r}_j - \tilde{r}_k)^{\rm T}   \left[  \int_0^\tau \text{d} t'    \left(S^{\rm T}_A(\tau - t') - S^{\rm T}_A(\tau)  \right) \Omega^{\rm T}  D(\tau - t') \Omega \left(S_A(\tau - t')  - S_A(\tau)  \right)   \right]   (\tilde{r}_j - \tilde{r}_k) \\
& \hspace{0.4cm} + \frac{1}{2} (\tilde{r}_j - \tilde{r}_k)^{\rm T}   \left[  \int_0^\tau \text{d} t'    \left(S^{\rm T}_A(\tau - t') - S^{\rm T}_A(\tau)  \right) \Omega^{\rm T}  D(\tau - t') \Omega \left(S_A(\tau)  - \mathds{1} \right)   \right]   (\tilde{r}_j - \tilde{r}_k) \\
&=   \frac{1}{4} \left[ (\tilde{r}_j (\tau) - \tilde{r}_j)  - (\tilde{r}_k (\tau) - \tilde{r}_k ) \right]^{\rm T} \Omega^{\rm T} \sigma (\tau)   \Omega  \left[ (\tilde{r}_j (\tau) - \tilde{r}_j)  - (\tilde{r}_k (\tau) - \tilde{r}_k ) \right]  \nonumber \\ 
& \hspace{2cm} + \frac{1}{4}     \int_0^\tau \text{d} t   \Big\{ \left[ (\tilde{r}_j (\tau - t) - \tilde{r}_j(\tau))  - (\tilde{r}_k (\tau - t) - \tilde{r}_k(\tau))  \right]^{\rm T}  \Omega^{\rm T}  D(\tau - t) \Omega \\&\hspace{4cm}\left[ (\tilde{r}_j (\tau - t)  + \tilde{r}_j(\tau) - 2 r_j)  - (\tilde{r}_k (\tau - t) + \tilde{r}_k(\tau) - 2 r_k)  \right]  \Big\}
\end{align*}
The phase is the same as in Table \ref{tab:unitary}, thus, concluding all the derivations of Table \ref{tab:open_app}.

\section{Doubling Solutions of Riccati Equations \label{app:doubling}}


Consider the ansatz for the covariance matrices $\sigma_{jk} = \mathds{Y}_{jk} \mathds{X}_{jk}^{-1}$, where $\mathds{X}_{jk}, Y_{jk}$ are two time dependent $2n$-dimensional matrices which are solutions to some coupled linear ODEs: 
$\dot{X}_{jk} = A_{jk} \mathds{X}_{jk} + B_{jk} \mathds{Y}_{jk} $ and   $\dot{Y}_{jk} = C_{jk} \mathds{X}_{jk} + D_{jk} \mathds{Y}_{jk} $. The initial condition $\sigma (0) = \sigma_0$ can be found as $\mathds{Y}_{jk} = \sigma_0$ and $\mathds{X}_{jk} = \mathds{1}$. Thus, computing  
\begin{align}
    \dot{\sigma}_{jk} &= \dot{\mathds{Y}}_{jk} \mathds{X}_{jk}^{-1}- \mathds{Y}_{jk} \mathds{X}_{jk}^{-1} \dot{\mathds{X}}_{jk} \mathds{X}_{jk}^{-1} = \left(  C_{jk} \mathds{X}_{jk} + D_{jk} \mathds{Y}_{jk} \right) \mathds{X}_{jk}^{-1} - \mathds{Y}_{jk} \mathds{X}_{jk}^{-1} \left(  A_{jk} \mathds{X}_{jk} + B_{jk} \mathds{Y}_{jk} \right) \mathds{X}_{jk}^{-1} \\
    & = C_{jk} + D_{jk} (\mathds{Y}_{jk} \mathds{X}_{jk}^{-1}) - (\mathds{Y}_{jk} \mathds{X}_{jk}^{-1})A_{jk}  - (\mathds{Y}_{jk} \mathds{X}_{jk}^{-1} ) B_{jk} (\mathds{Y}_{jk} \mathds{X}_{jk}^{-1}) \\
    & = D_{jk} {\sigma}_{jk} -  {\sigma}_{jk} A_{jk} - {\sigma}_{jk}^{\rm T} B_{jk} {\sigma}_{jk} + C_{jk} 
\end{align}
Comparing to the ode of the Table~\ref{tab:open},  one can infer that 
\begin{align}
    \frac{\text{d}}{\text{d} \tau}\begin{pmatrix}
        \mathds{X}_{jk} \\
        \mathds{Y}_{jk}
    \end{pmatrix} &= \begin{pmatrix}
         A_{jk} & B_{jk}\\
        C_{jk} & D_{jk}
    \end{pmatrix}
    \begin{pmatrix}
        \mathds{X}_{jk} \\
        \mathds{Y}_{jk}
    \end{pmatrix} =  \frac{1}{2} \begin{pmatrix}
          ( H_j + H_k   + \textcolor{blue}{2 {\rm E}^{\rm T}} ) \Omega  &  i ( H_j - H_k) \\
          i \Omega ( H_j - H_k)  \Omega^{\rm T }  \textcolor{blue}{+ D} & \ \Omega  ( H_j + H_k   + \textcolor{blue}{2 E} )
    \end{pmatrix} \begin{pmatrix}
        \mathds{X}_{jk} \\
        \mathds{Y}_{jk}
    \end{pmatrix} = \mathds{A}_{jk}
    \begin{pmatrix}
        \mathds{X}_{jk} \\
        \mathds{Y}_{jk}
    \end{pmatrix} \nonumber
\end{align}
where we also included general Gaussian noise. One can find the solution of these linear ODEs in terms of exponential of $4n*4n$ matrix $S_{\mathds{A}_{jk}}(\tau) = \exp(\tau \mathds{A}_{jk})$ and  $\begin{pmatrix}
    \mathds{X}_{jk} (\tau)\\
    \mathds{Y}_{jk}(\tau)
\end{pmatrix}  = S_{\mathds{A}_{jk}}(\tau) \begin{pmatrix}
    \mathds{1} \\
    \sigma_0
\end{pmatrix} $. 
One can now turn to solve the first moment ODE
\begin{equation}
    \dot{\bar{r}}_{jk} =\frac{1}{2}\left[ \left(  \Omega ( H_j + H_k \textcolor{blue}{+2 E} )  - i   \sigma_{jk} (H_j - H_k) \right)  \bar{r}_{jk} 
     + \Omega (\textcolor{blue}{2d} - r_j - r_k) + i  \sigma_{jk} ( r_j - r_k) \right] 
\end{equation}
via the propagator 
\begin{equation}
\label{eq:propagator_r_off}
    \mathds{P}_{jk}(\tau) = \mathcal{T} \left[ \exp \left(\frac{1}{2} \int_{0}^{\tau} \text{d} t \;  \Omega ( H_j + H_k +2 E)  - i   \sigma_{jk}(t) (H_j - H_k) \right)  \right] 
\end{equation}
where $ \mathcal{T}$ is the time ordering, such that
\begin{equation}
    \bar{r}_{jk} = \mathds{P}_{jk}(\tau) r_0 + \frac{1}{2} \mathds{P}_{jk}(\tau) \int_{0}^{\tau} \text{d} t \mathds{P}^{-1}_{jk}(t)  \left[ i  \sigma_{jk} (t) (r_j - r_k) + \Omega (\textcolor{blue}{2d } - r_j - r_k) \right]
\end{equation}
from which one can compute the time evolution of the QRDM exponent. We should note that, often, the exponent of $\mathds{P}_{jk}(\tau) $ commutes at different times, which makes the time ordering trivial. The QRDM elements $r^{(0)}_{jk}$ can then be computed by direct integration of the final trivial ODE.


\section{Protocols and Operations \label{app:operation}}

In this section, we derive the maps of GBCS under gaussian operations on the CV and common conditional operations on the qubit states . 

\subsection{Gaussian Operations and Deterministic CP Maps}

A Gaussian completely positive map $\Phi_g$ is uniquely determined by two matrices $X$ and $Y$, such that $Y+ i \Omega \geq i X \Omega X^{\rm T}$, to preserve the Heisenberg uncertainty principle. The action on a GBCS can be computed via the dual map $\Phi_g^*$, defined as $\text{Tr}_m \left[\Phi_{g} [\hat{A}]  \hat{B}\right] = \text{Tr}_m \left[\hat{A}  \Phi_{g}^* [\hat{B}]\right]$, on the Displacement operators basis, such that one can use the well known identity 
 \begin{equation}
     \Phi_g^* \left(\hat{D}_{\Omega r} \right) = {\rm e}^{-\frac{1}{4} r^{\rm T}  Y  r }\hat{D}_{\Omega X^{\rm T} r} \implies \Phi_g^* \left(\hat{D}_{- r} \right) = {\rm e}^{-\frac{1}{4} r^{\rm T} \Omega  Y \Omega^{\rm T}  r }\hat{D}_{- \Omega X^{\rm T} \Omega^{\rm T} r}
 \end{equation}
From definition, the branched characteristic functions of a GBCS it follwos the update on the characterstic function under the general Gaussian map 
\begin{align}
    \chi[\Phi_g(\hat{\varrho}_{jk})] (r) &= \text{Tr}[\Phi_g[\hat{\varrho}_{jk}] \hat{D}_{-r}] = \text{Tr}[\hat{\varrho}_{jk} \Phi^*_g[\hat{D}_{- r}]] = {\rm e}^{-\frac{1}{4} r^{\rm T} \Omega  Y \Omega^{\rm T}  r } \text{Tr}[\hat{\varrho}_{jk} \hat{D}_{- \Omega X^{\rm T} \Omega^{\rm T} r} ] =  {\rm e}^{-\frac{1}{4} r^{\rm T} \Omega  Y \Omega^{\rm T}  r } \chi_{jk} (\Omega X^{\rm T} \Omega^{\rm T} r) \nonumber \\
    &=   
  \exp \left( -\frac{1}{4}  r^{\rm T}  \Omega^{\rm T}  \left( X^{\rm T} \sigma_{jk}  X + Y \right) \Omega r  +  i  r^{\rm T} \Omega^{\rm T} X \bar{r}_{jk} + r^{(0)}_{jk}   \right) \;,
\end{align} 
which is a GBCS (Eq.~\ref{eq:chara_jk}) under the update rules for the phase space parameters given in Eq.~(\ref{eq:Gaussian_CP_update}).

Specific examples present in laboratories are given in the main text, and here we reports other examples~\cite{serafini_quantum_2017}. Specifically, common coherent operations are the phase shift and beam splitter, defined according to the transformations
\begin{equation}
    R_{\varphi} = \begin{pmatrix} \cos\varphi & \sin\varphi \\ -\sin\varphi & \cos\varphi \end{pmatrix} \hspace{1cm}
    R_{\theta} = \begin{pmatrix} \cos\theta & 0 & \sin\theta & 0 \\ 0 & \cos\theta & 0 & \sin\theta \\ -\sin\theta & 0 & \cos\theta & 0 \\ 0 & -\sin\theta & 0 & \cos\theta \end{pmatrix}
\end{equation}
and the single and two mode squeezing transformations are given by 
\begin{equation}
    S_{1m} = \begin{pmatrix} s & 0 \\ 0 & 1/s \end{pmatrix} \hspace{1cm}
    S_{2m} = \begin{pmatrix} \cosh r & 0 & \sinh r & 0 \\ 0 & \cosh r & 0 & -\sinh r \\ \sinh r & 0 & \cosh r & 0 \\ 0 & -\sinh r & 0 & \cosh r \end{pmatrix}
\end{equation}


\subsection{Conditional Rotations}

%Let us consider the general conditional rotation described by Eq.~(\ref{eq:ODE_conditional_rot}). For an initial coherent state, one fins that\begin{equation} X(\tau) = Y(\tau) = \ \begin{pmatrix} {\rm e}^{2 i \beta  \tau } & 0 \\  0 & {\rm e}^{2 i \beta  \tau } \\ \end{pmatrix} \implies \sigma^{\text{off}}(\tau)= \mathds{1} \end{equation}  Substituting into the equation for the first moments, and being diagonal, one finds the solution in terms of the exponential \begin{equation}   \bar{r}^{\text{off}}(\tau) = \exp \left(  \mp i \beta \int_{0}^{\tau} {\rm d} t \sigma^{\text{off}} (t) \right) r_0 = {\rm e}^{- i \beta \tau}\begin{pmatrix}      x_0  \\    p_0  \end{pmatrix} \end{equation}such that the integral for the QRDM exponent is \begin{equation}   \dot{r}^{(0)} = - i \beta  |r_0|^2 \int_{0}^{\tau} {\rm d} t   \left(1+{\rm e}^{-2 i \beta  \tau } \right) \implies r^{(0)}= \frac{1}{2}|r_0|^{2} \left(1-{\rm e}^{-2 i \beta  \tau } \right) - i \beta \tau \end{equation}


For an initial sqeueezed state, one finds the solution of the complex Riccati equation using the method outlined in Sec.~\ref{subsec:riccati}, with the two matrices (at $\tau = 1$)
\begin{equation}
    X = \left(
\begin{array}{cc}
 \cos (\bar{\theta}) \left(\cos \theta_\Delta +\frac{i}{s} \sin \theta_\Delta \right) & \sin (\bar{\theta}) (\cos \theta_\Delta +i s \sin \theta_\Delta ) \\
 -\frac{1}{s} \sin (\bar{\theta}) (s \cos \theta_\Delta +i \sin \theta_\Delta ) & \cos (\bar{\theta}) (\cos \theta_\Delta +i s \sin \theta_\Delta ) \\
\end{array}
\right)\nonumber
\end{equation}
\begin{equation}
   Y =  \left(
\begin{array}{cc}
 \frac{1}{s}\cos (\bar{\theta}) (\cos \theta_\Delta +i s \sin \theta_\Delta ) & \sin (\bar{\theta}) (s \cos \theta_\Delta +i \sin \theta_\Delta ) \\
 -\frac{1}{s}\sin (\bar{\theta}) (\cos \theta_\Delta +i s \sin \theta_\Delta ) & \cos (\bar{\theta}) (s \cos \theta_\Delta +i \sin \theta_\Delta ) \\
\end{array}
\right)\nonumber
\end{equation}
such that one finds 
\begin{equation}
    \sigma^{\text{off}} = \left(
\begin{array}{cc}
 \frac{\cos ^2(\bar{\theta}) (\cos \theta_\Delta +i s \sin \theta_\Delta )}{s \cos \theta_\Delta +i \sin \theta_\Delta }+\frac{\sin ^2(\bar{\theta}) (s \cos \theta_\Delta +i \sin \theta_\Delta )}{\cos \theta_\Delta +i s \sin \theta_\Delta } & \frac{\left(s^2-1\right) \sin (\bar{\theta}) \cos (\bar{\theta})}{(s \cos \theta_\Delta +i \sin \theta_\Delta ) (\cos \theta_\Delta +i s \sin \theta_\Delta )} \\
 \frac{\left(s^2-1\right) \sin (\bar{\theta}) \cos (\bar{\theta})}{(s \cos \theta_\Delta +i \sin \theta_\Delta ) (\cos \theta_\Delta +i s \sin \theta_\Delta )} & \frac{\cos ^2(\bar{\theta}) (s \cos \theta_\Delta +i \sin \theta_\Delta )}{\cos \theta_\Delta +i s \sin \theta_\Delta }+\frac{\sin ^2(\bar{\theta}) (\cos \theta_\Delta +i s \sin \theta_\Delta )}{s \cos \theta_\Delta +i \sin \theta_\Delta } \\
\end{array}
\right)
\end{equation}
Which can be rewritten to the form of Eq.~(\ref{eq:cov_rotation}), i.e. $\sigma^{\text{off}} = R_{\bar{\theta}} D (\theta_\Delta) R_{\bar{\theta}}^{\rm T}$, where $D (\theta_\Delta) = \text{diag} (e^{i \psi (\theta_\Delta)} |\zeta (\theta_\Delta)|, e^{-i \psi(\theta_\Delta)} /|\zeta(\theta_\Delta)|) $ is diagonal in its polar decomposition. Hence, by moving in the rotating frame with respect of the angle $\bar{\theta}$, i.e. $R_{\bar{\theta}} \bar{r}^{\text{off}} \to \bar{r}^{\text{off}}$, and changing the variable from $\tau$ to $\theta_\Delta$, one finds that the ODE for the first moment is 
\begin{equation}
   \frac{ \text{d} \bar{r}^{\text{off}}}{\text{d} \theta_{\Delta}} =   - i  D (\theta_\Delta) \bar{r}^{\text{off}} \implies \bar{r}^{\text{off}} = \exp \left( - i \int_{0}^{\theta_{\Delta}} \text{d} t D (t) \right) r_0 = \begin{pmatrix}
\frac{s  x_0}{s \cos \theta_\Delta +i \sin \theta_\Delta }  \\
 \frac{p_0}{\cos \theta_\Delta +i s \sin \theta_\Delta )}  
   \end{pmatrix}
\end{equation}
Given that the trace of the covariance and the norm of the first moments are both invariant under rotations, the QRDM ODE in the rotating frame reads
\begin{align}
    \frac{ \text{d} \bar{r}^{\text{(0) off}}}{\text{d} \theta_{\Delta}}   &= - i \left[ \frac{ x_0^2}{\left(  \cos \theta_\Delta +i / s \sin \theta_\Delta  \right)^2 } +\frac{p_0^2}{\left( \cos \theta_\Delta +i s \sin \theta_\Delta \right)^2} + \frac{1}{2} \left(\frac{s \cos \theta_\Delta +i \sin \theta_\Delta }{\cos \theta_\Delta +i s \sin \theta_\Delta }+\frac{\cos \theta_\Delta +i s \sin \theta_\Delta }{s \cos \theta_\Delta +i \sin \theta_\Delta }\right) \right] \nonumber
\end{align}
The off diagonal exponent reads
\begin{align}  
\bar{r}^{\text{(0) off}} = -i \sin  \theta_\Delta \left(\frac{p_0^2}{\cos  \theta_\Delta+i s \sin  \theta_\Delta}+\frac{s x_0^2}{s \cos  \theta_\Delta+i \sin  \theta_\Delta}\right)-\frac{1}{2} \log \left(\frac{1}{s}(s \cos  \theta_\Delta+i \sin  \theta_\Delta) (\cos  \theta_\Delta+i s \sin  \theta_\Delta)\right)+\log(q_0)\nonumber
\end{align}
where $q_0$ is the initial state of the QRDM that, after the rotation reads 
\begin{equation}
    \rho_q^{\text{off}} = q \frac{\sqrt{s} \exp \left( -i \sin  \theta_\Delta \left(\frac{p_0^2}{\cos  \theta_\Delta+i s \sin  \theta_\Delta}+\frac{s x_0^2}{s \cos  \theta_\Delta+i \sin  \theta_\Delta}\right) \right) }{ \sqrt{s \cos  \theta_\Delta+i \sin  \theta_\Delta} \sqrt{\cos  \theta_\Delta+i s \sin  \theta_\Delta }} 
\end{equation}


\subsection{Conditional Displacements}

The case of conditional displacements, with set of vector $d_j$, one can directly compute 
\begin{align}
    \Phi^d[\hat{\varrho}] &= \frac{1}{(2 \pi)^n} \sum_{j, k  = 0}^{N-1}  \int_{\mathds{R}^{2n}} \text{d} r \; \chi_{jk} (r)  \hat{D}_{d_j} \hat{D}_{r} \hat{D}_{d_k}^\dag \otimes \ket{j} \bra{k}  \\
    &= \frac{1}{(2 \pi)^n} \sum_{j, k  = 0}^{N-1} \int_{\mathds{R}^{2n}} \text{d} r \; \chi_{jk} (r)  {\rm e}^{ \frac{i}{2} r^{\rm T} \Omega (d_j + d_k) + \frac{i}{2}  d_j^{\rm T}  \Omega d_k }\hat{D}_{r +d_j - d_k}  \otimes \ket{j} \bra{k} \nonumber
\end{align}
where we repetitivelly used $ \hat{D}_{r_1} \hat{D}_{r_2} = {\rm e}^{- \frac{i}{2} r_1^{\rm T} \Omega r_2} \hat{D}_{r_1+r_2} =  {\rm e}^{\frac{i}{2} r_2^{\rm T} \Omega r_1} \hat{D}_{r_1+r_2}$. By changing the integration variable to $r' = r+ r_j - r_k$, one finds that the updated characteristic function after the operation is
\begin{equation}
    \chi_{jk}^{\Phi} (r') = \chi_{jk} (r' - d_j + d_k) \exp \left( { - \frac{i}{2} \bar{r'}^{\rm T} \Omega^{\rm T} (d_j+d_k) - \frac{i}{2}  d_j^{\rm T}  \Omega d_k  } \right)
\end{equation}
Substituting in Gaussian $\chi_{jk}$, it reads
\begin{align}
    \chi_{jk}^{\Phi} (r') =  \exp \Big[ -\frac{1}{4} &r'^{\rm T} \Omega^{\rm T}\sigma_{jk} \Omega r'  +  i  r'^{\rm T} \Omega^{\rm T} \left( \bar{r}_{jk} - \frac{1}{2} (d_j +d_k) - \frac{i}{2} \sigma_{jk} \Omega  (d_j - d_k) \right)  \\
    &+ r^{(0)}_{jk} - \frac{1}{4} \left( d_j - d_k\right)^{\rm T} \Omega^{\rm T} \sigma_{jk}  \Omega \left( d_j - d_k\right) +  i (d_j - d_k)^{\rm T} \Omega \bar{r}_{jk} - \frac{i}{2}  d_j^{\rm T}  \Omega d_k  \Big)
\end{align}
from which one can read the update rules given in Eq.~(\ref{eq:cond_dis_update}). 


\section{Perturbation Theory \label{app:perturbation}}



Greek symbols labels the components of vectorial representation of some quantity and the Einstein summations are performed up to $N^2$. 
For a single qubit, examples of DV non-orthogonal superoperator that frequency arises are the $\hat{\sigma}_x$ Hamiltonian term and $T_1$ Lindbladian term, which takes the form in the vectorized phase space as
%\begin{equation}
%   \text{vec} ( i \left[ \cdot, \hat{\sigma}_x \right] ) \to i \left( \mathds{1}  \otimes \hat{\sigma}_x  - \hat{\sigma}_x  \otimes \mathds{1} \right)  \dket{\chi} = \begin{pmatrix}
%       0 &-i & i & 0 \\
%       -i & 0 & 0 & i \\ 
%       i & 0 & 0 & -i \\ 
%       0 &i & -i & 0 
%   \end{pmatrix} \dket{\chi}
%\end{equation}
\begin{equation}
    \text{vec} \left[ 
        \sigma_{-}\rho\,\sigma_{+}
        -\tfrac{1}{2}\sigma_{+}\sigma_{-}\rho
        -\tfrac{1}{2}\rho\,\sigma_{+}\sigma_{-}
     \right] \longleftrightarrow - \gamma \frac{2 + j + k}{4} \chi_{jk}(\tilde{\xi}) +  \gamma \delta_{j,-1}\delta_{k,-1}\,\chi_{+1,+1}
\end{equation}

Consider the PDEs $\left( \partial_\tau - \mathcal{L}_{\alpha} \right) \chi_{\alpha} = V_{\alpha \beta} \chi_{\beta}$. The solution can be formulated in terms of Green's functions, which solve $\left( \partial_\tau - \mathcal{L}_{\alpha} \right) \left[\mathcal{G}_{\alpha \alpha} (\tau,  \tilde{\xi}, \tilde{\xi}') \right]  =0$ under the boundary conditions $\mathcal{G}_{\alpha \alpha} (0,  \tilde{\xi}, \tilde{\xi}') = \delta( \tilde{\xi}- \tilde{\xi}')$. Then, one can write the solution of the coupled PDEs as the sum of the homogenous soultion solving for some initial condition and a particular solution of the inhomogenous, such that 
\begin{equation}
    \chi_\alpha (\tilde{\xi}, \tau) = \chi^h_\alpha (\tilde{\xi}, \tau) +\chi^p_\alpha (\tilde{\xi}, \tau) =\chi^h_\alpha (\tilde{\xi}, \tau) +  \int_0^{\tau} {\rm d} t \int_{\mathds{R}^{2n}} {\rm d} \tilde{\xi}'  \mathcal{G}_{\alpha \alpha} (\tau - t,  \tilde{\xi}, \tilde{\xi}') V^{\alpha \beta} \chi_{\beta} (\tilde{\xi}', t)
\end{equation}
Given that the homogenous solution solves $\left( \partial_\tau - \mathcal{L}_{\alpha} \right) \chi^h_{\alpha} = 0$, one finds
\begin{align}
    \left( \partial_\tau - \mathcal{L}_{\alpha} \right) \chi_{\alpha} (\tilde{\xi}, \tau) &= \left( \partial_\tau - \mathcal{L}_{\alpha} \right)  \chi^i_\alpha (\tilde{\xi}, \tau) \nonumber\\  
    &=  \int_{\mathds{R}^{2n}} {\rm d} \tilde{\xi}' \delta(\tilde{\xi} - \tilde{\xi}') V^{\alpha \beta} \chi_{\beta} (\tilde{\xi}', \tau)  + \int_0^{\tau} {\rm d} t \int_{\mathds{R}^{2n}} {\rm d} \tilde{\xi}'  V^{\alpha \beta} \chi_{\beta} (\tilde{\xi}', t)  \left( \partial_\tau - \mathcal{L}_{\alpha} \right) \left[\mathcal{G}_{\alpha \alpha} (\tau - t,  \tilde{\xi}, \tilde{\xi}') \right] = V^{\alpha \beta} \chi_{\beta} (\tilde{\xi}, \tau) \nonumber
\end{align}
thus, recovering the PDE. Let us now take a Gaussian ansatz for the Green's function
\begin{equation}
    \mathcal{G}_\alpha(\tilde{\xi}, \tilde{\xi}', \tau) =   \exp \left( -\frac{1}{4} \tilde{\xi}^{\rm T} \Sigma_\alpha (\tilde{\xi}',  \tau) \tilde{\xi}  +  i  \tilde{\xi}^{\rm T} R_\alpha (\tilde{\xi}', \tau) +R^{(0)}_\alpha (\tilde{\xi}', \tau)\right) 
\end{equation}
to solve for $\left( \partial_\tau - \mathcal{L}_{\alpha} \right) \left[\mathcal{G}_{\alpha \alpha} (\tau,  \tilde{\xi}, \tilde{\xi}') \right]  =0$. From the same argument used in this work to solve any dynamics, in order to be a valid Ansatz, $\Sigma_\alpha$, $R_\alpha$ and $R^{(0)}_\alpha$ has to solve the vectorial form of the ODEs derived in the other Sections. The boundary condition can be formulated as the limit of a Gaussian of covariance matrix $\tilde{\Sigma}$
\begin{align}
    \mathcal{G}_{\alpha}(\tilde{\xi}, \tilde{\xi}', 0) &= \delta(\tilde{\xi}-\tilde{\xi}') = \lim_{\tilde{\Sigma} \to 0}  \frac{2^n}{\pi^n \sqrt{\text{det} \tilde{\Sigma} }} \exp \left( - \left(\tilde{\xi}- \tilde{\xi}'\right)^{\rm T} \tilde{\Sigma}^{-1} \left(\tilde{\xi}- \tilde{\xi}'\right) \right) 
\end{align}
which implies that the one has to compute the dynamics with initial conditions, and 
\begin{equation}
    \Sigma_\alpha (\tilde{\xi}',0) = 4 \tilde{\Sigma}^{-1} \hspace{1cm} R_\alpha (\tilde{\xi}',0) = - i 2  \tilde{\Sigma}^{-1} \tilde{\xi}' \hspace{1cm} R^{0}_{\alpha}(\tilde{\xi}',0)  =  \ln \left[ \frac{2^n}{\pi^n \sqrt{\text{det} \tilde{\Sigma} }} \exp(- \tilde{\xi}'^{ \; \rm T } \tilde{\Sigma}^{-1} \tilde{\xi}') \right]
\end{equation}

The Volterra integral can be solved as von Neumann series with  the $N$-th order given by $0<t_N<\cdots<t_1<\tau$:
\begin{equation}
    \hspace{-1cm}
    \dket{\chi^{(N)}}(\tau,\tilde{\xi}) = \!\!\int\limits_{0<t_N<\cdots<t_1<\tau} \!\!\!\! {\rm d}t_1\cdots{\rm d}t_N \int_{(\mathds{R}^{2n})^N} \!\!\!\! {\rm d}\tilde{\xi}_1\cdots{\rm d}\tilde{\xi}_N\ \mathcal{G}(\tau-t_1,\tilde{\xi},\tilde{\xi}_1)\,\mathds{V}\, \mathcal{G}(t_1-t_2,\tilde{\xi}_1,\tilde{\xi}_2)\,\mathds{V}\cdots \mathcal{G}(t_{N-1}-t_N,\tilde{\xi}_{N-1},\tilde{\xi}_N)\,\mathds{V}\,\dket{\chi^{h}(t_N,\tilde{\xi}_N)}\ ,
\end{equation}

\subsection{DV Decay ($T_1$) \label{app:qubit_decay}}

Let us first consider the case of a single qubit. The jump operator $\hat{L}_1 = \sqrt{\Gamma_1}\, \hat{\sigma}_-$ generates the Lindbladian whose vectorized phase-space representation reads
\begin{equation}
    \Gamma_1 \Big( \sigma_-\rho\,\sigma_+ - \tfrac{1}{2}\{\sigma_+\sigma_-,\rho\} \Big) \;\longleftrightarrow\; -\Gamma_1 \frac{2+j+k}{4}\,\chi_{jk}(\tilde{r}) \;+\; \Gamma_1\,\delta_{j,-1}\delta_{k,-1}\,\chi_{+1,+1}(\tilde{r})\,, \qquad j,k \in \{\pm 1\}\,.
\end{equation}
The diagonal (first) term is absorbed into the homogeneous generator $\mathcal{L}$, providing the damping rates $\{\Gamma_1, \Gamma_1/2, \Gamma_1/2, 0\}$ of the four branches, while the source (second) term defines the perturbation which, in the ordered basis $\dket{\chi} = (\chi_{++}, \chi_{+-}, \chi_{-+}, \chi_{--})^{\rm T}$, reads
\begin{equation}
    \mathds{V} = \Gamma_1 \begin{pmatrix} 0&0&0&0 \\ 0&0&0&0 \\ 0&0&0&0 \\ 1&0&0&0 \end{pmatrix}\,.
\end{equation}
It is possible to note that $\mathds{V}\,\mathcal{G}\,\mathds{V} = 0$, hence $\mathds{k}^{*k} = 0$ for all $k \geq 2$, such that the Liouville--Neumann series truncates at first order, reflecting the fact that a single DV excitation supports at most one quantum jump. 

The diagonal terms, which we shall just label with $\pm$,  evolve under the Gaussian differential operators $\mathcal{L}^{\pm} = -\frac{1}{4} D^{op}\, \tilde{r}^o \tilde{r}^p + A^{op}_{\pm}\, \tilde{r}^o \partial_{\tilde{r}^p} + i\, r^m_{\pm} \Omega^{mn} \tilde{r}^n$, with $A_\pm = \Omega(E + H_\pm)$ and $r_{\pm} = r_m \pm r_q$. Inserting the Gaussian ansatz for $\mathcal{G}^{\pm}$ yields the Lyapunov-type ODEs for its covariance matrix and frist moments: $\dot{\Sigma}_\pm = A_\pm \Sigma_\pm + \Sigma_\pm A^{\rm T}_\pm + D$ and $\dot{R}_\pm = A_\pm R_\pm - \Omega r_\pm$. Thus, by imposing the initial conditions enforcing $\mathcal{G}^{\pm}(0, \tilde{r}, \tilde{r}') = \delta(\tilde{r} - \tilde{r}')$ previously described, one can write the prpagator in the closed form
\begin{align}
\label{eq:green_qubit}
    \mathcal{G}^{\pm}(\tau, \tilde{r}, \tilde{r}') &=  \exp \left( -\frac{1}{4} \tilde{r}^{\rm T}  \left[\int_{0}^{\tau} \text{d}t S_{A_\pm}(\tau - t) D S^{\rm T}_{A_\pm}(\tau - t) \right]  \tilde{r}  - i  \tilde{r}^{\rm T} \left(S_{A_\pm}(\tau) - \mathds{1} \right)  r_{\pm}   \right) \nonumber \\
&\hspace{0.5cm}\cdot \lim_{\tilde{\Sigma} \to 0}  \frac{2^n}{\pi^n \sqrt{\text{det} \tilde{\Sigma} }}  \exp \left( - \tilde{r}^{\rm T} S_{A_\pm}(\tau) \tilde{\Sigma}^{-1} S_{A_\pm}^{\rm T}(\tau) \tilde{r}  + 2  \tilde{r}^{\rm T} S_{A_\pm} (\tau)   \tilde{\Sigma}^{-1} \tilde{r}' - \tilde{r}^{' \; \rm T } \tilde{\Sigma}^{-1} \tilde{r}' \right) \nonumber \\  
&= \exp\!\left( -\tfrac{1}{4}\, \tilde{r}^{\rm T} D_\pm(\tau)\, \tilde{r} \;-\; i\, \tilde{r}^{\rm T} r^f_\pm(\tau) \right) \delta\!\left( S^{\rm T}_{A_\pm}(\tau)\, \tilde{r} - \tilde{r}' \right)\,,
\end{align}
with $S_{A_\pm}(\tau) = e^{A_\pm \tau}$, $D_\pm(\tau) = \int_0^\tau {\rm d}t\, S_{A_\pm}(\tau - t)\, D\, S^{\rm T}_{A_\pm}(\tau - t)$ and $r^f_\pm(\tau) = \left( S_{A_\pm}(\tau) - \mathds{1} \right) r_\pm$. The homogenous term can be re-derived by acting with Eq.~(\ref{eq:green_qubit}) on a Gaussian initial condition $\chi_0(\tilde{r}) = \exp( -\frac{1}{4} \tilde{r}^{\rm T} \sigma_0 \tilde{r} + i \tilde{r}^{\rm T} r_0 )$: the delta function trivializes the phase-space integral and one recovers the closed-form Gaussian solutions of Sec.~\ref{sec:dynamics}, with moments $\sigma_\pm(\tau) = S_{A_\pm} (\tau)\sigma_0 S^{\rm T}_{A_\pm} (\tau)+ D_\pm(\tau)$ and $\bar{r}_\pm(\tau) = S_{A_\pm}  (\tau) r_0 - r^f_\pm(\tau)$. Including its diagonal damping and initial population $p_+$, one finds  $\chi_{+}^{h \text{on}}(t, \tilde{r}) = p_+\, e^{-\Gamma_1 t}\, \exp( -\frac{1}{4} \tilde{r}^{\rm T} \sigma_+(t)\, \tilde{r} + i\, \tilde{r}^{\rm T} \bar{r}_+(t) )$. By the structure of $\mathds{V}$, the only non-vanishing component of Eq.~(\ref{eq:order1}) is the ground-state one,
\begin{align}
    \chi^{(1)}(\tau, \tilde{r}) &= \Gamma_1 \int_0^\tau {\rm d}t \int_{\mathds{R}^{2n}} {\rm d}\tilde{r}'\; \mathcal{G}^{-}(\tau - t, \tilde{r}, \tilde{r}')\; \chi^{h}_{++}(t, \tilde{r}') 
    = p_+\, \Gamma_1 \int_0^\tau {\rm d}t\; e^{ -\frac{1}{4} \tilde{r}^{\rm T} \sigma^{\Gamma}(\tau, t)\, \tilde{r} \,+\, i\, \tilde{r}^{\rm T} \bar{r}^{\Gamma}(\tau, t) \,-\, \Gamma_1 t }\,,
\end{align}
where the $\tilde{r}'$-integral is again resolved by the delta function of Eq.~(\ref{eq:green_qubit}), giving the moments of the open dynamics
\begin{align}
\label{eq:decay_moments_open}
    \sigma^{\Gamma}(\tau, t) &= S_{A_-}(\tau - t) \left[ S_{A_+}(t)\, \sigma_0\, S^{\rm T}_{A_+}(t) + D_+(t) \right] S^{\rm T}_{A_-}(\tau - t) + D_-(\tau - t)\,, \nonumber \\
    \bar{r}^{\Gamma}(\tau, t) &= S_{A_-}(\tau - t) \left[ S_{A_+}(t)\, r_0 - \left( S_{A_+}(t) - \mathds{1} \right) r_+ \right] - \left( S_{A_-}(\tau - t) - \mathds{1} \right) r_-\,.
\end{align}
Also in the open dynamics case, the state evolves in the excited branch ($+$) up to the jump time $t$, when is transferred to the ground branch by $\mathds{V}$, and evolves in the ground branch ($-$) for the remaining time $\tau - t$; the diffusion accumulated in each branch is transported accordingly. Setting $D = 0$ and $E = 0$ recovers Eq.~(\ref{eq:decay_moments_open}). 

For a $N$ level system, the decay Lindbladian terms maps to the phase space as 
\begin{equation}
    \text{Eq.(\ref{eq:Nlevel_lindblad})} \rightarrow  -\frac{\Gamma_j + \Gamma_k}{2}\, \chi_{jk}(\tilde{r}) \;+\; \delta_{jk}\, \Gamma_{j+1}\, \chi_{j+1,\,j+1}(\tilde{r})\,, \implies  V_{(j,k),(j',k')} = \Gamma_{j+1}\, \delta_{jk}\, \delta_{j'k'}\, \delta_{j',\, j+1}\,,
\end{equation}
with $j,k\in[1,...,N-1]$. The former terms can be reabsorbed in the diagonal solutions, leading to exponential decay of all the entries with the exception of the ground state ($j=k=0$),  while the later, only concerns population transfer between diagonal elements, for which the propagators can be expressed as delta functions, imposing the symplectic transformation update at each jump. On the analogy of the case of a single qubit, for a $N$ dimensional DV, there are $N$ possible jumps (from the most exited state to the ground state) which truncates the series. Each phase space jump can be accounted as phase space parameters as a recursive Eq.~(\ref{eq:two_times_parameters}). 



\section{Phase-Space Parameters in Physical Examples \label{app:examples}}

\subsection{Two Qubits and a squeezed Resonator}

On the diagonal terms, the dynamics can be computed via the two non-symplectic transformation
\begin{equation}
    S^{\text{on}}_\pm = {\rm e}^{\tau \Omega (H_{\pm} + E} ) = e^{-\kappa\tau/2}
\begin{pmatrix}
\cosh(\omega_{\mu, \lambda} \tau)+\dfrac{ \lambda}{\omega_{\mu, \lambda}}\sinh(\omega_{\mu, \lambda} \tau) & \pm\dfrac{\mu}{\omega_{\mu, \lambda}}\sinh(\omega_{\mu, \lambda} \tau) \\
\mp\dfrac{\mu}{\omega_{\mu, \lambda}}\sinh(\omega_{\mu, \lambda} \tau) & \cosh(\omega_{\mu, \lambda} \tau)-\dfrac{ \lambda}{\omega_{\mu, \lambda}}\sinh(\omega_{\mu, \lambda} \tau)
\end{pmatrix}
\nonumber
\end{equation}
with $\omega_{\mu, \lambda} = \sqrt{\mu^2- \lambda^2}$. From the solution to the Lyapunov equation  (Tab~\ref{tab:open} b), one can compute the time evolution of the covariance matrices and the first moments. At the steady state, one finds that 
  \begin{align}
       \sigma_{\pm}^\text{on} (\tau \to \infty) &=\frac{1}{4  \mu ^2+\kappa ^2-4 \lambda ^2}
\begin{pmatrix}
 4   \mu ^2+\kappa ^2+2 \kappa  \lambda  & \mp 4  \lambda  \mu  \\
 \mp 4  \lambda  \mu  & 4   \mu ^2+\kappa ^2+2 \kappa  \lambda  \\
\end{pmatrix} \\
\bar{r}_{\pm}^\text{on} (\tau \to \infty) &= \frac{\alpha  \sqrt{2 \kappa }}{4   \mu ^2+\kappa ^2-4 \lambda ^2}
\begin{pmatrix}
 2 j \mu  \cos (\phi )-(\kappa +2 \lambda ) \sin (\phi ) \\
 (\kappa -2 \lambda ) \cos (\phi )+2 j \mu  \sin (\phi ) \\
\end{pmatrix}
   \end{align}


Define the two branch frequencies $
    \omega_{\pm} = 
    \frac{1}{4} \left[ (\kappa \pm 2 \lambda)^{2} + 4 i \kappa \mu - 4 \mu^{2} \right] $. Solving the
Riccati equation by the doubling method of Sec.~\ref{subsec:riccati}, with the vacuum
initial condition $\sigma_{0} = \mathds{1}$, gives the diagonal off diagonal covariance
\begin{equation}
    \sigma^{\text{off}} (\tau) = \text{diag} \left( 
        \frac{(\kappa + 2 \lambda + 2 i \mu) \sinh \left( \tau \sqrt{\omega_{-}} \right)
        + 2 \sqrt{\omega_{-}} \cosh \left( \tau \sqrt{\omega_{-}} \right)}
        {(\kappa - 2 \lambda + 2 i \mu) \sinh \left( \tau \sqrt{\omega_{-}} \right)
        + 2 \sqrt{\omega_{-}} \cosh \left( \tau \sqrt{\omega_{-}} \right)} ,
       \frac{(\kappa - 2 \lambda + 2 i \mu) \sinh \left( \tau \sqrt{\omega_{+}} \right)
        + 2 \sqrt{\omega_{+}} \cosh \left( \tau \sqrt{\omega_{+}} \right)}
        {(\kappa + 2 \lambda + 2 i \mu) \sinh \left( \tau \sqrt{\omega_{+}} \right)
        + 2 \sqrt{\omega_{+}} \cosh \left( \tau \sqrt{\omega_{+}} \right)}
    \right) \nonumber
\end{equation}
The first moments follow from the propagator $\mathds{P}_{jk}$ of Eq.~(\ref{eq:propagator_r_off})
via $\bar{r}^{\text{off}} = \mathds{P}_{jk} (\tau) \int_{0}^{\tau} \text{d} t \,
\mathds{P}^{-1}_{jk} (t) \, \Omega d$, with $d = \sqrt{2 \kappa} \, \alpha (\cos \phi , \sin \phi)^{\rm T}$
and $\bar{r}_0 = (0,0)^{\rm T}$, 
\begin{equation}
    \bar{r}^{\text{off}} (\tau) = 2 \sqrt{2} \, \alpha \sqrt{\kappa}
    \begin{pmatrix}
        \sin \phi \;
        \dfrac{\dfrac{\kappa - 2 \lambda + 2 i \mu}{2 \sqrt{\omega_{-}}}
        \left[ \cosh \left( \tau \sqrt{\omega_{-}} \right) - 1 \right]
        + \sinh \left( \tau \sqrt{\omega_{-}} \right)}
        {(\kappa - 2 \lambda + 2 i \mu) \sinh \left( \tau \sqrt{\omega_{-}} \right)
        + 2 \sqrt{\omega_{-}} \cosh \left( \tau \sqrt{\omega_{-}} \right)} \\[24pt]
        - \cos \phi \;
        \dfrac{\dfrac{\kappa + 2 \lambda + 2 i \mu}{2 \sqrt{\omega_{+}}}
        \left[ \cosh \left( \tau \sqrt{\omega_{+}} \right) - 1 \right]
        + \sinh \left( \tau \sqrt{\omega_{+}} \right)}
        {(\kappa + 2 \lambda + 2 i \mu) \sinh \left( \tau \sqrt{\omega_{+}} \right)
        + 2 \sqrt{\omega_{+}} \cosh \left( \tau \sqrt{\omega_{+}} \right)}
    \end{pmatrix} ,
\end{equation}
The ODE of the QRDM is
\begin{align}
   \dot{r}^{(0), \text{off}} =& - 8 i \alpha^{2} \kappa \mu \Big[ 
    \cos^{2} \phi_d \;
    \frac{\left( \dfrac{\kappa + 2 \lambda + 2 i \mu}{2 \sqrt{\omega_{+}}}
    \left[ \cosh \left( \tau \sqrt{\omega_{+}} \right) - 1 \right]
    + \sinh \left( \tau \sqrt{\omega_{+}} \right) \right)^{2}}
    {\left( (\kappa + 2 \lambda + 2 i \mu) \sinh \left( \tau \sqrt{\omega_{+}} \right)
    + 2 \sqrt{\omega_{+}} \cosh \left( \tau \sqrt{\omega_{+}} \right) \right)^{2}} \\
    &\hspace{2cm}
    + \sin^{2} \phi_d \;
    \frac{\left( \dfrac{\kappa - 2 \lambda + 2 i \mu}{2 \sqrt{\omega_{-}}}
    \left[ \cosh \left( \tau \sqrt{\omega_{-}} \right) - 1 \right]
    + \sinh \left( \tau \sqrt{\omega_{-}} \right) \right)^{2}}
    {\left( (\kappa - 2 \lambda + 2 i \mu) \sinh \left( \tau \sqrt{\omega_{-}} \right)
    + 2 \sqrt{\omega_{-}} \cosh \left( \tau \sqrt{\omega_{-}} \right) \right)^{2}} \Big] \nonumber \\
 &- \frac{i \mu}{2} \left[ 
    \frac{(\kappa - 2 \lambda + 2 i \mu) \sinh \left( \tau \sqrt{\omega_{+}} \right)
    + 2 \sqrt{\omega_{+}} \cosh \left( \tau \sqrt{\omega_{+}} \right)}
    {(\kappa + 2 \lambda + 2 i \mu) \sinh \left( \tau \sqrt{\omega_{+}} \right)
    + 2 \sqrt{\omega_{+}} \cosh \left( \tau \sqrt{\omega_{+}} \right)} 
    + \;  \frac{(\kappa + 2 \lambda + 2 i \mu) \sinh \left( \tau \sqrt{\omega_{-}} \right)
    + 2 \sqrt{\omega_{-}} \cosh \left( \tau \sqrt{\omega_{-}} \right)}
    {(\kappa - 2 \lambda + 2 i \mu) \sinh \left( \tau \sqrt{\omega_{-}} \right)
    + 2 \sqrt{\omega_{-}} \cosh \left( \tau \sqrt{\omega_{-}} \right)} \right] \nonumber.
\end{align}
where the real part of the first two lines is $\dot{\mathcal{C}}_{\bar{r}}$ and  the last line is $\dot{\mathcal{C}}_{\sigma}$.
For conciseness, here, we have set $H_{q}^{0} = 0$, as a pure frequency shift that does not affect the dephasing rate.
In the limit $\tau \to \infty$ one has
$\dot{\mathcal{C}}_{\bar{r}}  \to \text{Re} \left( \sqrt{\omega_{+}} + \sqrt{\omega_{-}} \right) / 2 - \kappa / 2$
and
$\dot{\mathcal{C}}_{\bar{r}}  \to 2 \mu^{2} \kappa^{2} \alpha^{2}
\left( \cos^{2} \phi / | \omega_{+} |^{2} + \sin^{2} \phi / | \omega_{-} |^{2} \right)$,
recovering Eqs.~(2) and~(4) of Ref.~\cite{eddins_high_2019}.


$$\mathcal S(\tau);\equiv;\frac{(\Delta m)^2}{v_\uparrow+v_\downarrow}
=\frac{\left[u^{\rm T}\left(r_{\uparrow\uparrow}(\tau)-r_{\downarrow\downarrow}(\tau)\right)\right]^{2}}
{u^{\rm T}\left(\sigma_{\uparrow\uparrow}(\tau)+\sigma_{\downarrow\downarrow}(\tau)\right)u+2(1-\eta)/\eta},$$
which is your expression. In variance language $\mathcal S=(\Delta m)^2/\big[2(V^{(u)}\uparrow+V^{(u)}\downarrow)\big]$, i.e. half the conventional $\mathrm{SNR}^2=(\Delta m)^2/(V_\uparrow+V_\downarrow)$ — this is the convention mismatch I kept flagging, and it is why the comparison below is against $\mathcal S/2$ and not $\mathcal S$

 On the diagonal term, the steady state CV
\begin{align}
0 &=   \left( j \mu \Omega  + \lambda \sigma_x  - \kappa \mathds{1} \right)    \sigma_{jj}^s - \sigma_{jj}^s  \left( j \mu \Omega  + \lambda \sigma_x  - \kappa \mathds{1} \right)  + \kappa \mathds{1} \hspace{1cm} 
        0  = \left( j \mu \Omega  + \lambda \sigma_x  - \kappa \mathds{1} \right)    \bar{r}_{jj}^s + 
      \Omega  d 
   \end{align}
 
Similarly, for the off-diagonal, for $j= 1$ and $k=-1$,
\begin{align}
0 & = \left(  \lambda \sigma_x  - \kappa \mathds{1} \right)    \sigma_{jk}^s - \sigma_{jk}^s  \left(\lambda \sigma_x  - \kappa \mathds{1} \right)  + \kappa \mathds{1} + i  \mu \left( \sigma_{jk}^s \sigma_{jk}^s - \mathds{1} \right)
\nonumber  \\
        0 & = \left[  \frac{j+k}{2} \mu \Omega  + \lambda \sigma_x  - \kappa \mathds{1} - i  \frac{j-k}{2} \mu \sigma_{jk}^s \right]    \bar{r}_{jk}^s + 
      \Omega  d  
   \end{align}








\subsection{Stern-Gerlach Interferometer}

When decay in the DV is not considered, the time evolution of the phase space parameters of the Stern-Gerlach Interferometer undergoing diffusive dynamics with an unknown external force can be computed, from Eq.~(\ref{eq:Lyapunov_sol}) and Tab.~\ref{tab:open_app}. This is achieved by a single symplectic transformation describing the quadratic evolution of the quantum harmonic oscillator of the mass, which reads
\begin{equation}
    S_{\text{qho}}(\tau) := {\rm e}^{\tau \Omega H} = \begin{pmatrix}
        \cos \tau & - \sin \tau \\ 
        \sin \tau & \cos \tau
    \end{pmatrix} \;.
\end{equation}
describing a rotation of the phase space. Thus, the phase space parameters of an initial squeezed thermal state with first and second moments $r_0 = (0,0)^{\rm T}$  and $\sigma_0 = (1+2N_p) \text{diag}(s, 1/s)$, respectively, and a qubit in initial state given by the density matrix  $\varrho^q(0) = \ket{+}\bra{+}$ evolves to a GBCS with the time dependent phase space parameters are, a single covariance matrix
\begin{equation} \label{eq:covariance_SG}
    \hspace{-1cm}\sigma(\tau) = (1 + 2 N_p) 
\begin{pmatrix}
 s \cos ^2(\tau) +\sin ^2(\tau) /s &  \left(1/s - s \right) \sin \tau  \cos \tau  \\
 \left(1/s - s \right) \sin \tau  \cos \tau   & s \sin ^2(\tau) +\cos ^2(\tau) /s \\    
\end{pmatrix}
+ \Gamma_x
\begin{pmatrix}
     \tau -\sin \tau  \cos \tau  & \sin ^2(\tau)  \\
 \sin ^2(\tau)  & \tau +\sin \tau  \cos \tau  \\
\end{pmatrix}
\end{equation}
where we recognise the first term as the covariance matrix of a squeezed thermal state undergoing a unitary evolution in a QHO and the second being the additional term adding diffusion (with a linear increment in $\tau$). The tree independent vectors of first moments are
\begin{equation}
    \bar{r}^{\text{on}}_{\pm} (\tau)
    = \left(f_u \mp f_q\right)
\begin{pmatrix}
1 - \cos \tau  \\
  - \sin \tau  
\end{pmatrix} \nonumber
\end{equation}
\begin{equation}
 \bar{r}^{\text{off}}(\tau) = f_u \begin{pmatrix}
          1 - \cos \tau  \\
        \sin \tau 
    \end{pmatrix} + i f_q 
\left[ (1+2 N_p)
\begin{pmatrix}
  \sin \tau  \left( \left(1/s - s\right) \cos \tau  - 1/s \right) \\
 2 \sin ^2\left(\frac{\tau }{2}\right) \left( \left(s-1/s\right) \cos \tau + s \right) \\
\end{pmatrix}
+ \Gamma_x
\begin{pmatrix}
  - 4 \sin^4(\frac{\tau}{2}) \\
  \tau + \sin \tau (\cos \tau - 2)\\


  
\end{pmatrix}
\right] \nonumber
\end{equation}
and, the QRDM parameters are
\begin{equation}
     \mathcal{C} (\tau) = f_q^2 \left[2 (1 + 2 N_p)  \sin ^2\left(\frac{\tau }{2}\right) \left(\left(s-\frac{1}{s}\right) \cos (\tau )+s+\frac{1}{s}\right) +  \Gamma_x \left( 3 \tau + \sin \tau (\cos \tau - 4) \right) \right] + \Gamma_z \tau + \Gamma_1 \tau/2
\end{equation}
and $\phi(\tau) = 2 f_q f_u (\tau +\sin \tau ) + \frac{1}{2} \omega_q \tau $. 

In order to account for the contribution of the decay in the DV, let us first compute the phase-space parameters of the unravelling of the inhomogeneous term of Eq.~(\ref{eq:decay_qubit}) via Eq.~(\ref{eq:decay_moments_open}). We note that the covariance matrix of is given by Eq.~(\ref{eq:covariance_SG}) as the SG coupling, being linear, does not affect the second moments, i.e. $S_{A_\pm} = S_{\text{qho}}$. The first moment can be ready computed using $r_\pm = r_m \pm r_q = (f_u \pm f_q,\;0)^{\rm T}$ and $r_0 = (0,0,)^{\rm T}$, such that
\begin{align}
    \bar{r}^{\Gamma}(\tau, t) 
    &= -  \left( S_{\text{qho}}(\tau)  - \mathds{1}  \right) r_m -  \left( S_{\text{qho}}(\tau) -  2 S_{\text{qho}}(\tau - t) + \mathds{1}  \right) r_q \\
    &= f_u \begin{pmatrix}
1 - \cos \tau  \\
   - \sin \tau  
\end{pmatrix} + f_q \begin{pmatrix}
-1 - \cos \tau  \\
  - \sin \tau  
\end{pmatrix} + 2 f_q \begin{pmatrix}
 \cos ( \tau  - t) \\
   \sin (\tau-t)  
\end{pmatrix} =  \bar{r}^{\Gamma}(\tau) ++ 2 f_q \begin{pmatrix}
 \cos ( \tau  - t) \\
   \sin (\tau-t)  
\end{pmatrix}
\end{align}
The last term is the only quantity which is function of $t$, such that one can write the charactertic function of the decayed part as
\begin{equation}
    \chi^{(1) } (\tau, \tilde{\xi})  = p_+ \Gamma_1  {\rm e}^{-\frac{1}{4} \tilde{\xi}^{\rm T} \sigma (\tau) \tilde{\xi}  +  i  \tilde{\xi}^{\rm T} \bar{r}^{\Gamma} ( \tau ) }  \ I_\Gamma (\tau, \tilde{\xi}) \Gamma_1
\end{equation}
where 
one can compute the time integral via the Basel functions $J_{n}(z)$ as 
\begin{align}
\label{eq:I_bessel_series}
    \hspace{-0.5cm}I_\Gamma (\tau, \xi ) &=  \Gamma_1  \int_0^{\tau}  {\rm d} t {\rm e}^{ 2 i f_q \left(  p \cos(\tau - t) + x \sin(\tau - t)  \right)  - \Gamma_1 t } = \Gamma_1 \sum_{n=-\infty}^{\infty} i^n \, J_n(2 f_q |\xi| ) \, e^{i n\phi} \,
\frac{e^{in\tau} - e^{-\Gamma_1\tau}}{i n + \Gamma_1}  \nonumber
\end{align}
where  $ |\xi| =\sqrt{x^2 + p^2}$ and $\phi =\operatorname{atan2}(x,p) $. It follows that, for the initial state $\rho_q = \ket{+} \bra{+}$ \begin{equation}
    \rho^q(\tau)= \begin{pmatrix}
        \chi^{\text{on}}_+  (\tau, 0)  & \chi^{\text{off}} (\tau, 0) \\
        \chi^{\text{off}} (\tau, 0)^*  & \chi^{\text{on}}_-  (\tau, 0)  + \chi^{(1) } (\tau, 0) 
    \end{pmatrix} = \frac{1}{2} \begin{pmatrix}
        {\rm e}^{-\Gamma_1 \tau} &  {\rm e}^{ - \mathcal{C} (\tau) + i  \phi (\tau) } \\
        {\rm e}^{ - \mathcal{C} (\tau) - i  \phi(\tau)}   & 2 - {\rm e}^{-\Gamma_1 \tau} \\ 
    \end{pmatrix} \;.
\end{equation}
as $I_\Gamma (\tau, 0 )  = {\rm e}^{-\Gamma_1 \tau}$ (as for $n\neq 0$, $J_n(0) = 0 $). Thus, the outcome probabilities are given by Eq.~(\ref{eq:probability_SG}).


Specifically, at time $\tau = \pi$, they reduce to 
\begin{equation}
    \sigma(\pi) = \begin{pmatrix}
         s (1 + 2  N_p) + \pi  \Gamma_x & 0\\ 
        0 & \frac{1}{s} (1 + 2 N_p) + \pi  \Gamma_x
    \end{pmatrix} \hspace{1cm}
    \bar{r}^{\text{on}}_{\pm} (\pi)
    =
\begin{pmatrix}
2 \left(f_u \mp f_q\right) \\
 0
\end{pmatrix} \hspace{1cm}
\bar{r}^{\text{off}}(\pi) =  \begin{pmatrix}
          2 f_u - 4 i f_q \Gamma_x \\
       i f_q \left(2 (1 + 2 N_p)/s + \pi \Gamma_x \right)
    \end{pmatrix} \nonumber
\end{equation}
\begin{equation}
    \mathcal{C} (\pi) = f_q^2 \left(4 (1 +2 N_p)/s + 3 \pi \Gamma_x \right) +  \pi \Gamma_z  +  \Gamma_1 \pi /2 \hspace{1cm} \phi(\pi) = 2 \pi f_q f_u +  \frac{\pi}{2} \omega_q  \nonumber
\end{equation}
and at  $\tau = 2 \pi$
\begin{equation}
    \sigma(2 \pi)  = \begin{pmatrix}
         s (1 + 2  N_p) + 2 \pi  \Gamma_x & 0\\ 
        0 & \frac{1}{s} (1 + 2 N_p) + 2 \pi  \Gamma_x
    \end{pmatrix}  \hspace{1cm}
    \bar{r}^{\text{on}}_{\pm} (2 \pi)
    =
\begin{pmatrix}
0 \\
 0
\end{pmatrix} \hspace{1cm}
\bar{r}^{\text{off}}(2 \pi) =  \begin{pmatrix}
   0  \\
        2 i \pi f_q  \Gamma_x
    \end{pmatrix} \nonumber
\end{equation}
\begin{equation}
    \mathcal{C} (2 \pi) = 6 \pi  f_q^2 \Gamma_x + 2 \pi \Gamma_z  +  \Gamma_1 \pi \hspace{1cm} \phi(2 \pi) = 4 \pi f_q f_u +  \pi \omega_q   \nonumber
\end{equation}

One can note that, at $\tau = \pi$, 
\begin{align}
     I_\Gamma (\pi, \xi ) &= \Gamma_1 \sum_{n=-\infty}^{\infty} i^n \, J_n(z) \, e^{i n\phi} \,
\frac{(-1)^n - e^{-\Gamma_1\pi}}{i n + \Gamma_1}   \\
& =  J_0(z)\big(1-e^{-\Gamma\pi}\big) 
+ \sum_{n=1}^{\infty} J_n(z)\,\frac{2 i^n  \big[(-1)^n - e^{-\Gamma_1\pi}\big]}{1 + n^2/\Gamma_1^2}
\left[\cos( n\phi ) + n \sin( n\phi ) /\Gamma_1 \right]
\end{align}
where  we used the symmetry of the Bassel functions $J_{-n}(z) = (-1)^{n} J_{n}(z)  $. 

Having solved the dynamics, let us now focus on the Wigner negativity of the system at $\tau = \pi$. The classical contribution ($\mathcal{W}_c$) is the weighted sum of two Gaussian of covariance matrices given by $\sigma^{-1}(\tau) = (1+2N)\text{diag}(s,1/s) + \pi \Gamma_x \mathds{1}$ and centered at $\bar{r}^{\text{on}}_{\pm}(\pi) =  (2(f_u \pm f_q), 0)^{\rm T}$. They represent the statistical mixture of the Gaussian processes of the diagonal terms of the density matrix.
The quantum contribution leads to the typical fringes of GBCSs given, at time $\tau = \pi$, by
\begin{equation}
   \mathcal{W}_q(r, \pi) =  2 {\rm e}^{-(\mathcal{C}_W (r)+ \mathcal{C} (\pi) )} \cos (\phi_W(r) + \phi (\pi)),
\end{equation}
where the Wigner decay and phase, functions of space variables $\tilde{x}$ and $\tilde{p}$, are
\begin{align}
    \mathcal{C}_W =   (1+2N_p) \frac{8 f_q^2}{s} - \frac{\left(\tilde{x}-2 f_u\right)^2 - (4 \Gamma_x f_q)^2}{\pi  \Gamma_x+  (1 +2 N_p)s} 
     - \frac{\tilde{p}^2 - (\pi \Gamma_x f_q)^2 }{\pi  \Gamma_x + (1 + 2 N_p)/s }\;, \nonumber
\end{align}
and $\phi_W $ is given in the Sec.~\ref{sec:physical_examples}.






\end{document}
%
% ****** End of file apssamp.tex ******
